\chapter{Eerste Deel: Analytiek van het esthetisch oordeelsvermogen}

\kantbook{Eerste Boek: Analytiek van het Schone}

\moment{Eerste moment van het smaakoordeel\footnote{De definitie van smaak die hier wordt gehanteerd, is dat smaak het vermogen is om het mooie te beoordelen. Maar om te bepalen wat nodig is om een object 'mooi' te noemen, moet men de smaakoordelen analyseren.

Bij het zoeken naar de momenten waarop dit oordeelsvermogen in zijn reflectie let, heb ik mij laten leiden door de logische functies van het oordelen (want zelfs in een smaakoordeel is altijd een relatie tot het verstand aanwezig). Ik ben begonnen met het moment van kwaliteit, omdat het esthetische oordeel over het mooie daar als eerste aandacht aan schenkt.}, betreffende zijn kwaliteit}

\section{Het smaakoordeel is esthetisch.}

Om te beslissen of iets mooi is, verbinden we de voorstelling daarvan niet via het verstand met het object om er kennis van te krijgen. In plaats daarvan verbinden we die voorstelling via de verbeeldingskracht (eventueel samen met het verstand) met het subject en met diens gevoel van plezier of onplezier.

Een smaakoordeel is daarom geen kennisoordeel en dus geen logisch oordeel, maar een esthetisch oordeel. Daarmee wordt bedoeld: een oordeel waarvan de bepalende grond \textbf{alleen subjectief kan zijn}.

Elke relatie tussen voorstellingen kan wel objectief zijn --- zelfs die van zintuiglijke indrukken --- want die kan aangeven wat er feitelijk in een empirische voorstelling aanwezig is. Maar dat geldt niet voor de relatie met het gevoel van plezier of onplezier. Via dat gevoel wordt namelijk niets aan het object zelf aangewezen; het subject ervaart alleen hoe het door de voorstelling wordt geraakt.

Een regelmatige, doelmatige structuur begrijpen met onze kenvermogens (of die voorstelling nu helder of vaag is), is iets heel anders dan ons van diezelfde voorstelling bewust zijn met een gevoel van voldoening.

In het laatste geval wordt de voorstelling uitsluitend betrokken op het subject, en wel op zijn 'gevoel van leven', onder de naam gevoel van plezier of onplezier. Dit gevoel vormt de basis van een heel bijzondere beoordelingsvaardigheid. Die draagt niets bij aan kennis, maar houdt de gegeven voorstelling vast binnen het subject als geheel, in relatie tot het volledige voorstellingsvermogen. De geest wordt zich van zijn eigen toestand bewust via dit gevoel.

Voorstellingen die aan een oordeel ten grondslag liggen, kunnen empirisch zijn (en dus esthetisch van aard). Maar het oordeel dat ermee wordt geveld is logisch wanneer die voorstellingen in het oordeel worden betrokken op het object.

Omgekeerd geldt echter: zelfs als de gegeven voorstellingen rationeel zijn, maar in het oordeel uitsluitend worden betrokken op het subject (op diens gevoel), dan is dat oordeel in die mate altijd esthetisch.

\begin{mdframed}[leftmargin=0pt,rightmargin=0pt,backgroundcolor=blue!6,linecolor=white,innertopmargin=6pt,innerbottommargin=6pt]

Kant stelt dat we bij het beoordelen van schoonheid geen kennis over het object vormen. We verbinden de voorstelling van het object niet met het verstand om te bepalen wat het object is, maar met het subject en diens gevoel van plezier of onplezier. Een oordeel over schoonheid is daarom geen kennisoordeel en ook geen logisch oordeel, maar een esthetisch oordeel, omdat de grond ervan uitsluitend subjectief is.

Relaties tussen voorstellingen kunnen objectief zijn en iets zeggen over wat er feitelijk in een voorstelling aanwezig is, zelfs bij zintuiglijke indrukken. Dat geldt echter niet voor de relatie tussen een voorstelling en het gevoel van plezier of onplezier. Dit gevoel zegt niets over het object zelf, maar alleen hoe het subject door de voorstelling wordt geraakt.

Het begrijpen van een regelmatige of doelmatige structuur met onze kenvermogens is iets anders dan het ervaren van voldoening bij diezelfde voorstelling. In het geval van een esthetisch oordeel wordt de voorstelling uitsluitend betrokken op het subject en op diens gevoel van leven, dat zich uit als plezier of onplezier. Dit gevoel vormt de grondslag van een bijzondere beoordelingswijze die geen kennis oplevert, maar de voorstelling binnen het subject vasthoudt in relatie tot het gehele voorstellingsvermogen.

Voorstellingen kunnen empirisch of rationeel zijn. Een oordeel is logisch wanneer deze voorstellingen op het object worden betrokken. Maar wanneer een oordeel, zelfs bij rationele voorstellingen, uitsluitend op het subject en diens gevoel wordt betrokken, is het oordeel in die zin altijd esthetisch.

\bigskip
Kant opent zijn analyse van het schone niet met een definitie, maar met een probleemstelling: hoe kan er een oordeel zijn dat noch op zintuiglijk genot, noch op kennis of moraal berust, en dat toch aanspraak maakt op algemene geldigheid? Het esthetische oordeel lijkt tussen bekende categorieën in te vallen. Het is geen loutere gewaarwording, zoals het aangename, maar ook geen vaststelling van wat een object is of behoort te zijn. Daarmee wordt vanaf het begin duidelijk dat schoonheid niet begrepen kan worden binnen de gebruikelijke oppositie van subjectief gevoel en objectieve kennis.

Om dit te verduidelijken onderscheidt Kant drie typen oordelen. Een oordeel over het aangename berust op zintuiglijk genot en is direct verbonden met behoefte of verlangen; het is lichamelijk bepaald en strikt particulier. Een oordeel over het goede daarentegen is praktisch of moreel van aard: het verwijst naar een doel, een norm of de wil, en gaat altijd gepaard met belang. Het oordeel over het schone onderscheidt zich van beide doordat het weliswaar op een gevoel van plezier berust, maar niet op begeerte of doelgerichtheid. Het betreft de vorm of voorstelling van het object en het welgevallen dat deze oproept, zonder dat dit iets zegt over de objectieve eigenschappen van het ding. Tegelijk is het geen blind gevoel: het vereist reflectie, omdat we ons bewust zijn van de manier waarop het object ons behaagt.

Met dit onderscheid positioneert Kant het esthetische oordeel systematisch tussen perceptie en cognitie. Hij grijpt terug op bekende vermogens — zintuiglijkheid, verstand, wil — maar laat zien dat het smaakoordeel niet volledig door een van deze kan worden verklaard. Tegelijk bereidt hij de weg voor latere paragrafen, waarin deze bijzondere status van het esthetische oordeel cruciaal zal blijken voor begrippen als belangeloosheid, algemene geldigheid en uiteindelijk ook voor de verbinding met moraliteit en teleologie. §1 fungeert zo als een eerste afbakening van het terrein, waarop de rest van de Analytiek van het Schone voortbouwt.

Dat schoonheid geen kennis is, blijkt daarbij essentieel. Juist omdat het oordeel niet vastligt in concepten of regels, is het vrij en belangeloos: het vraagt niet wat een object doet, dient of betekent, maar hoe het zich toont in de ervaring. Tegelijk verklaart dit waarom een subjectief oordeel toch een beroep kan doen op instemming van anderen. In het smaakoordeel ervaren we niet alleen plezier, maar reflecteren we op die ervaring zelf, alsof zij in beginsel deelbaar is. Zo wordt schoonheid begrijpelijk als een bijzondere vorm van reflectieve ervaring, die noch opgaat in gevoel, noch in begrip, en precies daarom een eigen plaats inneemt binnen de menselijke vermogens.

\end{mdframed}

\section{De voldoening die het smaakoordeel bepaalt, is belangeloos}

De voldoening die wij verbinden met de voorstelling van het bestaan van een object, noemen we belang. Zo’n voldoening heeft daarom altijd tegelijk een relatie met het begeervermogen: hetzij als bepalende grond daarvan, hetzij als iets wat noodzakelijk met die bepalende grond samenhangt.

Maar wanneer de vraag is of iets mooi is, wil men niet weten of er voor ons of voor iemand anders iets op het spel staat -- of kan staan -- bij het bestaan van dat ding. Men wil alleen weten hoe wij het beoordelen in loutere beschouwing (aanschouwing of reflectie).

Als iemand mij vraagt of ik het paleis dat ik voor mij zie mooi vind, kan ik gerust zeggen dat ik niet houd van dat soort dingen, die alleen maar gemaakt zijn om naar te gapen; of ik kan, zoals de Irokezen-sachem, zeggen dat mij in Parijs niets beter beviel dan de eethuizen. In echte \textbf{Rousseauiaanse} stijl zou ik zelfs de ijdelheid van de groten kunnen hekelen, die het zweet van het volk verspillen aan zulke overbodige zaken. Ik zou mij er tenslotte ook gemakkelijk van kunnen overtuigen dat, als ik mij op een onbewoond eiland bevond, zonder enige hoop ooit nog mensen tegen te komen, en ik zo’n prachtig bouwwerk alleen door een wens zou kunnen oproepen, ik niet eens de moeite zou nemen dat te doen, zolang ik al een hut had die comfortabel genoeg voor mij was.

Dit alles kan mij worden toegegeven en goedgekeurd; maar daar gaat het hier niet om. Men wil alleen weten of de loutere voorstelling van het object bij mij met voldoening gepaard gaat, hoe onverschillig ik ook mag staan tegenover het bestaan van het object van die voorstelling.

Het is gemakkelijk te zien dat, wanneer ik zeg dat iets \textbf{mooi} is en daarmee wil aantonen dat ik smaak heb, het erop aankomt wat ik met deze voorstelling in mijzelf doe, en niet hoe ik afhankelijk ben van het bestaan van het object. Iedereen moet toegeven dat een oordeel over schoonheid waarin ook maar het minste belang is vermengd, zeer partijdig is en geen zuiver smaakoordeel. Men mag in het geheel niet bevooroordeeld zijn ten gunste van het bestaan van het ding, maar moet in dit opzicht volledig onverschillig zijn om als rechter in zaken van smaak te kunnen optreden.

Wij kunnen deze stelling --- die van het grootste belang is --- op geen betere manier verhelderen dan door de zuivere, belangeloze\footnote{Een oordeel over een object van voldoening kan volledig belangeloos zijn en toch zeer interessant, dat wil zeggen: het is niet op enig belang gebaseerd, maar wekt wel een belang op. Zo zijn alle zuiver morele oordelen. Maar het zuivere smaakoordeel schept op zichzelf zelfs geen enkel belang. Pas in het maatschappelijk leven wordt het interessant om smaak te hebben; de reden daarvoor zal later worden uiteengezet.} voldoening in het smaakoordeel te contrasteren met die voldoening die met belang is verbonden, vooral wanneer wij er zeker van kunnen zijn dat er niet meer soorten belang zijn dan die welke nu zullen worden genoemd.

\begin{mdframed}[leftmargin=0pt,rightmargin=0pt,backgroundcolor=blue!6,linecolor=white,innertopmargin=6pt,innerbottommargin=6pt]

Bij schoonheid speelt bovendien een cruciaal punt: belangeloosheid. We spreken van ‘belang’ wanneer onze voldoening samenhangt met het bestaan van een object, bijvoorbeeld omdat het nuttig is, plezier oplevert, begeerte oproept of een wens vervult. Zo’n gevoel is altijd verbonden met wat we willen of verlangen.

Maar wanneer we iets mooi noemen, doet dat er allemaal niet toe. Het gaat er niet om of het object bestaat, nuttig is, indruk maakt, moreel verdedigbaar is of maatschappelijk waardevol. Zelfs sterke persoonlijke, morele of praktische bezwaren tegen het object mogen geen rol spelen. Het enige wat telt, is of de zuivere voorstelling van het object ons een gevoel van voldoening geeft.

Wie zegt dat iets mooi is, laat daarmee niet zien dat hij het object nodig heeft of waardeert om wat het oplevert, maar dat hij op een bepaalde manier ervaart bij het beschouwen ervan. Elk oordeel over schoonheid waarin ook maar een beetje eigenbelang meespeelt, is daarom geen zuiver smaakoordeel. Om echt over smaak te oordelen, moet men volledig onverschillig zijn ten aanzien van het bestaan of nut van het object.

Deze belangeloosheid is zo fundamenteel, dat zij het best zichtbaar wordt wanneer we haar vergelijken met vormen van voldoening die wél met belang verbonden zijn, en dat is wat Kant vervolgens gaat doen.

\bigskip

De centrale kwestie in §2 is waarom een oordeel over schoonheid geen belang mag dienen. Kant stelt dat schoonheid alleen esthetisch kan worden beoordeeld wanneer het welgevallen vrij is van verlangen, nut of doelgerichtheid. Het plezier dat met schoonheid gepaard gaat, betreft uitsluitend de voorstelling zelf en niet wat het object ons kan opleveren, moreel, praktisch of persoonlijk. Een smaakoordeel verliest zijn eigen aard zodra het verbonden raakt met bezit, gebruik, voordeel of een doel van de wil.

Kant maakt hier expliciet duidelijk wat hij onder belangeloosheid verstaat. Het gaat niet om een ontkenning van plezier, noch om een vorm van ascese, maar om de structuur van het oordeel. In een zuiver esthetisch oordeel wordt niets nagestreefd en niets verlangd; het welgevallen is niet afhankelijk van omstandigheden buiten de voorstelling. Zodra een belang meespeelt, wordt het oordeel heteronoom: het plezier verwijst dan niet meer naar de vorm of verschijning, maar naar het effect of het nut van het object.

Deze afbakening heeft een voorbereidende functie binnen de Kritiek. Door het esthetische oordeel strikt te scheiden van het aangename en van het goede, kan Kant later verklaren hoe een subjectief oordeel toch aanspraak maakt op algemene geldigheid. Alleen een belangeloos oordeel kan worden gedacht als vrij van particuliere neigingen en daarmee geschikt om als voorbeeld voor anderen te dienen. De belangeloosheid vormt zo een noodzakelijke voorwaarde voor de universaliteitsclaim die Kant in de volgende paragrafen zal uitwerken.

Wat hier op het spel staat, is de bijzondere positie van schoonheid tussen verlangen en kennis. Wanneer schoonheid een belang zou dienen, zou zij ofwel samenvallen met genot, ofwel met nut of moraliteit. Door haar belangeloos te denken, behoudt Kant ruimte voor een vorm van ervaring die vrij is, reflectief en niet reducerbaar tot behoeftebevrediging of doelgericht handelen. Juist deze vrijheid maakt het esthetische oordeel tot een eigenstandige en onmisbare vorm van menselijk oordelen.

\end{mdframed}

\section{De voldoening \textit{in het aangename} is met belang verbonden}

Het aangename is datgene wat de zintuigen in de gewaarwording behaagt.
Hier doet zich meteen de gelegenheid voor om een zeer gangbare verwarring aan de kaak te stellen en erop te wijzen, namelijk de verwarring die voortkomt uit de dubbele betekenis van het woord ‘gewaarwording’. Men zegt of denkt namelijk: alle voldoening is zelf een gewaarwording (namelijk van genoegen). Daarom zou alles wat behaagt, juist omdat het behaagt, aangenaam zijn (en, naar gelang van de verschillende graden of relaties tot andere aangename gewaarwordingen, lieflijk, bekoorlijk, verrukkelijk, genietbaar, enz.).

Als men dit echter toestaat, dan zijn zintuiglijke indrukken die de neiging bepalen, beginselen van de rede die de wil bepalen, of louter gereflecteerde vormen van aanschouwing die het oordeelsvermogen bepalen, wat betreft hun werking op het gevoel van genoegen allemaal volkomen hetzelfde. Want dan zou dit steeds slechts de aangenaamheid zijn in de gewaarwording van de eigen toestand. En aangezien uiteindelijk alle inspanning van onze vermogens gericht is op het praktische en daarin als hun doel moet samenkomen, zou men van hen geen andere waardering van dingen en hun waarde kunnen verwachten dan die welke bestaat in de bevrediging die zij beloven. Hoe zij deze bevrediging tot stand brengen, doet er dan uiteindelijk helemaal niet toe. En aangezien hier alleen de keuze van middelen een verschil kan maken, zou men elkaar hoogstens dwaasheid of onbegrip kunnen verwijten, maar nooit laagheid of slechtheid: want iedereen zou, elk op zijn eigen manier kijkend, één en hetzelfde doel nastreven, namelijk bevrediging.

Wanneer een bepaling van het gevoel van genoegen of ongenoegen gewaarwording wordt genoemd, dan betekent deze uitdrukking iets geheel anders dan wanneer ik de voorstelling van een ding (door middel van de zintuigen, als een ontvankelijkheid die tot het kenvermogen behoort) gewaarwording noem. In het laatste geval wordt de voorstelling namelijk op het object betrokken, maar in het eerste geval uitsluitend op het subject, en dient zij tot helemaal geen kennis, zelfs niet tot die waardoor het subject zichzelf kent.

In de bovenstaande uitleg verstaan wij echter onder het woord ‘gewaarwording’ een objectieve zintuiglijke voorstelling. En om niet steeds het risico te lopen verkeerd begrepen te worden, zullen wij datgene wat altijd louter subjectief moet blijven en absoluut geen voorstelling van een object kan vormen, aanduiden met de overigens gebruikelijke naam ‘gevoel’. De groene kleur van de weiden behoort tot de objectieve gewaarwording, als waarneming van een zintuiglijk object; maar de aangenaamheid ervan behoort tot de subjectieve gewaarwording, waardoor geen object wordt voorgesteld, dat wil zeggen: tot het gevoel, waardoor het object wordt beschouwd als een object van voldoening (wat geen kennis ervan is).

Dat mijn oordeel over een object, waardoor ik het aangenaam noem, een belang in dat object uitdrukt, blijkt al daaruit dat het door middel van gewaarwording een verlangen naar objecten van dezelfde soort opwekt. De voldoening veronderstelt dus niet het loutere oordeel erover, maar de relatie van zijn bestaan tot mijn toestand, voor zover die door een dergelijk object wordt beïnvloed. Daarom zegt men van het aangename niet alleen dat het behaagt, maar dat het bevredigt. Het is niet enkel instemming die ik eraan geef; er wordt veeleer een neiging opgewekt. En een oordeel over de gesteldheid van het object behoort zo weinig tot datgene wat op de levendigste wijze aangenaam is, dat zij die uitsluitend op genot uit zijn (want dat is het woord dat de intensiteit van de bevrediging aanduidt) zich er graag boven verheven achten om te oordelen.

\begin{mdframed}[leftmargin=0pt,rightmargin=0pt,backgroundcolor=blue!6,linecolor=white,innertopmargin=6pt,innerbottommargin=6pt]

Het aangename is alles wat onze zintuigen plezier doet ervaren. Kant waarschuwt dat het woord gewaarwording vaak verwarring veroorzaakt: men bedoelt soms de voorstelling van een object en soms het gevoel van genoegen. Het is belangrijk dit onderscheid te maken, want alles wat ons plezier doet, lijkt aangenaam, maar dat betekent nog niet dat het object zelf een bepaalde waarde of kennis heeft.

Zintuiglijke indrukken die verlangen of neiging oproepen, beginselen van de rede die de wil sturen, of gereflecteerde vormen van aanschouwing die ons oordeelsvermogen vormen, hebben qua effect op het gevoel van genoegen allemaal hetzelfde: ze geven bevrediging. Uiteindelijk is alle activiteit van onze vermogens gericht op praktische voldoening, en hoe die wordt bereikt doet er niet toe. Daarom kan men elkaar hoogstens onbegrip of dwaasheid verwijten, maar nooit slechtheid; iedereen streeft immers hetzelfde doel na: bevrediging.

Wanneer Kant over gewaarwording praat in deze context, bedoelt hij het subjectieve gevoel, niet de objectieve waarneming van het ding zelf. De kleur van een weiland bijvoorbeeld, is een objectieve gewaarwording; de aangenaamheid die we ervan ervaren is een subjectief gevoel, een belevenis van voldoening die niets zegt over het object zelf.

Een oordeel dat iets aangenaam is, drukt daarom altijd belang uit: het wekt een verlangen naar soortgelijke objecten op. Het plezier gaat niet alleen over het object zelf, maar over hoe het onze toestand beïnvloedt. Daarom zeggen we dat het aangename bevredigt, niet slechts dat het goedgekeurd wordt. Het roept een neiging of verlangen op. Wie vooral op genot uit is, stelt zich vaak boven het beoordelen van het object; voor hen gaat het alleen om de intensiteit van de bevrediging, niet om kennis of oordeel over het object.

\bigskip
Hier raakt Kant een punt dat vandaag misschien nog zichtbaarder is dan in de achttiende eeuw.

De kern is dat het aangename onvermijdelijk particulier is. 
Wat mij smaakt, prikkelt mijn zintuigen, wekt verlangen op en is verbonden met mijn lichaam, mijn gewoonten, mijn stemming en mijn geschiedenis. 
Het oordeel “dit is aangenaam” zegt daarom in wezen: dit doet mij goed. Daarmee kan het nooit meer zijn dan een verslag van mijn toestand. 
Kant bekritiseert dit genot niet, integendeel, hij erkent het volledig, maar hij begrenst zijn reikwijdte. 
Vanuit het aangename kan geen aanspraak op instemming worden gemaakt. Hier geldt werkelijk: over smaken valt niet te twisten.

Een smaakoordeel daarentegen wil méér dan zeggen dat iets mij behaagt. 
Wanneer ik zeg “dit is mooi”, spreek ik alsof mijn oordeel niet louter persoonlijk is. 
Ik verwacht geen bewijs, maar wél instemming. Dat kan alleen als het oordeel niet berust op zintuiglijke prikkeling of verlangen. 
Zodra het aangename de grond is, valt die verwachting weg: verlangen maakt het oordeel onvrij, belanghebbend en onvermijdelijk subjectgebonden.

In hedendaagse esthetiek zien we deze verwarring voortdurend terug. Denk aan uitspraken als:
“Dit nummer is goed, want het geeft me energie.”
“Deze film is mooi omdat hij me raakte.”
“Dit interieur is prachtig, het voelt zo comfortabel.”

In al deze gevallen wordt intensiteit van ervaring verward met esthetische geldigheid. 
Het object wordt gewaardeerd om wat het doet -- ontroeren, prikkelen, troosten, stimuleren -- niet om hoe het vormelijk verschijnt in een belangeloze beschouwing. 
Vooral in de cultuur van streaming, lifestyle en designmarketing wordt het schone gereduceerd tot wat “lekker werkt”, “fijn voelt” of “je gelukkig maakt”.

Voor Kant is dat precies het probleem: zodra genot het criterium wordt, verschuift esthetiek naar consumptie. 
Het oordeel verliest zijn publieke, aansprekende karakter en wordt een privé-ervaring. 
Het aangename mag intens zijn, zelfs waardevol, maar het kan geen grond zijn voor smaak, omdat het niets vraagt van anderen -- en ook niets van ons, behalve genieten.

\end{mdframed}

\section{De voldoening in het goede is met belang verbonden}

Iets is \textbf{goed} wanneer het alleen door de rede behaagt, door het loutere concept.
We noemen iets \textbf{goed voor iets} (het nuttige) wanneer het alleen als middel plezier geeft; iets anders noemen we \textbf{goed op zichzelf}, wanneer het plezier geeft omwille van zichzelf. 
Beide vormen van goed hebben altijd te maken met het idee van een doel en dus met de relatie van de rede tot (minstens mogelijke) wil. 
Daardoor is er altijd ook een voldoening in het \textbf{bestaan} van een object of een handeling, met andere woorden: een vorm van belang.

Om iets goed te vinden, moet ik altijd weten wat voor ding het object is; ik moet er een concept van hebben. 
Dat is bij schoonheid niet nodig. 
Bloemen, vrije ontwerpen, lijnen die zonder doel in elkaar gevlochten zijn onder de naam bladeren, betekenen niets en hangen van geen bepaald concept af, en geven toch plezier. 
De voldoening in het mooie daarentegen is afhankelijk van reflectie over een object die tot een soort concept leidt (welk precies, is onbepaald) en verschilt daardoor duidelijk van het aangename, dat volledig op gevoel berust.

Vaak lijkt het aangename hetzelfde als het goede. 
Zo wordt wel gezegd dat alle bevrediging (vooral als ze duurzaam is) op zichzelf goed is; ongeveer hetzelfde als zeggen dat duurzaam aangenaam zijn gelijkstaat aan goed zijn. 
Maar dit is een taalkundige vergissing, want de concepten die bij deze woorden horen, zijn niet uitwisselbaar. 
Het aangename, dat het object slechts in relatie tot de zintuigen voorstelt, moet eerst onder de beginselen van de rede worden gebracht via het concept van een doel, voordat het goed kan worden genoemd als object van de wil.
Er is een duidelijk verschil in de relatie tot voldoening wanneer iets dat bevredigt tegelijkertijd \textbf{goed} wordt genoemd. 
Bij het goede speelt altijd de vraag of het middellijk goed (nuttig) of onmiddellijk goed (goed op zichzelf) is, terwijl deze vraag bij het aangename nooit kan spelen, want aangenaam betekent altijd iets dat onmiddellijk behaagt. 
Hetzelfde geldt precies voor het mooie.

Zelfs in alledaags taalgebruik wordt het aangename onderscheiden van het goede. 
Een gerecht dat door specerijen en smaakstoffen de smaak prikkelt, kan aangenaam zijn, maar toch niet goed; het \textbf{pleziert} de zintuigen direct, maar wanneer men het door de rede bekijkt, bijvoorbeeld in het licht van de gevolgen, kan het onaangenaam zijn. 
Ook bij gezondheid zien we dit verschil: het is onmiddellijk aangenaam voor wie gezond is (als afwezigheid van pijn), maar om te zeggen dat het goed is, moet het via de rede aan een doel worden gekoppeld, bijvoorbeeld dat het ons geschikt maakt voor al onze taken.
Wat geluk betreft, denkt men vaak dat de grootste hoeveelheid (qua aantal en duur) aangenaamheid van het leven een waar goed, zelfs het hoogste goed, is. 
Maar de rede protesteert daartegen. 
Aangenaamheid is puur genot. 
Als dit alles was, zou het niet uitmaken op welke manier men het genot verkrijgt: passief door de gaven van de natuur of actief door eigen inspanning. 
Een mens die alleen \textbf{voor genot} leeft, zou daardoor nooit absolute waarde hebben. 
Pas door wat iemand doet zonder oog voor genot, in volledige vrijheid en onafhankelijk van wat de natuur hem passief kan geven, krijgt zijn bestaan als persoon een absolute waarde. 
En geluk, met al zijn aangenaamheid, is daarmee ver verwijderd van een onvoorwaardelijk goed.\footnote{Een verplichting tot genot is een duidelijke absurditeit. Hetzelfde geldt voor een veronderstelde verplichting in alle handelingen die uitsluitend gericht zijn op genot, hoe geestelijk verfijnd of verfraaid dit ook mag zijn, zelfs al zou het een mystiek, zogenaamd hemels genot betreffen.}

Ondanks dit verschil tussen aangenaam en goed, hebben ze één overeenkomst: beide zijn altijd met belang verbonden. 
Niet alleen het aangename (§3) en het middellijk goede (het nuttige, dat plezier geeft als middel tot iets aangenaams), maar ook wat absoluut en in alle opzichten goed is --- het morele goede --- draagt het hoogste belang met zich mee. 
Want het goede is het object van de wil (dat wil zeggen: van een begeervermogen dat door de rede wordt bepaald). 
En iets willen en voldoening hebben in het bestaan ervan --- met andere woorden, er belang bij hebben --- is precies hetzelfde.

\begin{mdframed}[leftmargin=0pt,rightmargin=0pt,backgroundcolor=blue!6,linecolor=white,innertopmargin=6pt,innerbottommargin=6pt]

In §4 bespreekt Kant het verschil tussen het aangename en het goede en hoe ze samenhangen met belang. Iets is goed wanneer het alleen door de rede behaagt, door het idee ervan, niet door zintuiglijk genot. Soms gaat het om iets dat goed voor iets is, dat wil zeggen: het behaagt alleen als middel. Soms gaat het om iets dat goed op zichzelf is, dat wil zeggen: het behaagt omwille van zichzelf. Beide vormen van goed hebben altijd te maken met een doel en met de wil, en daardoor is er altijd ook een soort voldoening of belang verbonden aan het bestaan van het object of de handeling.

Om iets goed te vinden, moet je weten wat voor ding het is; je hebt er een concept van nodig. Bij schoonheid is dat niet het geval. Bloemen, vrije ontwerpen of lijnen die zonder doel in elkaar gevlochten zijn, kunnen ons plezier doen, ook al betekenen ze niets en hangen ze van geen concept af. De voldoening in het mooie is echter afhankelijk van reflectie op het object, waardoor er een soort concept ontstaat, en onderscheidt zich zo van het aangename, dat volledig op gevoel berust.

Vaak lijkt het aangename hetzelfde als het goede, vooral wanneer bevrediging langdurig is. Kant waarschuwt dat dit een vergissing is: het aangename werkt direct op de zintuigen en pleziert onmiddellijk, terwijl het goede via de rede wordt beoordeeld en steeds de vraag oproept of iets nuttig is of goed op zichzelf. Zo kan een smakelijk gerecht aangenaam zijn, maar toch niet goed, omdat het, wanneer men het door de rede bekijkt, misschien onaangename gevolgen heeft. Ook gezondheid voelt aangenaam (afwezigheid van pijn), maar is pas goed wanneer het ons geschikt maakt voor onze taken en doelen. Geluk in de vorm van aangenaamheid alleen is nooit een absoluut goed. Een mens die alleen voor genot leeft, heeft geen intrinsieke waarde, ook al helpt hij anderen hetzelfde genot te ervaren. Absolute waarde ontstaat pas door te handelen zonder oog voor genot, volledig vrij en onafhankelijk van wat de natuur passief kan geven.

Toch hebben het aangename en het goede één overeenkomst: beide zijn altijd met belang verbonden. Het aangename geeft directe voldoening, het goede geeft voldoening omdat het object of de handeling ons iets oplevert of een doel dient. Het goede is het object van de wil: iets willen en tevreden zijn met het bestaan ervan betekent dat we er belang bij hebben.

\bigskip
In §4 wordt duidelijk waarom Kant het schone zorgvuldig tussen het aangename en het moreel goede plaatst, en waarom het juist daardoor zijn eigen statuut behoudt. Het schone is niet aangenaam, omdat het niet op zintuiglijk genot of verlangen berust. Maar het is ook niet moreel goed, omdat het niet voortkomt uit rede, plicht of een begrip van wat behoort te zijn. Schoonheid bevalt zonder prikkel en zonder wet.

Het moreel goede kan, zoals Kant benadrukt, wel degelijk interesseloos zijn: een morele handeling wordt niet gemotiveerd door voordeel of genot. Maar zij is nooit begriploos. Moraal veronderstelt altijd een concept, van plicht, van wet, van doelmatigheid van de wil. Wie iets goed noemt, weet waarom het goed is, of moet dat tenminste kunnen verantwoorden. In het esthetisch oordeel ontbreekt precies dit: we kunnen niet uitleggen waarom iets mooi is zonder het oordeel zelf te ondergraven. Schoonheid vraagt geen gehoorzaamheid, maar instemming.

Daarom scheidt Kant schoonheid en moraal, zonder ze te ontkoppelen. Ze raken elkaar in hun interesseloosheid en in hun aanspraak op algemeenheid, maar ze blijven verschillend in structuur. Het schone oefent geen bevel uit; het nodigt uit. Het moreel goede verplicht; het schone overtuigt niet, maar stemt af.

Als we schoonheid volledig moraliseren, verliezen we precies deze vrijheid. Schoonheid wordt dan een instrument: zij moet iets zeggen, iets leren, iets rechtvaardigen. Kunst en esthetische ervaring worden beoordeeld op hun boodschap, hun intentie of hun morele correctheid, niet op hun vormelijke verschijning en de vrije harmonie die zij oproepen. Het gevolg is dat esthetische ervaring wordt gereduceerd tot illustratie van waarden die elders al vastliggen.

Daarmee verliezen we ook iets fundamenteels menselijks: een ruimte waarin we kunnen oordelen zonder te willen bezitten, zonder te willen gebruiken en zonder te moeten gehoorzamen. Bij Kant is schoonheid juist waardevol omdat zij ons laat ervaren dat er een vorm van zin en gemeenschappelijkheid bestaat zonder dwang. Door haar te moraliseren maken we haar veiliger, misschien, maar ook armer, en uiteindelijk onvrij.
\end{mdframed}

\section{Vergelijking van de drie specifiek verschillende soorten voldoening}

Het aangename en het goede hebben beide een relatie tot het vermogen tot verlangen, en brengen daardoor voldoening met zich mee. 
Bij het aangename gaat het om een pathologisch bepaalde voldoening, die ontstaat door prikkels, terwijl het goede een pure praktische voldoening geeft, die niet alleen voortkomt uit de voorstelling van het object, maar tegelijk uit de veronderstelde relatie van het subject tot het bestaan van het object. 
Het is dus niet slechts het object zelf dat behaagt, maar ook zijn bestaan.
Het smaakoordeel daarentegen is uitsluitend \textbf{beschouwend}: het is een oordeel dat, volledig onverschillig ten aanzien van het bestaan van een object, alleen de constitutie van het object verbindt met het gevoel van genoegen of ongenoegen. 
Deze contemplatie is bovendien niet gericht op concepten; het smaakoordeel is immers geen kennisgericht oordeel, noch theoretisch noch praktisch, en is dus noch op concepten \textbf{gebaseerd}, noch op hen \textbf{gericht}.

Het aangename, het mooie en het goede duiden daarom drie verschillende relaties aan van voorstellingen tot het gevoel van genoegen en ongenoegen. 
Aan de hand hiervan onderscheiden we objecten of soorten voorstellingen van elkaar. 
Ook de woorden die bij elk van deze relaties horen en waarmee men het plezier in elk aanduidt, verschillen. 
\textbf{Aangenaam} is datgene wat iedereen noemt wat hem \textbf{bevredigt}; \textbf{mooi} is datgene wat louter \textbf{plezier} geeft; \textbf{goed} is datgene dat men \textbf{waardeert} of \textbf{goedkeurt}, dat wil zeggen: dat waaraan men een objectieve waarde toekent.
Het aangename geldt ook voor niet-rationele dieren. 
Het mooie geldt alleen voor mensen, dat wil zeggen voor wezens die zowel dierlijk als rationeel zijn, niet uitsluitend rationeel (zoals geesten), maar als wezens die tegelijkertijd dierlijk zijn. 
Het goede geldt echter voor elk rationeel wezen in het algemeen, een uitspraak die pas volledig gerechtvaardigd en verklaard kan worden in de volgende hoofdstukken.
Men kan zeggen dat van deze drie soorten voldoening alleen de voldoening bij de smaak voor het mooie belangeloos en \textbf{vrij} is; geen enkel belang, noch dat van de zintuigen noch dat van de rede, dwingt tot goedkeuring. 
Zo kan men zeggen dat voldoening in de drie genoemde gevallen verband houdt met \textbf{neiging}, \textbf{sympathie} of \textbf{respect}. 
Want \textbf{sympathie} is de enige vrije voldoening. 
Een object van neiging of iets dat door een wet van de rede wordt opgelegd ten behoeve van verlangen, laat ons geen vrijheid om zelf iets tot object van plezier te maken. 
Elk belang veronderstelt een behoefte of wekt er een op, en als bepalende grond van goedkeuring laat het oordeel over het object geen vrijheid meer.

Wat de neiging betreft in het geval van het aangename, zegt men wel dat honger de beste kok is en dat mensen met een gezonde eetlust alles waarderen wat eetbaar is; zulke voldoening laat dus geen keus zien volgens smaak. 
Pas wanneer de behoefte is bevredigd, kan men onderscheiden wie smaak heeft en wie niet. 
Op dezelfde manier zijn er gedragingen zonder deugd, beleefdheid zonder welwillendheid, correctheid zonder eerbaarheid, enzovoort. 
Waar de morele wet spreekt, bestaat er objectief gezien geen vrije keuze meer over wat gedaan moet worden; smaak tonen in het eigen gedrag (of bij het beoordelen van anderen) is iets heel anders dan het uitdrukken van morele overtuiging: dit laatste bevat een bevel en wekt een behoefte op, terwijl modieuze smaak juist alleen speelt met de objecten van voldoening zonder zich aan een van hen te hechten.

\paragraph{Definitie van het mooie afgeleid uit het eerste moment}
\textbf{Smaak} is het vermogen om een object of soort voorstelling te beoordelen door middel van voldoening of ongenoegen \textbf{zonder enig belang}. 
Het object van een dergelijke voldoening wordt \textbf{mooi} genoemd.

\begin{mdframed}[leftmargin=0pt,rightmargin=0pt,backgroundcolor=blue!6,linecolor=white,innertopmargin=6pt,innerbottommargin=6pt]

In §5 vergelijkt Kant de drie soorten voldoening die hij heeft besproken: het aangename, het mooie en het goede. Zowel het aangename als het goede hebben een relatie met ons verlangen en geven daardoor voldoening. Bij het aangename gaat het om een zintuiglijk bepaalde voldoening, die ontstaat door prikkels en behoeften. Bij het goede is de voldoening praktisch en ontstaat ze niet alleen uit de voorstelling van het object, maar ook uit de veronderstelde relatie van het subject tot het bestaan van het object: niet alleen het object zelf behaagt, maar ook dat het bestaat.

Het smaakoordeel daarentegen is uitsluitend beschouwend. Het verbindt alleen de voorstelling van het object met het gevoel van plezier of ongenoegen, ongeacht of het object werkelijk bestaat. Dit oordeel is niet cognitief: het is niet gebaseerd op concepten en richt zich er ook niet op.

Kant benadrukt dat het aangename, het mooie en het goede elk een andere relatie tot genoegen en belang hebben. Aangenaam is wat ons direct bevredigt, mooi is wat puur plezier geeft, en goed is wat we waarderen of goedkeuren, met een zekere objectieve waarde. Daarbij geldt het aangename ook voor dieren, het mooie alleen voor mensen, omdat het reflectie en rede vereist, en het goede voor alle rationele wezens.

Het belangrijkste onderscheid is dat alleen de voldoening in de smaak voor het mooie vrij en belangeloos is. Bij het aangename en het goede speelt altijd een vorm van belang mee: het aangename beantwoordt aan behoefte, het goede dient een doel of wil. Vrije voldoening betekent dat geen enkele nood of verlangens ons oordeel beïnvloedt. Wie puur smaak beoefent, speelt met objecten van voldoening zonder dat hij eraan gebonden is.

Kant concludeert dat smaak het vermogen is om een object te beoordelen door middel van plezier of ongenoegen zonder dat er belang bij is. Het object van een dergelijke voldoening wordt mooi genoemd.

\bigskip
In §5 krijgt schoonheid haar plaats niet doordat Kant eindelijk zegt wat zij is, maar doordat hij laat zien wat zij niet is. Dat is geen omweg, maar een principiële keuze. Schoonheid verschijnt hier als een eigen categorie van waardering juist doordat zij zich onttrekt aan de twee domeinen waarin waardering normaal gesproken thuis is: begeerte (het aangename) en wil (het goede).

Het aangename waardeert vanuit behoefte: iets is waardevol omdat het mij goed doet, mij bevredigt, mij aantrekt. Het goede waardeert vanuit normativiteit: iets is waardevol omdat het behoort te zijn, omdat het rationeel of moreel te rechtvaardigen is. In beide gevallen is waardering functioneel: zij staat in dienst van iets anders, genot, handelen, doelmatigheid. Schoonheid daarentegen heeft geen functie buiten zichzelf. Zij waardeert zonder te gebruiken, zonder te willen, zonder te rechtvaardigen.

Dat maakt schoonheid tot een eigen categorie: niet omdat zij een positief kenmerk deelt met aangenaam of goed, maar omdat zij een andere relatie tot waardering instelt. Bij schoonheid gaat het niet om wat het object doet (bevredigen, motiveren), maar om hoe het verschijnt in een vrije beschouwing. Het oordeel is contemplatief, niet praktisch; het waardeert zonder inzet.

Waarom bouwt Kant dit op via uitsluiting in plaats van definitie? Omdat elke positieve definitie schoonheid meteen zou koloniseren door een ander domein. Zodra je schoonheid definieert als bron van genot, glijdt zij af naar het aangename. Zodra je haar definieert als uitdrukking van een idee of waarde, glijdt zij af naar het goede. Een definitie zou haar reduceren tot middel, symptoom of drager van iets anders.

Door schoonheid negatief te positioneren, niet zintuiglijk, niet moreel, niet geïnteresseerd, bewaart Kant haar autonomie. Pas als die autonomie veiliggesteld is, kan later worden onderzocht wat haar positieve grond is (de vrije harmonie van verbeelding en verstand). §5 is dus geen voltooiing, maar een afbakening: een filosofische vrijwaringszone.

Schoonheid wordt hier niet verklaard, maar vrijgemaakt. Dat is essentieel, omdat Kant wil laten zien dat er een vorm van waardering bestaat die noch private voorkeur is, noch normatieve verplichting, en die juist daarom een unieke plaats inneemt in het menselijk leven.
\end{mdframed}

\moment{Tweede moment van het smaakoordeel, betreffende zijn kwantiteit}

\section{Het schone is datgene wat, zonder begrippen, wordt voorgesteld als het object van een algemene voldoening.}

Deze definitie van het schone kan worden afgeleid uit de voorafgaande uitleg ervan als een object van voldoening zonder enig belang. 
Want men kan datgene waarvan men zich bewust is dat de voldoening eraan in zijn eigen geval zonder enig belang is, op geen enkele andere wijze beoordelen dan zo, dat het een grond van voldoening voor iedereen moet bevatten. 
Aangezien zij namelijk niet gegrond is in enige neiging van het subject (noch in enig ander daaraan ten grondslag liggend belang), maar degene die oordeelt zich geheel \textbf{vrij} voelt ten aanzien van de voldoening die hij aan het object toekent, kan hij als grond van die voldoening geen bijzondere voorwaarden ontdekken die slechts zijn eigen subject betreffen, en moet hij haar daarom beschouwen als gegrond in zulke voorwaarden die hij ook bij ieder ander kan veronderstellen; bijgevolg moet hij menen gronden te hebben om van iedereen een soortgelijk welgevallen te verwachten.
Daarom zal hij over het schone spreken alsof schoonheid een eigenschap van het object was en het oordeel logisch (een kennis van het object constituerend door begrippen ervan), hoewel het slechts esthetisch is en alleen een betrekking van de voorstelling van het object tot het subject bevat; omdat het toch deze overeenkomst met het logisch oordeel heeft, dat zijn geldigheid voor iedereen kan worden voorondersteld. Maar deze algemeenheid kan niet uit begrippen voortkomen. Want er is geen overgang van begrippen naar het gevoel van welgevallen of onwelgevallen (behalve bij zuiver praktische wetten, die echter een belang met zich meebrengen van een soort die niet met het zuivere smaakoordeel verbonden is). Bijgevolg moet aan het smaakoordeel, met het bewustzijn van een daarin vervatte abstractie van alle belang, een aanspraak op geldigheid voor iedereen worden verbonden zonder de algemeenheid die aan objecten toekomt, dat wil zeggen: het moet verbonden zijn met een aanspraak op subjectieve algemeenheid.

\begin{mdframed}[leftmargin=0pt,rightmargin=0pt,backgroundcolor=blue!6,linecolor=white,innertopmargin=6pt,innerbottommargin=6pt]
Iets is schoon wanneer het ons behaagt zonder dat we daarvoor een begrip nodig hebben, en wanneer dit welbehagen gepaard gaat met de aanspraak dat iedereen het zou moeten kunnen delen. Als ik iets mooi vind zonder er enig belang bij te hebben, dan ervaar ik mijn oordeel niet als iets puur persoonlijks of toevalligs. Ik voel me vrij ten opzichte van mijn voorkeuren, verlangens of doelen, en daarom kan ik de grond van mijn genoegen niet herleiden tot iets wat alleen voor mij geldt.

Juist omdat mijn voldoening niet op een particuliere neiging berust, moet ik aannemen dat zij ook voor anderen toegankelijk is. Daarom verwacht ik, vanzelfsprekend en zonder bewijs, dat anderen hetzelfde genoegen zouden kunnen ervaren. Zo ontstaat de neiging om over schoonheid te spreken alsof zij een eigenschap van het object zelf is, en alsof mijn oordeel een soort kennis is.

Toch is dat oordeel geen kennis. Het zegt niets over het object via begrippen, maar alleen iets over de verhouding tussen mijn voorstelling van het object en mijn gevoel. Het lijkt op een logisch oordeel doordat het aanspraak maakt op algemene geldigheid, maar die algemeenheid kan niet uit begrippen voortkomen, omdat begrippen op zichzelf geen gevoelens van genoegen of ongenoegen kunnen opleveren.

Daarom bevat het smaakoordeel een bijzondere aanspraak: het vraagt om algemene instemming zonder objectieve regels. Het is geen universele geldigheid zoals bij kennis, maar een subjectieve algemeenheid: een aanspraak op gedeeld welbehagen, gebaseerd op het feit dat we onze menselijke vermogens als gemeenschappelijk veronderstellen.

\bigskip
In §6 wordt scherp zichtbaar wat voor Kant misschien wel het meest paradoxale kenmerk van het smaakoordeel is: het is onmiskenbaar subjectief, maar het spreekt met een universele stem. De vraag is dus niet óf mensen het eens zijn, maar waarom iemand die oordeelt zich gerechtigd voelt instemming te verwachten.

Die aanspraak op algemene geldigheid is nadrukkelijk niet statistisch. Kant bedoelt niet: “de meeste mensen vinden dit mooi”, laat staan “iedereen vindt dit mooi”. Het gaat om een normatieve claim: zo zou men moeten oordelen, als men vrij en onbevangen waarneemt. De algemeenheid zit dus niet in feitelijke overeenstemming, maar in de vorm van het oordeel zelf.

De sleutel ligt in het belangeloze karakter van het esthetisch welbehagen. Juist omdat ik geen particulier verlangen, nut of doel volg, ervaar ik mijn oordeel niet als privé of toevallig. Ik kan het niet herleiden tot mijn smaak, opvoeding, stemming of voordeel. Daardoor lijkt het oordeel “losgemaakt” van mij als individu. Wat overblijft, is een manier van waarnemen en beoordelen die ik als menselijk in het algemeen ervaar.

Hieruit ontstaat het gevoel van gerechtigdheid: ik verwacht instemming niet omdat ik anderen wil overtuigen, maar omdat ik geen reden zie waarom zij, als mensen met dezelfde vermogens, anders zouden moeten oordelen. De aanspraak op universaliteit berust dus op een impliciete veronderstelling van een gedeelde menselijke structuur: verbeelding en verstand werken bij iedereen volgens dezelfde basisprincipes.

Tegelijk benadrukt Kant dat dit geen kennis oplevert. Er is geen regel, geen bewijs, geen concept dat instemming kan afdwingen. De aanspraak blijft een “alsof”: we spreken alsof schoonheid een eigenschap van het object is, terwijl zij in feite ontstaat in de relatie tussen object en subject. Juist dit “alsof” maakt het smaakoordeel zo bijzonder: het balanceert tussen persoonlijke ervaring en intersubjectieve geldigheid.

Kort gezegd: we voelen ons gerechtigd instemming te verwachten omdat het esthetisch oordeel, in zijn belangeloosheid, boven het particuliere uitstijgt zonder ooit objectief te worden. Het appelleert niet aan wat mensen feitelijk vinden, maar aan wat zij als vrije, gedeelde menselijke beoordelaars zouden kunnen ervaren.
\end{mdframed}

\section{Vergelijking van het schone met het aangename en het goede aan de hand van het hierboven genoemde kenmerk}

Wat het \textbf{aangename} betreft, is iedereen ermee tevreden dat zijn oordeel, dat hij op een privégevoel grondt en waarin hij van een object zegt dat het hem behaagt, slechts tot zijn eigen persoon wordt beperkt. Daarom is hij er ook volkomen content mee wanneer, als hij zegt dat mousserende wijn van de Canarische Eilanden aangenaam is, iemand anders zijn uitdrukking verbetert en hem eraan herinnert dat hij moet zeggen: “Het is voor mij aangenaam”. En dit geldt niet alleen voor de smaak van tong, gehemelte en keel, maar ook voor wat voor iemand aangenaam kan zijn voor ogen en oren. Voor de één is de kleur violet zacht en lieflijk, voor de ander doods en levenloos. De één houdt van de klank van blaasinstrumenten, de ander van die van strijkinstrumenten. Het zou dwaas zijn om in zulke gevallen het oordeel van een ander dat van het onze verschilt, te betwisten met de bedoeling het als onjuist te veroordelen, alsof het logisch met het onze in tegenspraak was. Met betrekking tot het aangename geldt dus het beginsel: \textbf{ieder heeft zijn eigen} smaak (van de zintuigen).

Met het schone is het echter geheel anders. Het zou belachelijk zijn als iemand die prat gaat op zijn smaak zichzelf aldus zou willen rechtvaardigen: “Dit object (het gebouw dat wij bekijken, de kleding die iemand draagt, het gedicht dat ter beoordeling wordt voorgelegd) is \textbf{voor mij} mooi.” Want hij mag het niet \textbf{mooi} noemen als het slechts hem behaagt. Veel dingen kunnen voor hem charme en aangenaamheid hebben, daar zal niemand aanstoot aan nemen; maar wanneer hij iets mooi noemt, dan verwacht hij van anderen dezelfde voldoening: hij oordeelt niet slechts voor zichzelf, maar voor iedereen, en spreekt over schoonheid alsof zij een eigenschap van de dingen was. Daarom zegt hij dat het \textbf{ding} mooi is, en rekent hij niet op de instemming van anderen met zijn oordeel van voldoening omdat hij vaak heeft ervaren dat zij met het zijne overeenstemmen, maar \textbf{eist} hij deze instemming van hen. Hij berispt hen wanneer zij anders oordelen en ontzegt hun smaak, hoewel hij toch verlangt dat zij die zouden moeten hebben. In dit opzicht kan men dus niet zeggen: “ieder heeft zijn eigen bijzondere smaak”. Dat zou evenveel betekenen als te zeggen dat er helemaal geen smaak bestaat, dat wil zeggen: geen esthetisch oordeel dat terecht aanspraak kan maken op de instemming van iedereen.

Niettemin treft men ook bij het aangename soms eenstemmigheid in het oordelen aan, met het oog waarop men sommigen smaak ontzegt en anderen toekent, niet in de betekenis van een organisch zintuig, maar als een vermogen om in het algemeen over het aangename te oordelen. Zo zegt men van iemand die weet hoe hij zijn gasten met aangename dingen kan vermaken (door genot via alle zintuigen), zodat zij allen tevreden zijn, dat hij smaak heeft. Maar hier wordt de algemeenheid slechts vergelijkenderwijs verstaan, en in dit geval zijn er slechts \textbf{algemene} regels (zoals alle empirische regels zijn), geen \textbf{universele}, zoals die waarop het smaakoordeel over het schone zich waagt of aanspraak maakt. Het is een oordeel in relatie tot gezelligheid, voor zover het op empirische regels berust.
Wat het goede betreft, maken oordelen daarover weliswaar eveneens terecht aanspraak op geldigheid voor iedereen; maar het goede wordt alleen als object van een universele voldoening voorgesteld \textbf{door middel van een begrip}, wat noch bij het aangename, noch bij het schone het geval is.

\begin{mdframed}[leftmargin=0pt,rightmargin=0pt,backgroundcolor=blue!6,linecolor=white,innertopmargin=6pt,innerbottommargin=6pt]
In §7 vergelijkt Kant het schone met het aangename en het goede om de unieke positie van schoonheid te verduidelijken.

Bij het aangename is het oordeel volledig privé: iemand zegt dat iets hem behaagt en rekent daar verder niet op dat anderen hetzelfde zullen ervaren. Het principe “ieder zijn smaak” geldt hier volledig. Of het nu gaat om wijn, kleuren of muziek, er is geen reden om het oordeel van iemand anders te betwisten; verschillen zijn normaal en geen teken van gebrek aan smaak. Soms kan er toch overeenstemming ontstaan over wat aangenaam is, bijvoorbeeld wanneer iemand erin slaagt gasten op allerlei manieren tevreden te stellen. Maar deze instemming is slechts relatief en gebaseerd op empirische ervaring, niet op een universele regel.

Bij het schone ligt dit heel anders. Wanneer iemand iets mooi noemt, verwacht hij automatisch dat anderen dezelfde voldoening zullen ervaren. Het oordeel is niet beperkt tot het eigen genoegen, maar heeft de pretentie van universaliteit. Men spreekt over schoonheid alsof het een eigenschap van het object zelf is, ook al is het oordeel slechts esthetisch en subjectief. Anders dan bij het aangename, eist het schone instemming: wie anders oordeelt, wordt als “zonder smaak” gezien, hoewel Kant niet letterlijk beweert dat iedereen het eens moet zijn. Het verschil zit in de subjectieve algemeenheid: schoonheid richt zich op gedeelde menselijke vermogens en veronderstelt dat anderen, net als wij, vrij en onbevangen zouden oordelen.

Het goede tenslotte maakt ook aanspraak op universele geldigheid, maar op een heel andere manier: het berust op begrippen en het rationele idee van doel en morele waarde. Het wordt dus via een concept universeel, terwijl het aangename en het schone zonder begrippen ervaren worden. Schoonheid staat hier tussenin: het is universeel qua aanspraak, maar niet via begrippen of praktische doelen, en verschilt daardoor fundamenteel van zowel het aangename als het goede.

\bigskip
In §7 laat Kant zien dat het smaakoordeel op het eerste gezicht lijkt op een logisch oordeel, maar dat deze gelijkenis heel specifiek is en niet inhoudelijk. Net zoals een logisch oordeel aanspraak maakt op geldigheid voor iedereen, zo doet ook het esthetisch oordeel dat -- het spreekt met een pretentie van universaliteit.

Het gaat hierbij niet om waarheid of feitelijke overeenstemming, zoals bij kennis. Het smaakoordeel zegt niet: “dit object is objectief mooi,” noch vereist het bewijs of concepten. De gelijkenis zit puur in de vorm van het oordeel: het doet alsof het object een eigenschap bezit die voor iedereen geldig is, en het verwacht instemming van anderen.

Deze gelijkenis vertelt iets wezenlijks over onze menselijke behoefte aan overeenstemming. Zelfs in het subjectieve domein van smaak voelen we ons gerechtigd om een soort consensus te verwachten. We ervaren ons oordeel niet als louter persoonlijk, maar als iets dat, mits vrij van eigenbelang, door ieder gedeeld zou kunnen worden. De aanspraak op algemene geldigheid weerspiegelt dus een sociale en cognitieve behoefte: we zoeken bevestiging dat ons vrije oordeel iets universeelmenselijks raakt, iets dat anderen met dezelfde vermogens zouden kunnen ervaren.

Kortom: het smaakoordeel lijkt op een logisch oordeel in zijn vorm, omdat het een subjectieve algemeenheid veronderstelt, maar het verschilt radicaal in inhoud: het gaat niet om kennis of waarheid, maar om gedeeld genot en esthetische ervaring.
\end{mdframed}

\section{De universaliteit van de voldoening wordt in een smaakoordeel slechts als subjectief voorgesteld}

Deze bijzondere bepaling van de universaliteit van een esthetisch oordeel, zoals die in een smaakoordeel wordt aangetroffen, is iets opmerkelijks, niet zozeer voor de logicus, maar zeker voor de transcendentaal filosoof. Het ontdekken van de oorsprong ervan vergt van hem niet weinig inspanning, maar onthult tegelijk een eigenschap van ons kenvermogen die zonder deze analyse onbekend zou zijn gebleven.

Ten eerste moet men er volledig van overtuigd zijn dat men door middel van het smaakoordeel (over het schone) de voldoening aan een object aan iedereen toeschrijft, zonder deze echter op een begrip te baseren (want dan zou het om het goede gaan). En deze aanspraak op algemene geldigheid behoort zo wezenlijk tot een oordeel waarin wij iets schoon noemen, dat zonder deze gedachte het niemand ooit zou invallen deze uitdrukking te gebruiken. Alles wat zonder begrip behaagt, zou dan eenvoudig tot het aangename worden gerekend, waarbij iedereen zijn eigen mening kan hebben en niemand instemming van anderen met zijn oordeel verwacht, wat echter bij oordelen over schoonheid steeds wél het geval is.

Het eerste kan ik de smaak van de zintuigen noemen, het tweede de smaak van de reflectie, voor zover de eerste slechts privé-oordelen over een object velt, terwijl de tweede vermeend algemeen geldige (publieke) oordelen uitspreekt. Beide zijn echter esthetische (niet praktische) oordelen over een object, en betreffen enkel de relatie van zijn voorstelling tot het gevoel van genoegen en ongenoegen. Het is nu merkwaardig dat bij de smaak van de zintuigen de ervaring niet alleen leert dat haar oordeel (van genoegen of ongenoegen aan iets) niet algemeen geldig is, maar ook dat iedereen daarin van nature zo bescheiden is dat hij deze instemming niet eens aan anderen toeschrijft (hoewel in zulke oordelen vaak een vrij grote eensgezindheid wordt aangetroffen). Daarentegen kan de smaak van de reflectie, die, zoals de ervaring leert, vaak genoeg in haar aanspraak op algemene geldigheid wordt afgewezen, het toch mogelijk achten (en doet dat ook daadwerkelijk) om oordelen voor te stellen die universeel instemming zouden kunnen eisen, en verwacht die instemming ook daadwerkelijk van iedereen bij elk van haar oordelen. De mensen die zulke oordelen vellen, zien daarbij geen tegenspraak in de mogelijkheid van deze aanspraak, maar slechts onenigheid over de juiste toepassing van dit vermogen in afzonderlijke gevallen.

Hier moet allereerst worden opgemerkt dat een universaliteit die niet berust op begrippen van objecten (zelfs niet op louter empirische) in het geheel niet logisch is, maar esthetisch; dat wil zeggen: zij bevat geen objectieve kwantiteit van het oordeel, maar slechts een subjectieve. Hiervoor gebruik ik ook de uitdrukking gemeenschappelijke geldigheid, die niet de geldigheid van de relatie van een voorstelling tot het kenvermogen voor elk subject aanduidt, maar tot het gevoel van genoegen en ongenoegen. (Dezelfde uitdrukking kan overigens ook voor de logische kwantiteit van het oordeel worden gebruikt, mits men eraan toevoegt: objectief universele geldigheid, om haar te onderscheiden van de louter subjectieve, die altijd esthetisch is.)

Een objectief universeel geldig oordeel is altijd ook subjectief universeel geldig: als een oordeel geldig is voor alles wat onder een bepaald begrip valt, dan is het ook geldig voor iedereen die een object door middel van dit begrip voorstelt. Maar uit een subjectief universele geldigheid, dat wil zeggen: uit een esthetische universaliteit die op geen enkel begrip berust, kan in het geheel geen gevolgtrekking naar logische universaliteit worden gemaakt, omdat dit soort oordeel het object helemaal niet betreft. Juist daarom moet de esthetische universaliteit die aan een oordeel wordt toegekend van een bijzondere aard zijn: want het predicaat van schoonheid wordt niet met het begrip van het object verbonden in zijn volledige logische sfeer, en toch strekt het zich uit over de gehele sfeer van degenen die oordelen.

Wat de logische kwantiteit betreft, zijn alle smaakoordelen enkelvoudige oordelen. Want aangezien ik het object onmiddellijk aan mijn gevoel van genoegen of ongenoegen moet voorleggen, en niet via begrippen, kan het niet de kwantiteit hebben van een objectief algemeen geldig oordeel. Wanneer echter de enkelvoudige voorstelling van het object van een smaakoordeel, overeenkomstig de voorwaarden die dit oordeel bepalen, door vergelijking tot een begrip wordt omgevormd, kan daaruit een logisch universeel oordeel ontstaan. Zo verklaar ik door een smaakoordeel dat de roos die ik aanschouw mooi is. Daarentegen is het oordeel dat uit de vergelijking van vele enkelvoudige oordelen voortkomt, dat rozen in het algemeen mooi zijn, niet langer louter een esthetisch oordeel, maar een logisch oordeel dat esthetisch is gefundeerd.

Ook het oordeel dat de roos (in haar gebruik) aangenaam is, is weliswaar een esthetisch en enkelvoudig oordeel, maar geen smaakoordeel, doch een oordeel van de zintuigen. Het verschilt van het eerste doordat het smaakoordeel een esthetische kwantiteit van universaliteit met zich meebrengt, dat wil zeggen: geldigheid voor iedereen, die in het oordeel over het aangename niet wordt aangetroffen. Alleen oordelen over het goede hebben, hoewel zij eveneens de voldoening aan een object bepalen, logische en niet louter esthetische geldigheid, omdat zij als kennissen van het object voor iedereen geldig zijn.

Wanneer men objecten louter volgens begrippen beoordeelt, gaat elke voorstelling van schoonheid verloren. Daarom kan er ook geen regel bestaan volgens welke iemand zou kunnen worden gedwongen iets als schoon te erkennen. Of een kledingstuk, een huis of een bloem mooi is: niemand laat zich in zijn oordeel daarover overtuigen door gronden of principes. Men wil het object aan zijn eigen ogen voorleggen, alsof zijn voldoening van gewaarwording afhing; en toch, wanneer men het object dan mooi noemt, meent men met een universele stem te spreken en maakt men aanspraak op de instemming van iedereen, terwijl een louter privé-gewaarwording alleen voor hemzelf en zijn eigen voldoening beslissend zou zijn.

Hieruit blijkt dat in het smaakoordeel niets anders wordt gepostuleerd dan een dergelijke universele stem met betrekking tot voldoening, zonder bemiddeling van begrippen, dus de mogelijkheid van een esthetisch oordeel dat tegelijk als voor iedereen geldig kan worden beschouwd. Het smaakoordeel postuleert niet zelf de instemming van iedereen (dat kan alleen een logisch universeel oordeel doen, omdat het daarvoor gronden kan aanvoeren); het schrijft deze instemming slechts aan iedereen toe, als een geval van de regel, waarvan het bevestiging verwacht, niet uit begrippen, maar alleen uit de instemming van anderen. Deze universele stem is dus slechts een idee (waarop zij berust, zal hier nog niet worden onderzocht).

Of iemand die meent een smaakoordeel te vellen dit oordeel ook daadwerkelijk overeenkomstig deze idee vormt, kan onzeker zijn. Maar dát hij het erop betrekt, en dus dat het een smaakoordeel behoort te zijn, maakt hij kenbaar door de uitdrukking schoon. Daarvan kan hij voor zichzelf zeker zijn door het loutere bewustzijn dat hij alles wat tot het aangename en het goede behoort van de voldoening die hem overblijft heeft afgescheiden. En dát is alles waarvoor hij zich de instemming van iedereen belooft: een aanspraak die hij onder deze voorwaarden ook terecht zou mogen maken, als hij daar niet zo vaak tegen zou zondigen en daardoor een verkeerd smaakoordeel zou vellen.

\begin{mdframed}[leftmargin=0pt,rightmargin=0pt,backgroundcolor=blue!6,linecolor=white,innertopmargin=6pt,innerbottommargin=6pt]

In §8 probeert Kant duidelijk te maken wat voor soort “algemeenheid” het smaakoordeel eigenlijk heeft. Wanneer we iets mooi noemen, doen we alsof ons oordeel voor iedereen geldt, maar die geldigheid is niet objectief of logisch. Ze berust niet op een begrip van het object, zoals bij oordelen over het goede, maar ook niet op een louter privégevoel, zoals bij het aangename. Het is een bijzondere, subjectieve vorm van algemeenheid.

Kant wijst erop dat dit vanuit logisch oogpunt weinig opvalt, maar filosofisch gezien zeer opmerkelijk is. In een smaakoordeel schrijven we ons gevoel van genoegen toe aan iedereen, zonder dat we daarvoor een regel, definitie of concept kunnen geven. Als dit element van aanspraak op algemene instemming zou ontbreken, zouden we het woord “mooi” helemaal niet gebruiken; dan zouden we alles wat ons bevalt simpelweg “aangenaam” noemen, iets waarbij ieder voor zichzelf mag spreken en geen instemming van anderen verwacht.

Daarom onderscheidt Kant twee vormen van smaak. De “smaak van de zintuigen” levert alleen particuliere oordelen op: wat mij bevalt, bevalt mij, en daar houd ik het bij. De “smaak van reflectie” daarentegen doet alsof haar oordeel algemeen geldig is en richt zich tot iedereen. Beide zijn esthetische oordelen, omdat ze gaan over genoegen en ongenoegen, maar alleen de tweede maakt aanspraak op een soort publieke geldigheid. Opmerkelijk is nu dat mensen bij zintuiglijke voorkeuren vanzelf bescheiden zijn, terwijl bij oordelen over schoonheid die bescheidenheid ontbreekt, ook al blijkt in de praktijk vaak dat anderen het niet eens zijn.

Die algemene geldigheid van het smaakoordeel is echter niet logisch. Ze zegt niets over het object zelf, maar alleen iets over hoe het object zich verhoudt tot het gevoel van genoegen. Kant noemt dit een “gemeenschappelijke geldigheid”: geen geldigheid voor alle objecten onder een begrip, maar een geldigheid die zich uitstrekt tot alle beoordelaars. Juist omdat het oordeel niet op een begrip berust, kan men er geen logisch universeel oordeel uit afleiden.

Smaakoordelen zijn daarom altijd enkelvoudig: ze hebben betrekking op dit concrete object, zoals deze roos die ik mooi noem. Pas wanneer men zulke oordelen vergelijkt en veralgemeent, ontstaat er een logisch oordeel, bijvoorbeeld dat rozen in het algemeen mooi zijn. Dat laatste is geen zuiver esthetisch oordeel meer, maar een logisch oordeel dat op esthetische ervaring steunt. Het oordeel dat iets aangenaam is, blijft daarentegen altijd particulier en mist die aanspraak op algemeenheid.

Omdat schoonheid niet via begrippen kan worden vastgesteld, bestaat er ook geen regel die iemand kan dwingen iets mooi te vinden. Men moet het object zelf zien en beoordelen. Toch meent degene die het object mooi noemt, dat hij met één stem spreekt namens iedereen. Die “algemene stem” is geen feit en geen bewijsbare regel, maar een idee: een verwachting van instemming zonder conceptuele onderbouwing. Dat idee hoort wezenlijk bij het smaakoordeel, ook al blijkt in de praktijk vaak dat we ons daarin vergissen.

\bigskip

Een smaakoordeel voelt normatief omdat het niet blijft steken in een beschrijving van een persoonlijk gevoel, maar zich uitspreekt alsof het voor iedereen zou moeten gelden. Wie iets mooi noemt, zegt niet: “dit bevalt mij”, maar spreekt in een vorm die anderen aanspreekt en meeneemt. Die normativiteit zit niet in een regel of bewijs, maar in de vorm van het oordeel zelf: het presenteert zich als iets dat gedeeld kán en zou móéten worden.

Belangrijk is dat dit geen dwang inhoudt. Kant benadrukt dat niemand kan worden overtuigd of overhaald met argumenten om iets mooi te vinden. De normativiteit van het smaakoordeel is zacht en impliciet: zij uit zich als een verwachting van instemming, niet als een eis die afdwingbaar is. Evenmin is het een kwestie van meerderheid of statistiek. Dat veel mensen iets mooi vinden, maakt het oordeel niet geldiger; het oordeel draagt zijn aanspraak al in zich, nog vóór enige feitelijke overeenstemming.

Die bijzondere normativiteit ontstaat juist doordat het smaakoordeel geen beroep doet op begrippen. Omdat ik geen regel kan geven waarom iets mooi is, kan ik ook geen uitzondering formuleren: ik kan niet zeggen voor mij geldt dit, voor anderen misschien niet. Het oordeel staat daardoor meteen in een publieke ruimte, ook al komt het voort uit een subjectief gevoel. Het is precies deze combinatie – subjectief van oorsprong, algemeen van strekking – die het normatieve karakter verklaart.

In §8 wordt duidelijk dat het smaakoordeel altijd al relationeel is. Wie oordeelt, richt zich impliciet tot anderen en veronderstelt een gedeeld vermogen tot voelen en beoordelen. In die zin is het esthetisch oordeel een oefening in het aanspreken van een gemeenschappelijkheid die niet vastligt, maar telkens opnieuw moet blijken. We oefenen ons in het spreken namens een “wij” dat niet vooraf gegarandeerd is.

Dat maakt esthetisch oordelen tot een kwetsbare vorm van samen mens-zijn. Het veronderstelt vertrouwen in een gedeelde gevoeligheid, zonder dat die ooit volledig bevestigd kan worden. Wanneer instemming uitblijft, is dat geen weerlegging, maar een aanleiding tot gesprek, herbezinning of verfijning van het eigen oordeel. Zo bezien is esthetisch oordelen geen afsluiten van meningsverschillen, maar juist een manier om ze betekenisvol te maken.

Het normatieve gevoel dat bij schoonheid hoort, wijst dan niet op een verborgen wet, maar op een menselijke behoefte aan afstemming: de hoop dat ons gevoel niet louter privé is, maar in principe deelbaar. In die zin laat §8 zien dat esthetische oordelen niet alleen over objecten gaan, maar ook over de mogelijkheid van gemeenschap – een gemeenschap die niet wordt opgelegd, maar telkens opnieuw wordt uitgeprobeerd.

\end{mdframed}

\section{Onderzoek van de vraag of in het smaakoordeel het gevoel van genoegen voorafgaat aan het oordelen over het object, of omgekeerd.}

De oplossing van dit probleem is de sleutel tot de kritiek van de smaak en verdient daarom de volle aandacht.

Als het genoegen aan het gegeven object het eerst zou komen, en men in het smaakoordeel slechts de algemene mededeelbaarheid van dit genoegen aan de voorstelling van het object zou toeschrijven, dan zou deze gang van zaken innerlijk tegenstrijdig zijn. Want een dergelijk genoegen zou niets anders zijn dan louter zintuiglijke aangenaamheid, en zou daarom naar zijn aard slechts particuliere geldigheid kunnen hebben, omdat het onmiddellijk afhankelijk is van de voorstelling waardoor het object wordt \textbf{gegeven}.

Daarom moet de algemeen mededeelbare gesteldheid van de geest bij de gegeven voorstelling, als subjectieve voorwaarde van het smaakoordeel, als grond dienen, en moet het genoegen aan het object daaruit volgen. Nu kan echter niets algemeen mededeelbaar zijn behalve kennis en voorstellingen, voor zover zij tot kennis behoren. Want alleen in zoverre is de voorstelling objectief, en alleen daardoor heeft zij een algemeen referentiepunt waaraan ieders voorstellingsvermogen zich moet aanpassen.
Wanneer dus de bepalende grond van het oordeel wordt gedacht als deze algemene mededeelbaarheid van de voorstelling, maar louter subjectief, dat wil zeggen zonder een begrip van het object, dan kan dit niets anders zijn dan de gesteldheid van de geest die ontstaat uit de verhouding van de kenvermogens tot elkaar, voor zover zij een gegeven voorstelling op \textbf{kennis in het algemeen} betrekken.

De kenvermogens die door deze voorstelling in werking worden gezet, verkeren daarbij in een vrij spel, omdat geen bepaald begrip hen aan een specifieke regel van kennis bindt. De geestestoestand bij deze voorstelling moet dus bestaan in een gevoel van het vrije spel van de voorstellingsvermogens bij een gegeven voorstelling, met het oog op kennis in het algemeen.
Bij een voorstelling waardoor een object wordt gegeven, is – om überhaupt kennis mogelijk te maken – de \textbf{verbeeldingskracht} nodig voor het samenstellen van het veelvoud van de aanschouwing, en het \textbf{verstand} voor het begrip dat deze voorstellingen tot een eenheid brengt. Deze toestand van een \textbf{vrij spel} van de kenvermogens bij een voorstelling waardoor een object wordt gegeven, moet algemeen mededeelbaar zijn, omdat kennis – als bepaling van het object waaraan gegeven voorstellingen (in welk subject dan ook) moeten beantwoorden – de enige soort voorstelling is die voor iedereen geldt.

De subjectieve algemene mededeelbaarheid van de wijze van voorstelling in een smaakoordeel, die plaatsvindt zonder vooronderstelling van een bepaald begrip, kan dus niets anders zijn dan de geestestoestand in het vrije spel van verbeelding en verstand (voor zover zij met elkaar overeenstemmen, zoals vereist is voor \textbf{kennis in het algemeen}). Want wij zijn ons ervan bewust dat deze subjectieve verhouding, die geschikt is voor kennis in het algemeen, voor iedereen geldig moet zijn en daarom algemeen mededeelbaar is, net zoals elke bepaalde kennis dat is, die altijd op deze verhouding als haar subjectieve voorwaarde berust.

Dit louter subjectieve (esthetische) oordelen over het object, of over de voorstelling waardoor het object wordt gegeven, gaat vooraf aan het genoegen eraan, en vormt de grond van dit genoegen in de harmonie van de kenvermogens. Alleen op deze universaliteit van de subjectieve voorwaarden van het oordelen over objecten berust de algemene subjectieve geldigheid van de voldoening die wij met de voorstelling van het object verbinden dat wij mooi noemen.

Dat het vermogen om zijn geestestoestand mee te delen – zelfs slechts met betrekking tot de kenvermogens – met genoegen gepaard gaat, kan gemakkelijk empirisch en psychologisch worden vastgesteld uit de natuurlijke neiging van de mens tot sociabiliteit. Maar dat is voor ons doel niet voldoende. Want wanneer wij iets mooi noemen, verwachten wij in het smaakoordeel dit genoegen als noodzakelijk ook bij ieder ander, alsof het een eigenschap van het object zou zijn die daarin volgens begrippen wordt bepaald; terwijl schoonheid op zichzelf niets is, los van de betrekking tot het gevoel van het subject. De bespreking van deze kwestie moeten wij echter uitstellen tot wij eerst een andere vraag hebben beantwoord: hoe en of esthetische oordelen a priori mogelijk zijn.

Voorlopig houden wij ons bezig met de meer beperkte vraag: op welke wijze worden wij ons bewust van een wederzijdse subjectieve overeenstemming van de kenvermogens in het smaakoordeel – esthetisch, door louter innerlijke zin en gevoel, of intellectueel, door het bewustzijn van onze doelgerichte activiteit waardoor wij deze vermogens in werking stellen?

Als de gegeven voorstelling die het smaakoordeel uitlokt een begrip zou zijn, dat verstand en verbeelding in het oordelen over het object verenigt tot kennis van het object, dan zou het bewustzijn van deze verhouding intellectueel zijn (zoals bij het objectieve schematisme van het oordeelsvermogen, dat in de kritiek is behandeld). Maar in dat geval zou het oordeel geen betrekking hebben op genoegen of ongenoegen, en zou het dus geen smaakoordeel zijn.

Het smaakoordeel daarentegen bepaalt het object, onafhankelijk van begrippen, met betrekking tot voldoening en het predicaat van schoonheid. Daarom kan deze subjectieve eenheid van de verhouding zich alleen door middel van gevoel kenbaar maken. De bezieling van beide vermogens (verbeelding en verstand) tot een onbepaalde, maar door de prikkel van de gegeven voorstelling toch eensgezinde activiteit – namelijk die welke tot kennis in het algemeen behoort – is het gevoel waarvan het smaakoordeel de algemene mededeelbaarheid postuleert.

Een objectieve verhouding kan natuurlijk alleen gedacht worden; maar voor zover zij subjectief is wat haar voorwaarden betreft, kan zij toch in haar werking op de geest worden gevoeld. En bovendien is, in het geval van een verhouding die niet op een begrip berust (zoals die van de voorstellingsvermogens tot een kenvermogen in het algemeen), geen ander bewustzijn daarvan mogelijk dan door het gevoel van haar werking, dat bestaat in het vergemakkelijkte spel van beide geestvermogens (verbeelding en verstand), bezield door hun wederzijdse overeenstemming.

Een voorstelling die, hoewel enkelvoudig en zonder vergelijking met andere, toch overeenstemt met de voorwaarden van universaliteit – een overeenstemming die het werk van het verstand in het algemeen uitmaakt – brengt de kenvermogens in die welafgewogen gesteldheid die wij voor alle kennis nodig hebben en die wij daarom ook als geldig beschouwen voor iedereen (voor ieder mens) die door verstand en zintuiglijkheid in samenhang wordt bepaald om te oordelen.

\begin{mdframed}[leftmargin=0pt,rightmargin=0pt,backgroundcolor=blue!6,linecolor=white,innertopmargin=6pt,innerbottommargin=6pt]
In een smaakoordeel komt het genoegen niet eerst, om daarna pas als algemeen geldig te worden opgevat. Als dat zo was, zou het oordeel slechts over zintuiglijke aangenaamheid gaan, en die is altijd privé. Een dergelijk genoegen kan nooit aanspraak maken op algemene geldigheid.

Wat wél eerst komt, is een bepaalde gesteldheid van de geest die algemeen mededeelbaar is. Deze geestestoestand ontstaat wanneer onze kenvermogens – verbeelding en verstand – in een vrij spel met elkaar samenwerken, zonder dat een vast begrip hen stuurt. Dat vrije samenspel is gericht op kennis in het algemeen, maar levert nog geen concrete kennis op.

Omdat kennis in principe voor iedereen geldig moet zijn, is ook deze onderliggende geestestoestand algemeen mededeelbaar. Het genoegen dat wij bij het schone ervaren, volgt pas als gevolg van deze harmonie tussen de kenvermogens.

Het smaakoordeel is dus esthetisch en subjectief: het bepaalt het object niet door begrippen, maar door het gevoel van voldoening. Toch claimt dit oordeel algemene geldigheid, omdat het berust op voorwaarden die bij ieder mens gelden.

Deze overeenstemming tussen verbeelding en verstand wordt niet intellectueel herkend (zoals bij kennis), maar gevoeld. Het gevoel dat hierbij optreedt, is het teken van een goed afgestemde samenwerking van onze kenvermogens, zoals die voor alle kennis nodig is. Daarom beschouwen wij deze toestand als geldig voor iedereen, en verbinden wij haar met het predicaat “mooi”.

\bigskip
In §9 zegt Kant expliciet dat het esthetisch genoegen niet de oorzaak is van het smaakoordeel, maar het gevolg ervan. Wat voorafgaat, is een bepaalde gesteldheid van onze geest: een harmonisch samenspel van verbeelding en verstand, dat niet door een begrip wordt vastgelegd. Deze harmonie is geen kennis, maar ook geen louter zintuiglijk effect. Zij is de structurele voorwaarde waaronder een object als mooi kan verschijnen.

Deze paragraaf vormt daarmee een beslissend scharnierpunt in het betoog. Tot nu toe heeft Kant laten zien dat smaakoordelen aanspraak maken op algemene geldigheid zonder begrippen; hier laat hij zien hoe dat mogelijk is. De sleutel ligt in het feit dat het vrije spel van de kenvermogens precies die toestand is die ook aan alle kennis voorafgaat. Omdat deze verhouding van verbeelding en verstand constitutief is voor “kennis in het algemeen”, mag zij als algemeen mededeelbaar worden gedacht. Zo ontstaat de brug tussen subjectiviteit en universaliteit: niet via het object, maar via de structuur van ons kennen zelf.

De kernvraag van deze paragraaf is dus niet wat schoonheid is, maar wat er in ons gebeurt wanneer we haar ervaren. Kant beschrijft dit niet psychologisch, alsof hij een innerlijke beleving in kaart brengt, maar transcendentaal: hij zoekt naar de noodzakelijke voorwaarden waaronder zo’n ervaring überhaupt mogelijk is. Het vrije spel is geen toevallige gemoedstoestand, maar een structurele mogelijkheid van het menselijk kenvermogen. Het esthetisch genoegen is het voelbare teken dat deze vermogens elkaar tijdelijk zonder dwang ondersteunen.

Dat verklaart ook waarom dit genoegen normatief aanvoelt. We ervaren niet zomaar een prettige sensatie, maar een toestand waarvan we impliciet aannemen dat ieder mens die zou moeten kunnen delen, omdat zij voortkomt uit vermogens die we als gemeenschappelijk veronderstellen. Schoonheid raakt daarmee aan onze rationaliteit zonder haar te reduceren tot kennis.

De vraag of dit vrije spel herkenbaar is in eigen esthetische ervaring, laat zich voorzichtig beantwoorden. Vaak herkennen we bij schoonheid juist een gevoel van soepelheid, van “kloppen”, van aandacht die niet geforceerd wordt maar ook niet wegzakt. Er is betrokkenheid zonder doel, concentratie zonder inspanning. Precies daarin lijkt Kant het vrije spel te situeren: niet als een innerlijke gebeurtenis die we kunnen aanwijzen, maar als een structuur die zich laat voelen wanneer alles even op zijn plaats valt, zonder dat we weten waarom.
\end{mdframed}

\paragraph{Definitie van het mooie afgeleid uit het tweede moment}

Mooi is wat zonder een begrip algemeen bevalt.

\moment{Derde moment van het smaakoordeel, betreffende de verhouding tot de doeleinden die daarin in aanmerking worden genomen.}

\section{Over doelmatigheid in het algemeen}

Als men zou willen definiëren wat een doel is volgens zijn transcendentale bepalingen (dus zonder iets empirisch te veronderstellen, zoals het gevoel van genot), dan is een doel het object van een begrip, voor zover dit begrip wordt opgevat als de oorzaak van dat object (de reële grond van zijn mogelijkheid); en de causaliteit van een begrip met betrekking tot zijn object is doelmatigheid (forma finalis).
Daarom: waar niet alleen de kennis van een object, maar het object zelf (zijn vorm of zijn bestaan) als effect slechts als mogelijk wordt gedacht door middel van een begrip ervan, daar denkt men een doel. De voorstelling van het effect is hier de bepalende grond van zijn oorzaak en gaat daaraan vooraf. Het bewustzijn van de causaliteit van een voorstelling met betrekking tot de toestand van het subject, namelijk om die toestand in stand te houden, kan in het algemeen worden aangeduid met wat men genoegen noemt; daartegenover staat ongenoegen als die voorstelling die de grond bevat om de toestand van de voorstellingen naar hun tegendeel te bepalen (ze te belemmeren of op te heffen).

Het begeervermogen, voor zover het alleen door begrippen kan worden bepaald, dat wil zeggen: om te handelen overeenkomstig de voorstelling van een doel, zou de wil zijn. Een object, een gemoedstoestand of zelfs een handeling wordt echter ook dan doelmatig genoemd wanneer zijn mogelijkheid niet noodzakelijk de voorstelling van een doel veronderstelt, maar wanneer wij zijn mogelijkheid toch alleen kunnen verklaren en begrijpen door als grond daarvan een causaliteit volgens doeleinden aan te nemen, dat wil zeggen een wil die het overeenkomstig de voorstelling van een bepaalde regel zo heeft ingericht.

Doelmatigheid kan dus bestaan zonder een doel, voor zover wij de oorzaken van deze vorm niet in een wil plaatsen, maar de verklaring van haar mogelijkheid voor onszelf toch alleen denkbaar maken door haar van een wil af te leiden.

Nu hebben wij niet altijd noodzakelijk inzicht door de rede (betreffende haar mogelijkheid) in wat wij waarnemen. Daarom kunnen wij ten minste een doelmatigheid wat de vorm betreft waarnemen, ook zonder haar op een doel te baseren (als inhoud van de doelverhouding), en haar in objecten opmerken, zij het op geen andere wijze dan door reflectie.

\begin{mdframed}[leftmargin=0pt,rightmargin=0pt,backgroundcolor=blue!6,linecolor=white,innertopmargin=6pt,innerbottommargin=6pt]
Kant verduidelijkt wat hij onder doel, doelmatigheid en hun onderlinge verhouding verstaat, los van alle empirische aannames zoals gevoel van genot. Een doel is het object van een begrip dat wordt opgevat als de oorzaak van dat object; doelmatigheid is de causaliteit van een begrip met betrekking tot zijn object. Wanneer niet alleen onze kennis van een object, maar het object zelf slechts als mogelijk wordt gedacht door middel van een begrip, spreken we van een doel: het effect wordt dan gedacht als voorafgaand aan en bepalend voor zijn oorzaak.

Het bewustzijn van de werking van een voorstelling op de toestand van het subject kan genoegen heten; het tegendeel daarvan is ongenoegen. Het begeervermogen dat door begrippen wordt bepaald en handelt volgens de voorstelling van een doel, heet de wil.

Een object, toestand of handeling kan echter ook doelmatig genoemd worden zonder dat zijn mogelijkheid noodzakelijkerwijs een doel veronderstelt, zolang wij zijn mogelijkheid slechts kunnen begrijpen door een causaliteit volgens doeleinden aan te nemen, alsof er een wil volgens een regel heeft gehandeld. Daarom kan er doelmatigheid zonder een doel zijn: ook wanneer wij geen werkelijk doel of wil aannemen, kunnen wij de vorm van iets alleen begrijpen door haar als doelmatig te beschouwen.

Ten slotte merkt Kant op dat wij niet altijd inzicht hebben in de mogelijkheid van wat wij waarnemen. Toch kunnen wij doelmatigheid in de vorm van objecten waarnemen, ook zonder deze aan een werkelijk doel te koppelen, en dit gebeurt uitsluitend door reflectie.

\bigskip
Kant stelt hier expliciet dat schoonheid als doelmatig verschijnt: zij toont een orde, samenhang en gepastheid van vorm die wij normaal met doeleinden verbinden. Tegelijk benadrukt hij dat deze doelmatigheid geen werkelijk doel veronderstelt. Er is geen functie die vervuld moet worden, geen intentie die gerealiseerd wordt, geen concept waaronder het object valt. De doelmatigheid ligt niet in wat het object doet of bedoelt, maar in hoe het zich aan ons toont.

Dit is geen bijzaak, maar raakt het hart van de Kritiek van het oordeelsvermogen. Kant zoekt hier een middenpositie tussen mechanistische natuurverklaring en teleologische duiding. In de natuurteleologie zullen we organismen moeten begrijpen alsof zij doelmatig zijn; in de esthetiek ervaren we vormen alsof zij op een doel zijn afgestemd, zonder dat we er een doel aan mogen toeschrijven. In beide gevallen fungeert doelmatigheid niet als kennis, maar als regulatief principe van reflectie. Het esthetisch oordeel is daarmee de “lichte” of vrije variant van een denkbeweging die in de natuurteleologie zwaarder wordt ingezet.

De kernvraag, hoe kan iets doelmatig lijken zonder een doel te hebben?, wordt beantwoord door de verschuiving van intentie naar verschijningsvorm. Doelmatigheid zit hier niet in de oorzaak, maar in de ervaring van samenhang: de delen lijken op elkaar afgestemd, niets lijkt willekeurig, alles “klopt”. Dat dit als doelmatigheid wordt ervaren, zegt primair iets over de verhouding tussen het object en onze kenvermogens, niet over het object als product van een wil. “Zonder doel” betekent dus niet zinloos, maar: zonder extern opgelegd eindpunt.

Juist daarom raakt deze formule zo diep aan religie en symboliek. Wat doelmatig lijkt zonder doel, roept vanzelf de gedachte op aan betekenis zonder dat die wordt vastgelegd. Het opent een ruimte waarin vorm zinvol is zonder dat die zin wordt uitgelegd, waarin orde ervaren wordt zonder dat zij verklaard wordt. Dat is precies de ruimte waarin symboliek en religieuze ervaring zich bewegen: het gevoel dat iets “meer” betekent, zonder dat dit meer tot een concept of leer kan worden herleid. Kant ontkent hier geen diepte of betekenis, maar bewaakt streng dat zij niet tot kennis wordt verheven. Daarmee maakt hij esthetische ervaring tot een plek waar betekenis voelbaar is, zonder dogma, en dat verklaart haar blijvende kracht.
\end{mdframed}

\section{Het smaakoordeel heeft niets anders dan de \emph{vorm van de doelmatigheid} van een object (of van de wijze waarop het wordt voorgesteld) als zijn grond.}

Elk doel, wanneer het als grond van voldoening wordt beschouwd, brengt altijd een belang met zich mee, als bepalende grond van het oordeel over het object van het genoegen. Daarom kan geen subjectief doel het smaakoordeel funderen. Maar bovendien kan ook geen voorstelling van een objectief doel, dat wil zeggen van de mogelijkheid van het object zelf volgens beginselen van doelmatige samenhang, en dus geen begrip van het goede, het smaakoordeel bepalen. Want het smaakoordeel is een esthetisch oordeel en geen ken-oordeel; het heeft dus geen betrekking op enig begrip van de constitutie en de innerlijke of uiterlijke mogelijkheid van het object door deze of gene oorzaak, maar uitsluitend op de verhouding van de voorstellingsvermogens tot elkaar, voor zover zij door een voorstelling worden bepaald.

Nu is deze verhouding, bij de bepaling van een object als schoon, verbonden met een gevoel van genoegen dat door het smaakoordeel tegelijk als voor iedereen geldig wordt verklaard; daarom kan noch een aangenaamheid die de voorstelling begeleidt, noch de voorstelling van de volmaaktheid van het object en het begrip van het goede de bepalende grond bevatten. Niets anders dan de subjectieve doelmatigheid in de voorstelling van een object, zonder enig doel (hetzij objectief, hetzij subjectief), dus louter de vorm van doelmatigheid in de voorstelling waardoor een object ons gegeven wordt, voor zover wij ons daarvan bewust zijn, kan de voldoening uitmaken die wij, zonder een begrip, als algemeen mededeelbaar beoordelen, en dus de bepalende grond van het smaakoordeel vormen.

\begin{mdframed}[leftmargin=0pt,rightmargin=0pt,backgroundcolor=blue!6,linecolor=white,innertopmargin=6pt,innerbottommargin=6pt]
Kant zegt hier dat een oordeel over schoonheid niet wordt bepaald door een doel of functie, maar alleen door de vorm waarin iets doelmatig lijkt.

Als iets ons behaagt omdat het een doel dient, of dat doel nu persoonlijk is (het geeft mij plezier) of objectief is (het is nuttig of moreel goed), dan speelt er altijd een belang mee. En zodra er een belang is, kan er geen sprake meer zijn van een zuiver smaakoordeel. Daarom kan schoonheid niet gebaseerd zijn op genot, nut, volmaaktheid of morele goedheid.

Een smaakoordeel is geen kennis-oordeel: het zegt niets over hoe een object is opgebouwd, waar het voor dient of waardoor het mogelijk is. Het zegt alleen iets over hoe onze verbeelding en ons verstand zich tot elkaar verhouden wanneer we het object waarnemen.

Wat ons dan behaagt, is niet een doel, maar de vorm van doelmatigheid: het object lijkt precies zo ingericht dat onze kenvermogens er vrij en harmonieus mee kunnen omgaan, zonder dat er een concreet doel achter zit. Dit gevoel van doelmatigheid zonder doel is wat wij als schoonheid ervaren en waarvan wij menen dat het voor iedereen geldig kan zijn, ook al berust het niet op een begrip.

\bigskip
Kant stelt hier scherp dat alleen de vorm van een object, de wijze waarop het aan ons verschijnt, beslissend is voor een smaakoordeel. Niet wat het object doet, oproept, symboliseert of beoogt, maar hoe het is gevormd in relatie tot onze kenvermogens. Decoratie, charme, effect, ontroering of verleiding zijn niet verboden, maar zij kunnen geen grond zijn voor schoonheid, omdat zij altijd een belang binnensmokkelen: ze willen iets met ons, ze willen ons ergens toe bewegen of aan iets binden.

Dat past precies in Kants bredere betoog. De hele Kritiek van het oordeelsvermogen is een poging om een domein te denken dat noch door natuurwetenschap, noch door moraal wordt gedomineerd. Door schoonheid te funderen in vormelijke doelmatigheid zonder doel, kan hij een criterium formuleren dat niet afhangt van nut, moraal, emotionele manipulatie of sociale conventie. Daarmee ontstaat ook een norm voor kunstkritiek: kunst is niet beter omdat zij harder raakt, meer ontroert of overtuigender verleidt, maar omdat zij een vrije, niet-afgedwongen samenklank van verbeelding en verstand mogelijk maakt.

De kernvraag, waarom vorm essentieel is, raakt precies dit punt. Vorm is datgene waardoor zintuiglijkheid geordend verschijnt zonder vastgelegd te worden door een begrip. Vorm is geen abstractie die het zinnelijke wegfiltert; integendeel, zij is de manier waarop het zinnelijke voor ons toegankelijk wordt zonder te worden gereduceerd tot functie of boodschap. Kant sluit het zintuiglijke dus niet uit, maar onderwerpt het aan reflectie: niet het prikkelende telt, maar het reflecteerbare.

In relatie tot moderne kunst wordt dit spanningsvol. Veel moderne en hedendaagse kunst problematiseert juist vorm, of gebruikt vorm doelbewust om effecten te sorteren: shock, ontregeling, morele urgentie, politieke mobilisatie. Vanuit Kant bezien ontstaat dan de vraag of zulke werken nog aanspraak maken op een zuiver smaakoordeel, of dat zij zich bewegen in het domein van het aangename of het moreel-appellerende. Tegelijk kan Kants nadruk op vorm ook verrassend actueel blijken: juist waar moderne kunst afstand neemt van narratief, symboliek of boodschap, en de toeschouwer confronteert met hoe iets verschijnt, sluit zij nauw aan bij Kants idee van esthetische vrijheid. In die zin blijft zijn vormbegrip geen beperking, maar een blijvende uitdaging.
\end{mdframed}

\section{Het smaakoordeel berust op a-priorigronden.}

Het is volstrekt onmogelijk om a priori een verbinding vast te stellen tussen het gevoel van genoegen of ongenoegen als gevolg en een bepaalde voorstelling (hetzij een gewaarwording, hetzij een begrip) als oorzaak. Zo’n verbinding zou namelijk een causale relatie zijn, en die kan, voor zover het om ervaringsobjecten gaat, altijd alleen a posteriori, dus door ervaring zelf, worden gekend.
Weliswaar hebben wij in de Kritiek van de praktische rede het gevoel van eerbied (als een bijzondere en eigen wijziging van dit gevoel, die niet precies samenvalt met noch het genoegen noch het ongenoegen dat wij aan empirische objecten ontlenen) a priori afgeleid uit algemene morele begrippen. Maar daar konden wij ook de grenzen van de ervaring overschrijden en een beroep doen op een causaliteit die berust op een bovenzinnelijke eigenschap van het subject, namelijk vrijheid. Zelfs daar hebben wij dit gevoel echter niet werkelijk als een gevolg uit het morele idee afgeleid; wat uit dat idee werd afgeleid, was slechts de bepaling van de wil. De gemoedstoestand van een wil die door iets is bepaald, is op zichzelf reeds een gevoel van genoegen en valt daarmee samen; zij volgt er dus niet als een effect uit voort. Dat laatste zou alleen het geval zijn als het begrip van het morele als een goed voorafging aan de bepaling van de wil door de wet, want dan zou het zinloos zijn het met dat begrip verbonden genoegen als louter uit een kennis daarvan af te leiden.

Iets dergelijks geldt voor het genoegen in het esthetische oordeel, met dit verschil dat het hier louter beschouwend is en geen belang bij het object voortbrengt, terwijl het in het morele oordeel praktisch is. Het bewustzijn van de louter formele doelmatigheid in het spel van de kenvermogens van het subject, bij een voorstelling waardoor een object wordt gegeven, is zelf het genoegen, omdat het een bepalende grond bevat voor de activiteit van het subject met betrekking tot de verlevendiging van zijn kenvermogens. Het bevat dus een innerlijke causaliteit (die doelmatig is) ten aanzien van kennis in het algemeen, zonder echter op een bepaalde kennis te zijn beperkt. Het bevat daarmee niets anders dan de vorm van de subjectieve doelmatigheid van een voorstelling in een esthetisch oordeel.

Dit genoegen is op geen enkele wijze praktisch: noch zoals het genoegen dat voortkomt uit de pathologische grond van het aangename, noch zoals dat uit de intellectuele grond van het voorgestelde goede. Toch bezit het een eigen causaliteit, namelijk die van het in stand houden van de toestand van de voorstelling in het gemoed en van de bezigheid van de kenvermogens, zonder verder doel. Wij blijven bij de beschouwing van het schone stilstaan, omdat deze beschouwing zichzelf versterkt en herhaalt. Dat is vergelijkbaar, al is het niet identiek, met de manier waarop wij blijven hangen bij een bekoorlijkheid in de voorstelling van een object, wanneer die onze aandacht telkens opnieuw aantrekt en het gemoed daarbij passief blijft.

\begin{mdframed}[leftmargin=0pt,rightmargin=0pt,backgroundcolor=blue!6,linecolor=white,innertopmargin=6pt,innerbottommargin=6pt]
    \sloppy
Kant stelt hier dat het genoegen in een smaakoordeel niet vooraf wordt veroorzaakt door een object, maar voortkomt uit wat er in ons gebeurt tijdens het waarnemen. Je kunt niet van tevoren, los van ervaring, vastleggen dat een bepaalde voorstelling noodzakelijk plezier zal geven; zulke causale verbanden leer je normaal alleen via ervaring kennen.

Bij het esthetisch oordeel ligt het anders dan bij morele oordelen. In de moraal is het gevoel verbonden met handelen en wil; bij schoonheid gaat het juist om zuivere beschouwing, zonder praktisch doel of belang.

Het genoegen bij schoonheid is identiek aan het bewustzijn van een vrije, harmonische activiteit van onze kenvermogens (verbeelding en verstand). Dat genoegen is dus niet een extra gevolg, maar valt samen met die toestand zelf. Het heeft geen praktisch effect, maar het houdt de geest in die toestand vast: we blijven kijken en beschouwen, omdat dit spel van onze vermogens zichzelf in stand houdt en herhaalt.

Daarom blijven we bij het schone graag stilstaan: de ervaring versterkt zichzelf, zonder dat we iets willen gebruiken, bezitten of bereiken.

\bigskip
Kant stelt hier onmiskenbaar een ideaal van zuiverheid voorop: een smaakoordeel mag niet worden bepaald door zintuiglijk belang, praktisch nut of morele overwegingen. Zodra genot, begeerte, charme of effectbejag meespelen, verliest het oordeel zijn zuiver esthetisch karakter. Het esthetisch genoegen moet samenvallen met het vormmatige vrije spel van de kenvermogens, niet met wat het object voor ons doet of oplevert.
Maar Kant zegt daarmee niet dat zintuiglijkheid moet verdwijnen. Integendeel: zonder zintuiglijke voorstelling is er geen esthetisch oordeel. Wat uitgesloten wordt, is niet de zintuiglijke input, maar de binding aan belang die ervan kan uitgaan.

Deze zuiverheid functioneert als een trans\-cen\-dentaal ideaaltype. Kant zoekt niet naar een beschrijving van hoe mensen feitelijk oordelen, maar naar de voor\-waarden waaronder een oordeel aan\-spraak op algemene geldig\-heid kan maken. Alleen een oordeel dat vrij is van particuliere be\-langen kan die aan\-spraak doen.
Tegelijk erkent Kant impliciet het spanningsveld: in de praktijk zijn onze esthetische oordelen vaak vermengd met genot, voorkeur, culturele codes en morele verwachtingen. De theorie abstraheert hiervan om zichtbaar te maken wat een smaakoordeel als zodanig zou moeten zijn.

Bestaan er zuivere oordelen van smaak? In strikte zin: zelden, misschien nooit volledig. Kant weet dit. Zuiverheid is geen empirisch gegeven, maar een regulatief begrip: een norm waaraan we onze oordelen kunnen toetsen. Het is wat het smaakoordeel wil zijn wanneer het meer claimt dan “dit vind ik prettig”.
Juist doordat zuiverheid zo moeilijk te realiseren is, maakt ze zichtbaar wanneer een oordeel eigenlijk over genot, status, moraal of effect gaat, en niet over schoonheid.

Zuiverheid is bij Kant geen fictie in de zin van een leeg verzinsel, maar ook geen beschrijving van een alledaagse realiteit. Ze is een normatief oriëntatiepunt: iets wat ons oordeel richting geeft en ons dwingt tot zelfkritiek.
Je zou kunnen zeggen: zuiverheid is geen toestand die we bereiken, maar een maatstaf die ons oordeel disciplineert. Ze bewaart de autonomie van het esthetische domein, juist door onbereikbaar te blijven.
\end{mdframed}

\section{Het zuivere smaakoordeel is onafhankelijk van charme en emotie.}

Elk belang bederft het smaakoordeel en ontneemt het zijn onpartijdigheid, vooral wanneer de doelmatigheid niet aan het gevoel van genoegen voorafgaat, zoals bij het belang van de rede, maar er juist op berust, wat altijd gebeurt bij het esthetisch oordeel over iets, voor zover het behaaglijk of pijnlijk is. Oordelen die daardoor worden beïnvloed, kunnen daarom helemaal geen aanspraak maken op algemene instemming, of hoogstens zo weinig aanspraak als er kan worden gemaakt wanneer dergelijke gewaarwordingen tot de bepalende gronden van de smaak behoren. De smaak is altijd nog barbaars wanneer hij voor zijn voldoening de toevoeging van \textbf{charmes} en \textbf{emoties} nodig heeft, laat staan wanneer hij deze tot maatstaf van zijn goedkeuring maakt.

En toch worden charmes niet alleen vaak samen met schoonheid, die eigenlijk alleen de vorm zou moeten betreffen, opgenomen als een bijdrage aan de esthetische algemene instemming, maar worden zij zelfs als schoonheden op zichzelf voorgesteld, zodat de stof van het genoegen voor de vorm wordt gehouden, een misverstand dat, zoals vele andere die toch altijd iets waars als grond hebben, door een zorgvuldige bepaling van deze begrippen kan worden opgeheven.

Een smaakoordeel waarop charme en emotie geen invloed hebben, ook al kunnen deze met het genoegen in het schone samengaan, en dat dus als bepalende grond uitsluitend de doelmatigheid van de vorm heeft, is een \textbf{zuiver smaakoordeel}.

\begin{mdframed}[leftmargin=0pt,rightmargin=0pt,backgroundcolor=blue!6,linecolor=white,innertopmargin=6pt,innerbottommargin=6pt]

Zodra genoegen of afkeer berust op zintuiglijke aantrekkelijkheid of gevoel, kan het oordeel geen aanspraak maken op algemene geldigheid.
Schoonheid betreft eigenlijk alleen de vorm van het object, niet de aangename werking ervan.
Een zuiver smaakoordeel baseert zich daarom uitsluitend op de doelmatigheid van de vorm, ook al kunnen charme en emotie er feitelijk mee samengaan.

\bigskip
Kant stelt in §13 expliciet dat charme (Reiz) en ontroering (Rührung) het smaakoordeel onzuiver maken. Zodra een oordeel steunt op zintuiglijke prikkeling of emotionele beroering, wordt het bepaald door een belang, hoe subtiel ook. Schoonheid mag daarom niet gebaseerd zijn op wat aangenaam treft of affectief meesleept, maar uitsluitend op de vorm van de voorstelling en de daarin ervaren doelmatigheid zonder doel. Alleen een oordeel dat door deze formele doelmatigheid wordt gedragen, kan aanspraak maken op algemene geldigheid.

Binnen de opbouw van de Kritiek van de oordeelskracht functioneert deze paragraaf als een scherpe afbakening van het esthetische domein tegenover sentimentalisme en effectesthetiek. Kant zuivert hier het terrein waarop een transcendentale analyse van smaak mogelijk is. Tegelijk bereidt §13 het onderscheid voor tussen vrije en aanhangende schoonheid (§16): zodra charme, emotie of functie meebepalend zijn, is schoonheid niet langer vrij maar gebonden aan bijkomende voorwaarden.

De kernvraag waarom emotionele of zintuiglijke werking geen grond kan zijn voor een zuiver smaakoordeel, beantwoordt Kant structureel. Dergelijke werking berust op particuliere gevoeligheid en introduceert altijd een belang, namelijk het verlangen naar voortzetting of herhaling van het aangename. Daarmee verplaatst het oordeel zich van de vorm van het object naar het effect op het subject. Een oordeel dat zo tot stand komt, kan niet normatief zijn, omdat het geen beroep doet op een gedeelde structuur van kenvermogens, maar op individuele affectieve respons.

Wanneer schoonheid afhankelijk wordt van effect, ontroering of impact, verliezen we volgens Kant precies datgene wat haar tot schoonheid maakt in strikte zin: haar vermogen om een vrije, belangeloze instemming te vragen. Schoonheid wordt dan gereduceerd tot ervaring, intensiteit of beleving, en daarmee tot iets dat ons overkomt in plaats van iets waarvoor we elkaar redelijkerwijs kunnen aanspreken. Wat verloren gaat, is niet emotie, maar de mogelijkheid van een gemeenschappelijk esthetisch oordeel dat verder reikt dan persoonlijke indruk of momentane ontroering.
\end{mdframed}

\section{Verheldering aan de hand van voorbeelden.}

Esthetische oordelen kunnen, net als theoretische (logische) oordelen, worden verdeeld in empirische en zuivere. De eerste zijn die welke aangenaamheid of onaangenaamheid vaststellen, de tweede die welke de schoonheid van een object of van de wijze waarop het wordt voorgesteld vaststellen; de eersten zijn zintuiglijke oordelen (materiële esthetische oordelen), de laatsten (als formele) zijn als enigen eigenlijke oordelen van smaak.

Een smaakoordeel is dus slechts zuiver voor zover er geen louter empirische voldoening in zijn bepalingsgrond wordt vermengd. Dit gebeurt echter altijd wanneer charme of ontroering enig aandeel heeft in het oordeel waarmee iets mooi wordt verklaard.

Nu kunnen hier vele bezwaren rijzen, die beweren dat charme niet slechts een noodzakelijk bestanddeel van schoonheid is, maar zelfs geheel op zichzelf voldoende om mooi genoemd te worden. Een loutere kleur, bijvoorbeeld het groen van een grasveld, een loutere toon, onderscheiden van geluid en geruis, bijvoorbeeld die van een viool, wordt door de meeste mensen op zichzelf al mooi genoemd, hoewel beide slechts de materie van de voorstellingen, namelijk loutere gewaarwording, tot grond lijken te hebben, en daarom eigenlijk alleen aangenaam genoemd zouden mogen worden. Toch zal men tegelijk zeker opmerken dat de gewaarwordingen van kleur zowel als van toon terecht slechts dan als mooi gelden, voor zover zij \textbf{zuiver} zijn, een bepaling die reeds de vorm betreft, en die bovendien het enige is wat met zekerheid algemeen kan worden meegedeeld over deze voorstellingen: want van de kwaliteit van de gewaarwordingen zelf kan niet worden aangenomen dat zij bij alle subjecten overeenstemmen, en evenmin kan gemakkelijk worden aangenomen dat de aangenaamheid van de ene kleur boven de andere, of van de toon van het ene muziekinstrument boven die van een ander, door iedereen op dezelfde wijze zal worden beoordeeld.

Indien men met Euler aanneemt dat kleuren trillingen (\emph{pulsus}) van de lucht zijn die elkaar onmiddellijk opvolgen, zoals tonen trillingen van de door geluid verstoorde lucht zijn, en, wat het belangrijkste is, dat de geest niet slechts door de zintuigen hun uitwerking op de prikkeling van het orgaan waarneemt, maar ook door reflectie het regelmatige spel van de indrukken waarneemt, en dus de vorm in de combinatie van verschillende voorstellingen, waaraan ik zeer weinig twijfel, dan zouden kleuren en tonen niet louter gewaarwordingen zijn, maar reeds een formele bepaling van de eenheid van een veelheid ervan, en in dat geval zouden zij ook als schoonheden op zichzelf kunnen gelden.

De zuiverheid van een eenvoudige soort gewaarwording betekent echter dat haar gelijkmatigheid door geen enkele vreemde gewaarwording wordt verstoord of onderbroken, en uitsluitend tot de vorm behoort; want in dat geval kan men abstraheren van de kwaliteit van die soort gewaarwording, van de vraag of en welke kleur of of en welke toon zij voorstelt. Daarom worden alle eenvoudige kleuren, voor zover zij zuiver zijn, als mooi beschouwd; gemengde kleuren hebben dit voordeel niet, omdat men juist doordat zij niet eenvoudig zijn, geen maatstaf heeft om te oordelen of zij zuiver of onzuiver genoemd moeten worden.

Wat betreft de opvatting dat de schoonheid die aan een object op grond van zijn vorm wordt toegeschreven, best kan worden verhoogd door charme, dit is een algemene dwaling en een dwaling die zeer schadelijk is voor een echte, onbedorven en goed gefundeerde smaak, hoewel charmes zeker naast de schoonheid kunnen worden toegevoegd om de geest door de voorstelling van het object boven de droge voldoening uit te interesseren, en zo te dienen om smaak en haar cultivering aan te bevelen, vooral wanneer deze nog ruw en ongeoefend is. Maar zij schaden het smaakoordeel daadwerkelijk wanneer zij de aandacht op zichzelf vestigen als gronden voor het oordelen over schoonheid. Want het is er zo ver van dat zij tot de smaak bijdragen, dat zij, wanneer de smaak nog zwak en ongeoefend is, veeleer slechts met voorzichtigheid moeten worden toegelaten, als vreemdelingen, en alleen voor zover zij die mooie vorm niet verstoren.

In de schilderkunst en de beeldhouwkunst, ja in alle beeldende kunsten, in de architectuur en de tuinbouw voor zover deze schone kunsten zijn, is de \textbf{tekening} datgene wat essentieel is, waarin datgene wat de grond vormt van alle schikkingen voor de smaak niet is wat in de gewaarwording behaagt, maar enkel wat door zijn vorm behaagt. De kleuren die de omtrek verlevendigen behoren tot de charme; zij kunnen het object op zichzelf voor de zintuigen verlevendigen, maar zij kunnen het niet waardig maken om aanschouwd en mooi te zijn, ja zij worden vaak zelfs aanzienlijk beperkt door wat door de mooie vorm wordt vereist, en zelfs waar charme is toegestaan, wordt zij alleen door deze laatste veredeld.

Alle vorm van de objecten van de zintuigen, zowel van de uiterlijke als middellijk ook van de innerlijke, is ofwel \textbf{gestalte} ofwel \textbf{spel}: in het laatste geval ofwel spel van gestalten, in de ruimte, in mimiek en dans, ofwel louter spel van gewaarwordingen, in de tijd. De \textbf{charme} van kleuren of van de aangename tonen van instrumenten kan worden toegevoegd, maar de \textbf{tekening} in het eerste geval en de compositie in het tweede vormen het eigenlijke object van het zuivere smaakoordeel; en dat de zuiverheid van kleuren zowel als van tonen, evenals hun veelheid en hun contrast, tot de schoonheid schijnen bij te dragen, betekent niet dat zij als het ware een aanvulling van gelijke rang leveren aan de voldoening in de vorm doordat zij op zichzelf aangenaam zijn, maar veeleer dat zij deze slechts nauwkeuriger, bepaalder en vollediger aanschouwelijk maken, en bovendien de voorstelling door hun charme verlevendigen, en zo de aandacht voor het object zelf wekken en in stand houden.

Zelfs wat men \textbf{ornamenten} (\emph{parerga}) noemt, dat wil zeggen datgene wat niet innerlijk tot de volledige voorstelling van het object als bestanddeel behoort, maar er slechts uiterlijk als toevoeging bij hoort en de voldoening van de smaak vermeerdert, doet dit nog steeds alleen door zijn vorm: zoals de lijsten van schilderijen, de draperieën van standbeelden of de colonnades rond prachtige gebouwen. Maar als het ornament zelf niet uit mooie vorm bestaat, als het, zoals een vergulde lijst, slechts wordt aangebracht om door zijn charme instemming met het schilderij aan te bevelen -- dan heet het textbf{versiering}, en doet het afbreuk aan de echte schoonheid.

\textbf{Ontroering}, een gewaarwording waarin aangenaamheid slechts wordt voortgebracht door middel van een ogenblikkelijke remming gevolgd door een sterkere uitstorting van de levenskracht, behoort helemaal niet tot de schoonheid. Het verhevene, waarmee het gevoel van ontroering is verbonden, vereist daarentegen een andere maatstaf voor het oordelen dan die waarop de smaak berust; en zo heeft een zuiver smaakoordeel noch charme noch ontroering, met één woord geen gewaarwording, als materie van het esthetische oordeel, tot zijn bepalingsgrond.

\begin{mdframed}[leftmargin=0pt,rightmargin=0pt,backgroundcolor=blue!6,linecolor=white,innertopmargin=6pt,innerbottommargin=6pt]
    
In deze passage onderzoekt Kant het onderscheid tussen empirische en zuivere esthetische oordelen. Net als bij theoretische oordelen kan men ook bij esthetische oordelen een scheiding aanbrengen tussen oordelen die gebaseerd zijn op louter zintuiglijke aangenaamheid of onaangenaamheid, en oordelen die betrekking hebben op de schoonheid van een object of van de wijze waarop het wordt voorgesteld. Een smaakoordeel is slechts zuiver wanneer het niet wordt beïnvloed door empirische gevoelens zoals bekoring of ontroering. Zodra dergelijke gevoelens meespelen, verliest het oordeel zijn zuivere esthetische karakter.

Kant bespreekt vervolgens de wijdverbreide opvatting dat bepaalde kleuren of tonen uit zichzelf mooi zouden zijn. Hoewel mensen bijvoorbeeld het groen van een grasveld of de klank van een viool vaak als mooi beschouwen, zijn dit volgens hem in feite slechts aangename sensaties. Toch merkt hij op dat men kleuren en tonen alleen mooi noemt wanneer zij zuiver zijn, en zuiverheid is een eigenschap die tot de vorm behoort. Daarmee verschuift de waardering van het materiële aspect van de gewaarwording naar een formeel aspect dat voor iedereen toegankelijk kan zijn. Kant verwijst in dit verband naar de theorie van Euler, die kleuren en tonen opvat als regelmatige trillingen. Indien de geest niet alleen de sensatie waarneemt, maar ook de regelmaat van deze trillingen, dan bevatten kleuren en tonen een formeel element dat hen tot mogelijke objecten van een smaakoordeel maakt.

Vervolgens benadrukt Kant dat de vorm van een object het enige is dat een zuiver smaakoordeel kan bepalen. Bekoring kan wel interesse wekken of het object aantrekkelijker maken, maar draagt niet bij aan het smaakoordeel en kan het zelfs schaden, vooral wanneer de smaak nog onontwikkeld is. In de beeldende kunsten, waaronder schilderkunst, beeldhouwkunst, architectuur en tuinarchitectuur, is daarom de tekening of compositie het wezenlijke element. Kleur kan het object levendiger maken, maar bepaalt niet de schoonheid ervan en moet zich vaak voegen naar de eisen van de vorm.

Kant maakt ook een onderscheid tussen ornamenten en decoratie. Ornamenten zijn elementen die niet tot de innerlijke voorstelling van het object behoren, maar deze wel kunnen vergezellen en de smaak kunnen verhogen, mits zij zelf een mooie vorm bezitten. Decoratie daarentegen bestaat uit toevoegingen die slechts door hun bekoring willen behagen en daardoor de aandacht afleiden van de vorm. Zulke toevoegingen verminderen volgens Kant de echte schoonheid van het object.

Tot slot stelt Kant dat emotie, die ontstaat door een tijdelijke onderbreking en daaropvolgende versterking van de levensenergie, niet tot de sfeer van schoonheid behoort. Het sublieme, dat wel met emotie gepaard gaat, vereist een geheel andere wijze van beoordelen dan die van de smaak. Daarom heeft een zuiver smaakoordeel geen sensatie, noch bekoring noch emotie, als grondslag, maar uitsluitend de vorm van het object.

\bigskip
Kant maakt in deze paragraaf expliciet dat kleuren en klanken weliswaar kunnen behagen, maar dat dit behagen vrijwel altijd berust op louter zintuiglijke prikkeling. Zij behoren tot het domein van het aangename, niet tot dat van de zuivere schoonheid. Een zuiver smaakoordeel kan slechts gebaseerd zijn op vorm, dat wil zeggen op datgene in de voorstelling dat onafhankelijk is van individuele gevoeligheden. Zodra bekoring, emotie of sensatie een rol speelt, wordt het oordeel vermengd met subjectieve neigingen en verliest het zijn aanspraak op algemeen geldige instemming.

Binnen het bredere betoog van de *Kritiek van het oordeelsvermogen* vervult deze passage een precisering van Kants centrale these dat schoonheid niet in de materie van de gewaarwording ligt, maar in de vorm van de voorstelling. Kant scherpt hier de grens aan tussen wat tot de sfeer van de smaak behoort en wat slechts decoratief is. Kleuren, klanken en andere zintuiglijke kwaliteiten kunnen een voorstelling aantrekkelijker maken, maar zij dragen niet bij aan de waardering van de vorm zelf. Daarmee verdedigt Kant de autonomie van het esthetisch oordeel tegenover elke poging om het te funderen op prikkel, emotie of zintuiglijke intensiteit.

De centrale vraag van deze paragraaf betreft de rol die zintuiglijke kwaliteiten mogen spelen in een smaakoordeel. Kant erkent dat zij een voorstelling kunnen verlevendigen en de aandacht kunnen vasthouden, maar hun functie blijft secundair. Zij mogen het object ondersteunen, maar niet bepalen. De zuivere schoonheid ligt in de ordening, de regelmaat, de compositie, kortom in datgene wat door iedereen kan worden ingezien zonder beroep op persoonlijke gevoeligheid. Zintuiglijke kwaliteiten mogen dus aanwezig zijn, maar slechts als begeleidend verschijnsel dat de vorm beter waarneembaar maakt, niet als grond van het oordeel zelf.

Toch voelen wij ons sterk aangetrokken tot kleur, klank en sensatie. Dat komt doordat deze elementen direct inspelen op onze zintuiglijke ontvankelijkheid en onmiddellijk plezier verschaffen. Zij vragen geen reflectie, geen afstand, geen oordeel: zij treffen ons rechtstreeks in onze levenskracht. Waar vorm een beroep doet op ons vermogen tot vrije, belangeloze reflectie, spreken kleuren en klanken onze meest directe gevoeligheid aan. Daarom zijn zij verleidelijk, soms zelfs overweldigend. Maar juist omdat zij zo dicht bij onze subjectieve gesteldheid liggen, kunnen zij volgens Kant nooit de basis vormen van een oordeel dat aanspraak maakt op algemene geldigheid.

\end{mdframed}

\section{Het smaakoordeel is volledig onafhankelijk van het begrip van volmaaktheid.}

\textbf{Objectieve} doelmatigheid kan alleen worden gekend door middel van de verhouding van het veelvoudige tot een bepaald doel, dus alleen door middel van een begrip. Hieruit alleen al wordt duidelijk dat het schone, waarvan het oordeel een louter formele doelmatigheid, dat wil zeggen een doelmatigheid zonder doel, tot grond heeft, geheel onafhankelijk is van de voorstelling van het goede, aangezien deze laatste een objectieve doelmatigheid vooronderstelt, dat wil zeggen de betrekking van het object op een bepaald doel.

Objectieve doelmatigheid is ofwel uiterlijk, dat wil zeggen het \textbf{nut} van het object, ofwel innerlijk, dat wil zeggen zijn \textbf{volmaaktheid}. Dat de voldoening aan een object, op grond waarvan wij het mooi noemen, niet kan berusten op de voorstelling van zijn nut, blijkt voldoende duidelijk uit de twee voorafgaande hoofdafdelingen\footnote{LAH: het eerste en tweede moment van de 'Analytiek van het schone'}, aangezien het in dat geval geen onmiddellijke voldoening aan het object zou zijn, wat juist de wezenlijke voorwaarde is van het oordeel over schoonheid. Maar een objectieve innerlijke doelmatigheid, dat wil zeggen volmaaktheid, komt reeds dichter bij het predicaat van schoonheid, en is daarom zelfs door gezaghebbende filosofen met schoonheid vereenzelvigd, zij het met de voorwaarde \textbf{dat zij verward gedacht wordt}. In een kritiek van de smaak is het van het grootste belang te beslissen of schoonheid werkelijk tot het begrip van volmaaktheid kan worden herleid.

Om objectieve doelmatigheid te beoordelen hebben wij altijd het begrip van een doel nodig, en, indien die doelmatigheid niet een uiterlijke (nut) maar een innerlijke is, het begrip van een innerlijk doel, dat de grond van de innerlijke mogelijkheid van het object bevat. Nu is een doel in het algemeen datgene waarvan het \textbf{begrip} kan worden beschouwd als de grond van de mogelijkheid van het object zelf; dus om een objectieve doelmatigheid in een ding voor te stellen moet het begrip van \textbf{wat voor soort ding het behoort te zijn} voorafgaan; en de overeenstemming van het veelvoudige in het ding met dit begrip, dat de regel levert voor de verbinding van het veelvoudige daarin, is de \textbf{kwalitatieve volmaaktheid} van een ding. \textbf{Kwantitatieve} volmaaktheid, als de volledigheid van een ding in zijn soort, is hiervan geheel onderscheiden en is een louter grootheidsbegrip (van totaliteit), waarin \textbf{wat het ding behoort te zijn} reeds als bepaald wordt gedacht en slechts wordt gevraagd of \textbf{alles} wat daarvoor vereist is ook aanwezig is. Wat formeel is in de voorstelling van een ding, dat wil zeggen de overeenstemming van het veelvoudige met een eenheid, terwijl onbepaald blijft wat het ding behoort te zijn, laat op zichzelf helemaal geen kennis van objectieve doelmatigheid toe, omdat van deze eenheid \textbf{als doel}, van wat het ding behoort te zijn, wordt geabstraheerd, en er niets overblijft dan de subjectieve doelmatigheid van de voorstellingen in de geest van de beschouwer, die een zekere doelmatigheid van de voorstellingsstaat van het subject aanduidt, en daarin een gemak in het vatten van een gegeven vorm door de verbeelding, maar niet de volmaaktheid van enig object, dat hier door geen enkel begrip van een doel wordt gedacht. Wanneer ik bijvoorbeeld in het bos een grasveld aantref waaromheen de bomen in een kring staan, en ik mij daarvoor geen doel voorstel, bijvoorbeeld dat het dient voor een landelijke dans, dan wordt door de loutere vorm niet het minste begrip van volmaaktheid gegeven. Maar een formele \textbf{objectieve} doelmatigheid zonder doel voor te stellen, dat wil zeggen de loutere vorm van een volmaaktheid, zonder enige materie en zonder \textbf{enig begrip} van datgene waarmee zij zou moeten overeenstemmen, zelfs al was het slechts het idee van een wetmatigheid in het algemeen, is een ware tegenspraak.

Nu is het smaakoordeel een esthetisch oordeel, dat wil zeggen een oordeel dat op subjectieve gronden berust, en zijn bepalingsgrond kan geen begrip zijn, en dus ook geen begrip van een bepaald doel. Door schoonheid, als formele subjectieve doelmatigheid, wordt dus geen volmaaktheid van het object gedacht als een veronderstelde formele maar tegelijk ook objectieve doelmatigheid, en het onderscheid tussen de begrippen van het schone en het goede, alsof zij slechts in logische vorm verschilden, waarbij het eerste slechts een verward maar het tweede een duidelijk begrip van volmaaktheid zou zijn, terwijl zij overigens in inhoud en oorsprong identiek waren, vervalt. Want in dat geval zou er geen \textbf{specifiek} verschil tussen beide zijn, en zou een smaakoordeel evenzeer een kennisoordeel zijn als het oordeel waarmee iets goed wordt genoemd, net zoals wanneer de gewone man zegt dat bedrog onrechtvaardig is en zijn oordeel op verwarde beginselen grondt, terwijl de filosoof het op duidelijke grondt, maar beide in wezen identieke beginselen van de rede zijn. Ik heb echter reeds uiteengezet dat een esthetisch oordeel van een geheel eigen aard is en absoluut geen kennis, zelfs geen verwarde, van het object verschaft, wat alleen in een logisch oordeel gebeurt; terwijl het esthetische oordeel daarentegen de voorstelling waardoor een object wordt gegeven uitsluitend op het subject betrekt, en ons geen enkele eigenschap van het object voor ogen stelt, maar alleen de doelmatige vorm in de bepaling van de voorstellingsvermogens die ermee bezig zijn. Het oordeel heet ook esthetisch juist omdat zijn bepalingsgrond geen begrip is, maar het gevoel, van de innerlijke zin, van die overeenstemming in het spel van de vermogens van de geest, voor zover deze slechts gevoeld kan worden. Indien men daarentegen verwarde begrippen en het daarop berustende objectieve oordeel esthetisch zou noemen, dan zou men een verstand hebben dat door de zintuigen oordeelt, of een zintuig dat zijn object door begrippen voorstelt, wat beide innerlijk tegenstrijdig is. Het vermogen van begrippen, of zij nu verward of duidelijk zijn, is het verstand; en hoewel het verstand ook tot het smaakoordeel behoort als esthetisch oordeel, zoals tot alle oordelen, behoort het daar niet toe als vermogen tot kennis van een object, maar als vermogen tot bepaling van het oordeel en zijn voorstelling, zonder begrip, overeenkomstig de verhouding van de voorstelling tot het subject en tot diens innerlijk gevoel, en wel voor zover dit oordeel volgens een algemene regel mogelijk is.

\begin{mdframed}[leftmargin=0pt,rightmargin=0pt,backgroundcolor=blue!6,linecolor=white,innertopmargin=6pt,innerbottommargin=6pt]

In deze paragraaf onderzoekt Kant de verhouding tussen het smaakoordeel en het begrip van volmaaktheid. Hij begint met de vaststelling dat objectieve doelmatigheid slechts kan worden gekend door het verband van een veelheid met een bepaald doel, en dat dit noodzakelijk een concept veronderstelt. Hieruit volgt reeds dat het schone, waarvan het oordeel slechts een louter formele doelmatigheid zonder doel als grond heeft, volledig onafhankelijk is van de voorstelling van het goede, aangezien het goede altijd een objectieve doelmatigheid en dus een bepaald doel veronderstelt.

Kant onderscheidt vervolgens externe doelmatigheid, die betrekking heeft op de bruikbaarheid van een object, en interne doelmatigheid, die betrekking heeft op de volmaaktheid ervan. Dat schoonheid niet op bruikbaarheid kan berusten, is volgens hem evident, omdat een smaakoordeel een onmiddellijke voldoening in het object vereist, terwijl bruikbaarheid altijd een middel-doelrelatie impliceert. Interne doelmatigheid, ofwel volmaaktheid, lijkt dichter bij schoonheid te staan en is daarom door sommige filosofen met schoonheid vereenzelvigd, zij het onder de voorwaarde dat het begrip ervan verward wordt opgevat. Kant benadrukt dat het voor een kritiek van de smaak van groot belang is om te bepalen of schoonheid werkelijk tot het begrip van volmaaktheid kan worden herleid.

Om objectieve doelmatigheid te beoordelen is altijd het concept van een doel nodig. Wanneer het om interne doelmatigheid gaat, moet men bovendien het concept van een intern doel hebben, dat de grond bevat van de innerlijke mogelijkheid van het object. Een doel is immers dat waarvan het concept als grond van de mogelijkheid van het object kan worden beschouwd. Daarom moet men, om objectieve doelmatigheid te kunnen voorstellen, eerst weten wat voor soort ding het object geacht wordt te zijn. De overeenstemming van de veelheid in het object met dit concept vormt de kwalitatieve volmaaktheid ervan. Kwantitatieve volmaaktheid, daarentegen, betreft slechts de volledigheid van een ding binnen zijn soort en is een kwestie van omvang, niet van vorm.

Wanneer men abstraheert van het doel van een object, blijft slechts de formele overeenstemming van de veelheid met een eenheid over, maar deze formele structuur maakt geen kennis van objectieve doelmatigheid mogelijk. Wat dan overblijft is slechts een subjectieve doelmatigheid van de voorstellingen in het gemoed van de beschouwer, namelijk een gemak waarmee de verbeelding een gegeven vorm kan vatten. Dit heeft niets te maken met de volmaaktheid van een object, die altijd door een concept van een doel wordt bepaald. Kant illustreert dit met het voorbeeld van een cirkelvormige open plek in het bos: zonder een doel te veronderstellen, bijvoorbeeld dat zij voor dansen is aangelegd, levert de vorm op zichzelf geen enkel begrip van volmaaktheid op. Een formele objectieve doelmatigheid zonder doel is volgens Kant zelfs een tegenspraak.

Het smaakoordeel is echter een esthetisch oordeel, dat op subjectieve gronden berust en waarvan de bepalende grond geen concept kan zijn. Schoonheid, als formele subjectieve doelmatigheid, houdt daarom geen enkele volmaaktheid van het object in, ook niet in een vermeend formeel-objectieve zin. De opvatting dat het schone slechts een verward begrip van volmaaktheid zou zijn, terwijl het goede een duidelijk begrip daarvan is, verwerpt Kant als ongeldig. Indien dit zo was, zou er geen wezenlijk verschil bestaan tussen het smaakoordeel en het oordeel dat iets goed is, en zou het esthetische oordeel een cognitief oordeel worden. Kant benadrukt echter dat een esthetisch oordeel van geheel eigen aard is en geen enkele kennis van het object verschaft, zelfs geen verwarde. Het verwijst uitsluitend naar de wijze waarop de voorstelling het subject raakt, en niet naar eigenschappen van het object zelf.

Het smaakoordeel wordt esthetisch genoemd omdat zijn grond niet in een concept ligt, maar in een gevoel van innerlijke zintuiglijkheid, namelijk het gevoel van een harmonisch samenspel van de kenvermogens. Indien men verwarde concepten en de daarop gebaseerde objectieve oordelen esthetisch zou noemen, zou men ofwel een verstand dat door de zintuigen oordeelt, ofwel een zintuig dat door concepten voorstelt, moeten aannemen, wat beide onmogelijk is. Het verstand, dat de bron van alle concepten is, speelt in het smaakoordeel wel een rol, maar niet als kenvermogen van het object. Het functioneert slechts als vermogen dat het oordeel bepaalt in overeenstemming met de verhouding van de voorstelling tot het subject en diens innerlijk gevoel, en wel op een wijze die toch volgens een algemene regel mogelijk is.

\bigskip
Kant maakt in deze paragraaf ondubbelzinnig duidelijk dat schoonheid niet kan worden herleid tot volmaaktheid. Volmaaktheid veronderstelt altijd een begrip van wat een object moet zijn: een doel, een functie, een vooraf gegeven maatstaf waaraan het object moet beantwoorden. Schoonheid daarentegen is volgens Kant onafhankelijk van elk dergelijk concept. Zij berust op een louter formele doelmatigheid zonder doel, een ervaring waarin de vorm van het object een vrije harmonie van onze kenvermogens teweegbrengt zonder dat wij hoeven te weten wat het object behoort te zijn. Daarmee scheidt Kant het esthetische domein scherp af van het teleologische: waar volmaaktheid altijd een conceptuele beoordeling vereist, is schoonheid een oordeel dat juist geen concept als grondslag kan hebben.

Deze gedachte past naadloos in het bredere betoog van de *Kritiek van het oordeelsvermogen*. Kant bereidt hier het onderscheid voor tussen vrije schoonheid en aanhangende schoonheid, dat hij in §16 expliciet zal uitwerken. Vrije schoonheid is geheel onafhankelijk van een begrip van het object; aanhangende schoonheid daarentegen veronderstelt wel een idee van wat het object moet zijn, zoals bij een menselijk lichaam of een gebouw. Door in §15 te benadrukken dat schoonheid niet op volmaaktheid kan worden teruggebracht, maakt Kant duidelijk dat het zuivere smaakoordeel uitsluitend betrekking heeft op vrije schoonheid. Tegelijkertijd markeert hij de grens tussen esthetische en teleologische oordelen: teleologische oordelen beoordelen objecten op basis van hun doelmatigheid, terwijl esthetische oordelen geen enkel doel veronderstellen.

Vanuit deze achtergrond wordt begrijpelijk waarom schoonheid niet hetzelfde is als goed ontworpen of perfect functionerend. Een object kan technisch uitstekend functioneren, volledig beantwoorden aan zijn doel en in alle opzichten rationeel geordend zijn, zonder dat het daarom mooi is. Functionele volmaaktheid is een eigenschap die door het verstand wordt vastgesteld op basis van een concept van wat het object moet doen. Schoonheid daarentegen ontstaat wanneer de vorm van het object een vrije, niet-doelgerichte harmonie van verbeelding en verstand teweegbrengt. Een perfect ontworpen machine kan dus volmaakt zijn zonder mooi te zijn, terwijl een object zonder duidelijke functie toch als mooi kan worden ervaren. Schoonheid is geen eigenschap die uit het doel van het object volgt, maar een ervaring die ontstaat in het subject.

Toch verwarren wij schoonheid vaak met efficiëntie of optimalisatie. In het dagelijks leven worden objecten die goed functioneren, strak ontworpen zijn of een duidelijke doelmatigheid vertonen, gemakkelijk als mooi bestempeld. Dat komt doordat doelmatigheid een zekere helderheid en orde in de vorm met zich meebrengt, en deze orde kan de waarneming vergemakkelijken. Maar deze verwarring verhult het wezenlijke verschil dat Kant benadrukt: efficiëntie spreekt het verstand aan, schoonheid het vrije spel van onze vermogens. Wanneer wij schoonheid reduceren tot optimalisatie, verliezen wij uit het oog dat het esthetische oordeel juist daar begint waar het doelgerichte denken ophoudt.
\end{mdframed}

\section{Het smaakoordeel waardoor een object onder de voorwaarde van een bepaald begrip als mooi wordt verklaard, is niet zuiver.} % §16

Er zijn twee soorten schoonheid: vrije schoonheid (\emph{pulchritudo vaga}) of slechts aanhangende schoonheid (\emph{pulchritudo adhaerens}). De eerste veronderstelt geen begrip van wat het object behoort te zijn; de tweede veronderstelt zo’n begrip wel, en de volmaaktheid van het object overeenkomstig dat begrip. De eerste worden genoemd (op zichzelf staande) schoonheden van dit of dat ding; de laatste, als aan een begrip gehecht (voorwaardelijke schoonheid), worden toegeschreven aan objecten die onder het begrip van een bepaald doel staan.

Bloemen zijn vrije natuurlijke schoonheden. Nauwelijks iemand anders dan de botanicus weet wat voor soort ding een bloem behoort te zijn; en zelfs de botanicus, die haar herkent als het voortplantingsorgaan van de plant, schenkt aan dit natuurlijke doel geen aandacht wanneer hij de bloem naar de smaak beoordeelt. Dit oordeel is dus niet gegrond op enige volmaaktheid, enige innerlijke doelmatigheid waaraan de samenstelling van het veelvoudige zou zijn gerelateerd. Veel vogels, de papegaai, de kolibrie, de paradijsvogel, en een menigte zeekreeftachtigen, zijn schoonheden op zichzelf, die niet overeenkomstig begrippen aan een bepaald object met betrekking tot zijn doel zijn gehecht, maar vrij zijn en om zichzelf behagen. Evenzo betekenen ornamenten \emph{à la grecque}, loofwerk voor randen of op behang, enz., op zichzelf niets: zij stellen niets voor, geen object onder een bepaald begrip, en zijn vrije schoonheden. Tot dezelfde soort kan men ook rekenen wat in de muziek fantasieën worden genoemd, zonder thema, ja zelfs alle muziek zonder tekst.

Bij het beoordelen van een vrije schoonheid, naar loutere vorm, is het smaakoordeel zuiver. Er wordt geen begrip van enig doel verondersteld waartoe het veelvoudige het gegeven object zou moeten dienen en dat het dus zou moeten voorstellen, waardoor de verbeelding, die als het ware spelend is in het beschouwen van de gestalte, slechts beperkt zou worden.

Maar de schoonheid van een mens, en in deze soort die van een man, een vrouw of een kind, de schoonheid van een paard, van een gebouw, zoals een kerk, een paleis, een arsenaal of een tuinhuis, veronderstelt een begrip van het doel dat bepaalt wat het ding behoort te zijn, dus een begrip van zijn volmaaktheid, en is dus slechts aanhangende schoonheid. Zoals nu de vermenging van het aangename, van de gewaarwording, met schoonheid, die eigenlijk alleen de vorm betreft, de zuiverheid van het smaakoordeel belemmerde, zo schaadt ook de vermenging van het goede, dat wil zeggen de wijze waarop het veelvoudige goed is voor het ding zelf overeenkomstig zijn doel, met schoonheid, zijn zuiverheid.

Men zou aan een gebouw veel kunnen toevoegen dat in de aanschouwing behaaglijk zou zijn, als het maar niet bedoeld was als kerk; een figuur zou met allerlei krullen en lichte maar regelmatige lijnen verfraaid kunnen worden, zoals de Nieuw-Zeelanders dat doen met hun tatoeëringen, als het maar geen mens was; en deze laatste zou veel fijnere trekken en een aangenamer, zachter verloop van zijn gelaatstrekken kunnen hebben, als hij maar niet geacht werd een man, of zelfs een krijger, voor te stellen.

De voldoening aan het veelvoudige in een ding met betrekking tot het innerlijke doel dat zijn mogelijkheid bepaalt, is een voldoening die op een begrip berust; de voldoening aan schoonheid daarentegen is er een die geen begrip veronderstelt, maar onmiddellijk verbonden is met de voorstelling waardoor het object wordt gegeven, niet waardoor het wordt gedacht. Wanneer nu het smaakoordeel met betrekking tot het laatste afhankelijk wordt gemaakt van het doel in het eerste, als een oordeel van de rede, en daardoor wordt beperkt, dan is het geen vrij en zuiver smaakoordeel meer.

Weliswaar wint de smaak door deze verbinding van esthetische voldoening met het intellectuele, doordat hij vastheid krijgt en, hoewel niet algemeen, regels kan worden voorgeschreven met betrekking tot bepaalde doelmatig bepaalde objecten. Maar in dat geval zijn dit ook geen regels van de smaak, maar slechts regels voor de vereniging van smaak met rede, dat wil zeggen van het schone met het goede, waardoor het eerste bruikbaar wordt als instrument van de bedoeling met betrekking tot het laatste, zodat de gemoedstoestand die zichzelf draagt en van subjectief algemene geldigheid is, ten grondslag kan liggen aan datgene wat alleen door inspannende vastbeslotenheid kan worden volgehouden maar objectief algemeen geldig is. Strikt genomen wint echter de volmaaktheid niets door schoonheid, noch wint schoonheid iets door volmaaktheid; veeleer, aangezien wij bij het vergelijken van de voorstelling waardoor een object ons wordt gegeven met het object, met betrekking tot wat het behoort te zijn, niet kunnen vermijden deze tegelijkertijd met het subject samen te houden, wint het \textbf{gehele vermogen} van de voorstellingsvermogens wanneer beide gemoedstoestanden met elkaar in overeenstemming zijn.

Een smaakoordeel met betrekking tot een object met een bepaald innerlijk doel zou dus alleen dan zuiver zijn wanneer degene die oordeelt ofwel geen begrip van dit doel heeft, ofwel daarvan in zijn oordeel abstraheert. Maar in dat geval zou hij, hoewel hij een juist smaakoordeel heeft geveld, doordat hij het object als een vrije schoonheid beoordeelde, toch door een ander worden bekritiseerd en van een valse smaak beschuldigd, namelijk door iemand die de schoonheid in het object slechts als een aanhangende eigenschap beschouwt, en naar het doel van het object kijkt, hoewel beiden op hun wijze juist oordelen: de een op grond van wat hij voor zijn zintuigen heeft, de ander op grond van wat hij in zijn gedachten heeft.

Door middel van dit onderscheid kan men veel geschillen over schoonheid tussen smaakbeoordelaars beslechten, door hun te laten zien dat de een zich met vrije schoonheid bezighoudt, de ander met aanhangende schoonheid, dat de eerste een zuiver, de tweede een toegepast smaakoordeel velt.

\begin{mdframed}[leftmargin=0pt,rightmargin=0pt,backgroundcolor=blue!6,linecolor=white,innertopmargin=6pt,innerbottommargin=6pt]
In deze paragraaf maakt Kant een fundamenteel onderscheid tussen twee soorten schoonheid: vrije schoonheid en aanhangende schoonheid. Vrije schoonheid veronderstelt geen enkel begrip van wat het object zou moeten zijn; zij wordt beoordeeld zonder voorafgaande kennis van een doel of functie. Aanhangende schoonheid daarentegen veronderstelt wel een begrip van het doel van het object en daarmee ook een idee van zijn volmaaktheid. Vrije schoonheid is schoonheid die aan een object zelf toekomt zonder dat het onder een bepaald concept valt, terwijl aanhangende schoonheid slechts kan worden beoordeeld in relatie tot een vooraf gegeven einddoel.

Kant illustreert vrije schoonheid met voorbeelden uit de natuur en de kunst. Bloemen gelden als vrije natuurlijke schoonheden, omdat men bij hun beoordeling geen rekening houdt met hun biologische functie. Zelfs de botanist, die weet dat de bloem het voortplantingsorgaan van de plant is, laat dit doel buiten beschouwing wanneer hij haar schoonheid beoordeelt. Ook veel vogels en zeedieren zijn volgens Kant schoonheden op zichzelf, omdat hun vorm niet wordt beoordeeld aan de hand van een concept van wat zij zouden moeten zijn. Evenzo behoren ornamentale patronen, zoals Grieks geïnspireerde motieven of bladornamenten, tot de vrije schoonheden, omdat zij niets voorstellen en geen object onder een bepaald begrip representeren. In de muziek geldt hetzelfde voor fantasieën zonder thema en voor muziek zonder tekst.

Bij de beoordeling van vrije schoonheid is het smaakoordeel zuiver, omdat het uitsluitend op de vorm berust en geen concept van een doel veronderstelt. De verbeelding kan zich vrij bewegen in de waarneming van de vorm, zonder door een vooraf gegeven idee te worden beperkt. Anders ligt het bij de schoonheid van mensen, dieren of gebouwen. De schoonheid van een mens, een paard of een kerk veronderstelt altijd een begrip van wat deze objecten zouden moeten zijn. Een mens moet bijvoorbeeld bepaalde proporties hebben om als mens te worden herkend, en een kerk moet aan bepaalde eisen voldoen om haar functie te vervullen. In zulke gevallen is de schoonheid aanhangend, omdat zij afhankelijk is van een concept van volmaaktheid dat uit het doel van het object voortvloeit. Net zoals de vermenging van het aangename met schoonheid de zuiverheid van het smaakoordeel aantast, zo schaadt ook de vermenging van het goede met schoonheid de zuiverheid ervan.

Kant benadrukt dat men veel aan een object zou kunnen toevoegen dat het aangenamer maakt, maar dat deze toevoegingen niet passend zouden zijn wanneer het object een bepaald doel moet vervullen. Een gebouw zou visueel aantrekkelijker kunnen worden gemaakt als het geen kerk hoefde te zijn; een menselijk lichaam zou met sierlijke lijnen kunnen worden versierd als het niet als mens moest worden beoordeeld; en een gezicht zou zachtere en aangenamere trekken kunnen hebben als het niet een man of zelfs een krijger moest voorstellen. Deze voorbeelden tonen aan dat het doelbegrip de beoordeling van schoonheid kan beperken.

De voldoening die men ervaart in de overeenstemming van een object met zijn interne doel is een voldoening die op een concept berust. De voldoening in schoonheid daarentegen is onmiddellijk verbonden met de voorstelling van het object en veronderstelt geen concept. Wanneer het smaakoordeel afhankelijk wordt gemaakt van het doel van het object, wordt het een oordeel van de rede en verliest het zijn vrijheid en zuiverheid. Toch kan de smaak door deze verbinding met het intellectuele aan vastheid winnen, omdat zij dan regels kan krijgen voor objecten die door hun doel bepaald zijn. Maar deze regels zijn geen regels van de smaak zelf; het zijn regels voor de samenhang tussen smaak en rede, waardoor schoonheid dienstbaar kan worden gemaakt aan het goede. Schoonheid en volmaaktheid versterken elkaar echter niet wederzijds; zij blijven onderscheiden domeinen, al kan de geest als geheel winnen wanneer beide toestanden van het gemoed in harmonie zijn.

Een smaakoordeel over een object met een bepaald intern doel kan slechts zuiver zijn wanneer de beoordelaar dit doel niet kent of er in zijn oordeel van abstraheert. In dat geval beoordeelt hij het object als vrije schoonheid. Toch kan iemand anders hem verwijten dat hij een verkeerde smaak heeft, omdat deze ander de schoonheid uitsluitend als aanhangende eigenschap beschouwt en het object beoordeelt in relatie tot zijn doel. Beide oordelen kunnen echter op hun eigen wijze juist zijn: het ene vanuit de onmiddellijke zintuiglijke voorstelling, het andere vanuit het concept van wat het object moet zijn. Door dit onderscheid te maken kan men veel geschillen over schoonheid oplossen, omdat men kan aantonen dat de ene beoordelaar zich richt op vrije schoonheid en de andere op aanhangende schoonheid, waarbij de eerste een zuiver en de tweede een toegepast smaakoordeel velt.

\bigskip
Kant maakt in deze paragraaf duidelijk dat er twee fundamenteel verschillende manieren zijn om schoonheid te beoordelen. Vrije schoonheid ontstaat wanneer wij een object uitsluitend op basis van zijn vorm beoordelen, zonder enig begrip van wat het object zou moeten zijn. Bloemen, ornamenten en abstracte patronen zijn hiervan de voorbeelden: zij vragen geen kennis van hun functie of aard om als mooi te worden ervaren. Aanhangende schoonheid daarentegen veronderstelt dat wij weten wat het object is en welk doel het vervult. Bij mensen, dieren of gebouwen speelt altijd een idee van wat zij behoren te zijn mee, en dit idee beïnvloedt onvermijdelijk het oordeel van schoonheid.

Deze tweedeling past naadloos in Kants bredere betoog in de Kritiek van het oordeelsvermogen. Zij vormt een noodzakelijke voorbereiding op zijn latere analyses van menselijke schoonheid, architectuur en kunst. In al deze domeinen is het onmogelijk om de vorm volledig los te maken van het doel of het concept van het object. Kant moet daarom eerst het onderscheid tussen vrije en aanhangende schoonheid scherp formuleren om te kunnen uitleggen waarom sommige esthetische oordelen zuiver zijn en andere noodzakelijk door begrippen worden gekleurd. Tegelijkertijd markeert dit onderscheid de grens tussen het esthetische en het teleologische: waar het teleologische oordeel altijd een doel veronderstelt, kan het zuivere esthetische oordeel slechts bestaan wanneer het zonder doel wordt geveld.

De centrale vraag van deze paragraaf is wanneer een oordeel van schoonheid kennis van wat iets is vereist. Dat is het geval zodra het object onder een bepaald concept valt dat zijn identiteit en functie bepaalt. Een menselijk lichaam kan niet worden beoordeeld zonder een idee van menselijke proporties; een kerk kan niet worden beoordeeld zonder een begrip van wat een kerk behoort te zijn; een paard kan niet worden beoordeeld zonder kennis van de kenmerken die een paard definiëren. In zulke gevallen is het oordeel van schoonheid onvermijdelijk verbonden met een normatief kader dat bepaalt wat het object moet zijn. De schoonheid wordt dan niet uitsluitend in de vorm gevonden, maar in de mate waarin de vorm overeenkomt met het doel van het object.

Zodra normen, functies of idealen meespelen, verandert het oordeel van schoonheid van karakter. Het wordt minder vrij, minder spontaan en minder zuiver. De verbeelding kan zich niet langer onbelemmerd bewegen, omdat zij wordt teruggefloten door het begrip van wat het object behoort te zijn. De beoordeling wordt gedeeltelijk een rationele beoordeling: men vraagt zich niet alleen af of de vorm aangenaam is, maar ook of zij passend is, doelmatig, in overeenstemming met een ideaal. Hierdoor verschuift het oordeel van een zuiver esthetische ervaring naar een gemengde ervaring waarin zowel smaak als rede een rol spelen. Wat verloren gaat, is de onvoorwaardelijke vrijheid van het zuivere smaakoordeel; wat gewonnen wordt, is een zekere stabiliteit en normativiteit, maar tegen de prijs van het verdwijnen van de puurheid die Kant als essentieel voor schoonheid beschouwt.
\end{mdframed}

\section{Over het ideaal van het schone.} % §17

Er kan geen objectieve regel van de smaak bestaan die door middel van begrippen zou bepalen wat mooi is. Want ieder oordeel uit deze bron is esthetisch, dat wil zeggen: zijn bepalende grond is het gevoel van het subject en niet een begrip van een object. Het zoeken naar een principe van de smaak dat door bepaalde begrippen het universele criterium van het schone zou verschaffen, is een vruchteloze onderneming, omdat wat men zoekt onmogelijk en innerlijk tegenstrijdig is. De universele mededeelbaarheid van de gewaarwording (van voldoening of onvoldoening), en wel een die zonder begrippen plaatsvindt, de zoveel mogelijk eensgezindheid van alle tijden en volkeren omtrent dit gevoel bij de voorstelling van bepaalde objecten: hoewel zwak en nauwelijks voldoende voor een gissing, is dit het empirische criterium voor de afleiding van een smaak uit de gemeenschappelijke, diep in alle mensen verborgen grond van eensgezindheid in het oordelen over vormen waaronder objecten hun gegeven zijn.

Daarom worden sommige producten van de smaak als \textbf{voorbeeldig} beschouwd – niet alsof men smaak zou kunnen verwerven door anderen na te volgen. Want smaak moet een eigen vermogen zijn; wie echter een voorbeeld navolgt, toont, voor zover hij het goed treft, wel vaardigheid, maar hij toont smaak slechts voor zover hij dit voorbeeld zelf kan beoordelen. Hieruit volgt echter dat het hoogste voorbeeld, het archetype van de smaak, slechts een idee is, dat ieder in zichzelf moet voortbrengen en waaraan hij alles moet toetsen wat een object van de smaak is, of wat een voorbeeld is van oordelen door middel van smaak, zelfs de smaak van ieder ander. \textbf{Idee} betekent, strikt genomen, een begrip van de rede, en \textbf{ideaal} de voorstelling van een individueel wezen als adequaat aan een idee. Dat archetype van de smaak, dat weliswaar berust op de onbepaalde idee van de rede van een maximum, maar niet door begrippen kan worden voorgesteld, doch alleen in een individuele presentatie, zou daarom beter het ideaal van het schone genoemd worden, iets wat wij in onszelf trachten voort te brengen, ook al zijn wij er niet in het bezit van. Maar dit zal slechts een ideaal van de verbeelding zijn, juist omdat het niet op begrippen berust maar op voorstelling, en het vermogen van voorstelling is de verbeelding. – Hoe bereiken wij nu zo’n ideaal van schoonheid? \emph{A priori} of empirisch? En eveneens: welke soort schoonheid laat een ideaal toe?

Ten eerste moet worden opgemerkt dat de schoonheid waarvoor een idee gezocht moet worden, geen \textbf{vage} schoonheid mag zijn, maar een schoonheid die \textbf{vastligt} door een begrip van objectieve doelmatigheid; zij mag dus niet behoren tot het object van een geheel zuiver smaak­oordeel, maar tot dat van een gedeeltelijk geintellectualiseerd smaak­oordeel. Dat wil zeggen: bij welk soort beoordelingsgronden men ook een ideaal veronderstelt, daaraan moet een idee van de rede ten grondslag liggen overeenkomstig bepaalde begrippen, die \emph{a priori} het doel bepalen waarop de innerlijke mogelijkheid van het object berust. Een ideaal van mooie bloemen, van mooie meubelstukken, van een mooi landschap kan niet worden gedacht. Maar ook een ideaal van een schoonheid die aan bepaalde doelen hecht, bijvoorbeeld van een mooie woning, een mooie boom, mooie tuinen, enz., kan niet worden voorgesteld, vermoedelijk omdat de doelen niet voldoende bepaald en vastgelegd zijn door hun begrip, en de doelmatigheid daarom bijna even vrij is als bij \textbf{vage} schoonheid. Alleen datgene wat het doel van zijn bestaan in zichzelf heeft, de \textbf{mens}, die zijn doelen zelf door de rede bepaalt, of, waar hij ze uit uiterlijke waarneming moet afleiden, deze toch kan vergelijken met wezenlijke en universele doelen en in dat geval ook esthetisch hun overeenstemming daarmee kan beoordelen: deze \textbf{mens} alleen is tot een ideaal van \textbf{schoonheid} in staat, zoals ook de menselijkheid in zijn persoon, als intelligentie, als enige onder alle objecten in de wereld tot het ideaal van \textbf{volmaaktheid} in staat is.

Hierbij zijn echter twee elementen betrokken: \textbf{ten eerste} de esthetische \textbf{normale idee}, die een individuele aanschouwing (van de verbeelding) is, welke de maatstaf voor het oordelen ervan als een ding behorend tot een bepaalde diersoort voorstelt; \textbf{ten tweede} de \textbf{idee van de rede}, die de doelen van de menselijkheid, voor zover zij niet zintuiglijk kunnen worden voorgesteld, tot het principe maakt voor het oordelen over haar gestalte, waardoor deze, als hun verschijningseffect, zichtbaar worden. De normale idee moet haar elementen voor de gestalte van een dier van een bepaalde soort uit de ervaring halen; maar de grootste doelmatigheid in de bouw van de gestalte, die als universele maatstaf voor het esthetisch oordelen over ieder individu van deze soort zou kunnen dienen, het beeld dat als het ware opzettelijk de techniek van de natuur heeft gegrondvest, waaraan alleen de soort als geheel, maar geen enkel afzonderlijk individu, adequaat is, ligt slechts in de idee van degene die oordeelt, die echter met haar verhoudingen volledig \emph{in concreto} kan worden voorgesteld als een esthetische idee in een modelbeeld. Om enigszins begrijpelijk te maken hoe dit gebeurt (want wie kan dit geheim geheel aan de natuur ontfutselen?), zullen wij een psychologische verklaring proberen.

Men moet opmerken dat de verbeelding niet alleen weet hoe zij ons af en toe, zelfs na lange tijd, tekens van begrippen kan terugroepen op een wijze die ons geheel onbegrijpelijk is; zij weet ook hoe zij het beeld en de gestalte van een object kan reproduceren uit een enorme hoeveelheid objecten van verschillende soorten, of zelfs van één en dezelfde soort; ja, wanneer de geest erop gericht is vergelijkingen te maken, weet zij zelfs, naar alle waarschijnlijkheid werkelijk, zij het niet bewust, als het ware het ene beeld over het andere heen te leggen en door middel van de overeenstemming van meerdere van dezelfde soort tot een gemiddelde te komen dat allen als gemeenschappelijke maat kan dienen. Iemand heeft duizend volwassen mannen gezien. Nu, als hij zou willen oordelen wat als hun betrekkelijk normale grootte moet gelden, dan laat (naar mijn mening) de verbeelding een groot aantal beelden (misschien alle duizend) over elkaar heen leggen, en, als ik hier de analogie van de optische voorstelling mag gebruiken, in de ruimte waar het grootste aantal samenvalt en binnen de omtrek van de plaats die door de meest geconcentreerde kleuren wordt verlicht, daar wordt de \textbf{gemiddelde grootte} herkenbaar, die zowel in hoogte als in breedte even ver verwijderd is van de uiterste grenzen van de grootste en de kleinste gestaltes; en dit is de gestalte van een mooie man. (Men zou hetzelfde resultaat mechanisch kunnen verkrijgen als men alle duizend mannen zou meten, hun hoogten, breedten (en omtrekken) zou optellen en vervolgens de som door duizend zou delen. Maar de verbeelding doet precies dit door middel van een dynamisch effect, dat ontstaat uit de herhaalde apprehensie van zulke gestalten op het orgaan van de innerlijke zin.) Wanneer nu op vergelijkbare wijze bij deze gemiddelde man de gemiddelde kop, de gemiddelde neus, enz. wordt gezocht, dan is deze gestalte de grondslag voor de normale idee van de mooie man in het land waar deze vergelijking wordt gemaakt; daarom moet onder deze empirische voorwaarden een neger noodzakelijkerwijs een andere normale idee van de schoonheid van een gestalte hebben dan een blanke, een Chinees een andere dan een Europeaan. Precies zo is het gesteld met het model van een mooi paard of een mooie hond (van een bepaald ras). – Deze \textbf{normale idee} is niet afgeleid uit de uit ervaring genomen verhoudingen \textbf{als bepaalde regels}; veeleer worden volgens haar eerst regels voor het oordelen mogelijk. Zij is het beeld voor de hele soort, zwevend tussen alle bijzondere en op verschillende wijzen afwijkende aanschouwingen van de individuen, dat de natuur gebruikte als het archetype dat ten grondslag ligt aan haar voortbrengselen in dezelfde soort, maar dat zij in geen enkel individu volledig schijnt te hebben bereikt. Zij is geenszins het volledige \textbf{archetype} van \textbf{schoonheid} in deze soort, maar slechts de vorm die de onontbeerlijke voorwaarde van alle schoonheid uitmaakt, en dus slechts de juistheid in de voorstelling van de soort. Zij is, zoals gezegd is van \textbf{Polycletus’} beroemde \textbf{Doryphoros}, de regel (en \textbf{Myrons} koe zou op dezelfde wijze in haar soort kunnen dienen). Juist daarom kan zij niets specifiek karakteristieks bevatten, want dan zou zij niet de \textbf{normale idee} voor de soort zijn. Haar voorstelling behaagt ook niet wegens schoonheid, maar enkel omdat zij geen enkele voorwaarde tegenspreekt waaronder alleen een ding van deze soort mooi kan zijn. De voorstelling is slechts academisch juist.

Toch bestaat er nog een onderscheid tussen de \textbf{normale idee} van het schone en zijn \textbf{ideaal}, dat op de reeds aangevoerde gronden alleen bij de \textbf{menselijke gestalte} te verwachten is. In dit laatste bestaat het ideaal in de uitdrukking van het \textbf{morele}, zonder welke het object niet universeel, en bovendien positief (niet louter negatief in een academisch juiste voorstelling) zou bevallen. De zichtbare uitdrukking van morele ideeën, die de mens innerlijk besturen, kan natuurlijk alleen uit de ervaring worden geput; maar om als het ware in lichamelijke verschijning (als effect van het innerlijke) hun verbinding zichtbaar te maken met alles wat ons verstand met het moreel goede verbindt in de idee van de hoogste doelmatigheid – goedheid van ziel, of zuiverheid, of kracht, of rust, enz. – daarvoor zijn zuivere ideeën van de rede en een grote kracht van de verbeelding vereist, verenigd in degene die dit slechts zou oordelen, laat staan in degene die het zou voorstellen. De juistheid van zo’n ideaal van schoonheid blijkt daaruit dat geen enkele zintuiglijke bekoring met de voldoening in zijn object vermengd mag worden, terwijl het toch een groot belang toelaat dat men erin stelt, wat dan bewijst dat het oordelen volgens zo’n maatstaf nooit louter esthetisch kan zijn, en het oordelen volgens een ideaal van schoonheid geen louter smaak­oordeel is.

\begin{mdframed}[leftmargin=0pt,rightmargin=0pt,backgroundcolor=blue!6,linecolor=white,innertopmargin=6pt,innerbottommargin=6pt]

Kant opent deze paragraaf met de stelling dat er geen objectieve regel van smaak kan bestaan die door middel van begrippen bepaalt wat mooi is. Een smaakoordeel is immers esthetisch en berust op het gevoel van het subject, niet op een concept van het object. Het zoeken naar een universeel criterium van schoonheid op basis van begrippen is daarom volgens Kant een innerlijk tegenstrijdige onderneming. De enige aanwijzing voor een gemeenschappelijke grond van smaak ligt in de feitelijke overeenstemming van mensen in hun oordelen over bepaalde vormen, een overeenstemming die weliswaar zwak en empirisch is, maar toch wijst op een gedeelde menselijke aanleg om vormen op een bepaalde wijze te beoordelen.

Omdat mensen deze gemeenschappelijke aanleg delen, worden bepaalde uitingen van smaak als maatgevend gezien.
Dit betekent echter niet dat smaak door imitatie kan worden aangeleerd. Imitatie toont slechts vaardigheid; smaak toont zich pas wanneer iemand het voorbeeld zelf kan beoordelen. 
Hieruit volgt dat het hoogste model van smaak geen concreet voorbeeld is, maar een idee dat ieder mens in zichzelf moet voortbrengen. 
Dit idee, dat niet door begrippen kan worden weergegeven maar slechts in een individuele voorstelling, noemt Kant het ideaal van schoonheid. 
Het is een ideaal van de verbeelding, omdat het niet op begrippen berust maar op een voorstelling die de verbeelding voortbrengt.

Kant onderzoekt vervolgens welke soorten schoonheid een dergelijk ideaal toelaten. Een ideaal kan niet behoren tot vage schoonheid, die geen bepaald doel of concept veronderstelt. Het moet gaan om schoonheid die verbonden is met een objectieve doelmatigheid, en dus niet tot een volledig zuiver smaakoordeel behoort, maar tot een gedeeltelijk intellectueel oordeel. Toch kunnen zelfs objecten met een bepaald doel, zoals huizen, bomen of tuinen, geen ideaal van schoonheid hebben, omdat hun doeleinden niet voldoende vastliggen. Alleen de mens, die zijn doelen door de rede bepaalt en deze kan vergelijken met universele doeleinden, kan drager zijn van een ideaal van schoonheid. Zoals alleen de mens drager kan zijn van een ideaal van volmaaktheid, zo kan alleen de menselijke figuur een ideaal van schoonheid toelaten.

Kant onderscheidt vervolgens twee elementen die bij dit ideaal betrokken zijn. Ten eerste is er de esthetische normale idee, een individuele intuïtie van de verbeelding die dient als maatstaf voor wat tot een bepaalde diersoort behoort. Deze normale idee ontstaat door een psychologisch proces waarin de verbeelding talloze waarnemingen van individuen als het ware over elkaar heen legt en zo een gemiddelde vorm produceert die als maatstaf kan dienen. Deze gemiddelde vorm is niet het ideaal van schoonheid, maar slechts de correcte voorstelling van de soort, de noodzakelijke voorwaarde waaronder schoonheid mogelijk is. Zij bevat niets karakteristieks en behaagt niet door schoonheid, maar door haar overeenstemming met de voorwaarden van de soort.

Ten tweede is er de idee van de rede, die de morele doeleinden van de mens tot principe maakt voor de beoordeling van zijn uiterlijke verschijning. Het ideaal van de menselijke schoonheid bestaat in de zichtbare uitdrukking van morele ideeën, zoals goedheid, zuiverheid, kracht of rust. Deze morele kwaliteiten kunnen slechts uit ervaring worden afgeleid, maar hun zichtbare samenhang met de uiterlijke vorm vereist een verbeeldingskracht die door ideeën van de rede wordt geleid. Het ideaal van schoonheid is daarom niet louter een esthetische maatstaf, maar een samengaan van verbeelding en rede.

Kant besluit dat het ideaal van schoonheid nooit een zuiver esthetisch oordeel kan zijn. Het laat geen zintuiglijke bekoring toe, maar wekt toch een groot belang. Dit toont aan dat een oordeel volgens een ideaal van schoonheid geen zuiver smaakoordeel is, maar een oordeel waarin esthetische en morele elementen samenkomen. Het ideaal van schoonheid is daarmee een unieke grensfiguur in Kants esthetica: het behoort tot de sfeer van de smaak, maar overstijgt deze door zijn verbinding met de rede.

\bigskip
Kant stelt in deze paragraaf ondubbelzinnig dat er geen algemeen ideaal van schoonheid kan bestaan. Een ideaal veronderstelt een norm die niet louter uit de zintuiglijke verschijning voortkomt, maar uit een idee van de rede. Omdat schoonheid in haar zuivere vorm geen concept veronderstelt, kan er geen universeel model zijn dat voor alle objecten zou bepalen wat het volmaakte schone is. Alleen bij de mens is een ideaal van schoonheid denkbaar, omdat alleen de mens zijn doeleinden door de rede bepaalt en zijn uiterlijke verschijning kan worden beoordeeld in relatie tot morele ideeën. Kant werkt dit morele aspect in deze paragraaf nog niet volledig uit, maar hij maakt duidelijk dat het ideaal van menselijke schoonheid niet louter een gemiddelde of correcte vorm is, maar een zichtbare uitdrukking van innerlijke morele kwaliteiten.

Deze gedachte past naadloos in het bredere betoog van de *Kritiek van het oordeelsvermogen*. Kant bereidt hier de overgang voor naar de symbolische functie van schoonheid en naar de gedachte dat schoonheid een zintuiglijke analogie vormt voor moraliteit, een idee dat in §59 expliciet wordt uitgewerkt. De mens wordt zo het punt waar esthetiek en ethiek elkaar raken: in de menselijke figuur kan de verbeelding een zichtbare gestalte geven aan morele ideeën, waardoor schoonheid een brug wordt tussen natuur en vrijheid. Tegelijkertijd vormt deze analyse een opstap naar Kants bespreking van menselijke waardigheid, omdat het ideaal van schoonheid alleen mogelijk is bij een wezen dat zichzelf als doel kan bepalen.

De centrale vraag van deze paragraaf is waarom alleen de mens een ideaal van schoonheid kan hebben. Kant antwoordt impliciet dat een ideaal slechts mogelijk is wanneer de vorm van een object kan worden beoordeeld in relatie tot doeleinden die door de rede worden vastgesteld. Dieren, planten of artefacten hebben wel vormen en functies, maar zij bepalen hun doeleinden niet zelf. Hun doelmatigheid is gegeven door de natuur of door menselijke ontwerpkeuzes, en daarom kan hun schoonheid nooit meer zijn dan vrije of aanhangende schoonheid. De mens daarentegen is een wezen dat zijn doeleinden zelf formuleert en dat zijn uiterlijke verschijning kan laten corresponderen met innerlijke morele principes. Alleen bij hem kan de verbeelding een model vormen dat niet louter een gemiddelde van waarnemingen is, maar een zichtbare uitdrukking van een moreel ideaal.

Dit alles wijst op de bijzondere plaats van de mens tussen natuur en vrijheid. Als natuurwezen behoort hij tot de wereld van verschijningen, onderworpen aan causaliteit en biologische vorm. Als redelijk wezen behoort hij tot de wereld van vrijheid, waarin hij zichzelf doelen stelt en morele wetten volgt. Het ideaal van menselijke schoonheid ontstaat precies op het snijvlak van deze twee sferen: het lichaam wordt drager van een innerlijke morele betekenis, en de zintuiglijke vorm wordt een symbolische uitdrukking van vrijheid. In die zin toont het ideaal van schoonheid hoe de mens tegelijk een natuurlijk organisme en een moreel subject is, en hoe zijn verschijning een unieke plaats inneemt in de orde van de natuur doordat zij het zichtbare spoor draagt van een onzichtbare, vrije rede.
\end{mdframed}

\paragraph{Definitie van het mooie afgeleid uit dit derde moment}
\textbf{Schoonheid} is de vorm van de \textbf{doelmatigheid} van een object, voor zover deze daarin wordt waargenomen \textbf{zonder voorstelling van een doel}.

\moment{Vierde moment van het smaakoordeel, betreffende de modalieit van de voldoening in het object.}

\section{De modaliteit van het smaakoordeel.} % §18

Van iedere voorstelling kan ik zeggen dat het ten minste \textbf{mogelijk} is dat zij (als een kennis) met een genoegen wordt verbonden. Van datgene wat ik \textbf{aangenaam} noem, zeg ik dat het \textbf{daadwerkelijk} een genoegen in mij teweegbrengt. Van het schone daarentegen denkt men dat het een noodzakelijke betrekking tot voldoening heeft. Nu is deze noodzakelijkheid van een bijzondere soort: niet een theoretische objectieve noodzakelijkheid, waarbij \emph{a priori} kan worden gekend dat iedereen deze voldoening \textbf{zal voelen} bij het object dat door mij schoon wordt genoemd, noch een praktische noodzakelijkheid, waarbij door middel van begrippen van een zuivere wil, die als regels voor vrij handelende wezens dienen, deze voldoening een noodzakelijke gevolg is van een objectieve wet en niets anders betekent dan dat men absoluut (zonder een verder doel) op een bepaalde wijze behoort te handelen. Veeleer kan zij, als een noodzakelijkheid die in een esthetisch oordeel wordt gedacht, slechts \textbf{exemplarisch} worden genoemd, dat wil zeggen een noodzakelijkheid van de instemming van \textbf{allen} met een oordeel dat wordt beschouwd als een voorbeeld van een universele regel die men niet kan voortbrengen. Aangezien een esthetisch oordeel geen objectief en kennend oordeel is, kan deze noodzakelijkheid niet uit bepaalde begrippen worden afgeleid en is zij derhalve niet apodictisch. Nog veel minder kan zij worden afgeleid uit de universaliteit van de ervaring (uit een volledige eenstemmigheid in oordelen over de schoonheid van een bepaald object). Want niet alleen zou de ervaring daarvoor nauwelijks voldoende bewijsmateriaal leveren, maar bovendien is het onmogelijk om enig begrip van de noodzakelijkheid van deze oordelen op empirische oordelen te gronden.

\begin{mdframed}[leftmargin=0pt,rightmargin=0pt,backgroundcolor=blue!6,linecolor=white,innertopmargin=6pt,innerbottommargin=6pt]
Kant onderzoekt in deze paragraaf de bijzondere vorm van noodzakelijkheid die aan het smaakoordeel wordt toegeschreven. Hij begint met de vaststelling dat elke voorstelling in beginsel met een gevoel van genoegen kan worden verbonden, en dat het aangename datgene is waarvan men zegt dat het daadwerkelijk een genoegen teweegbrengt. Bij het schone echter denkt men dat het noodzakelijk met een gevoel van voldoening verbonden is. Deze noodzakelijkheid is echter van een geheel eigen soort. Zij is niet theoretisch, want men kan niet a priori aantonen dat iedereen hetzelfde genoegen zal ervaren bij een object dat ik mooi noem. Zij is evenmin praktisch, want zij berust niet op een wet van de wil die voorschrijft dat men op een bepaalde wijze moet handelen en waarin het genoegen een noodzakelijke consequentie van een morele wet zou zijn.

De noodzakelijkheid die in een esthetisch oordeel wordt gedacht, noemt Kant exemplarisch. Dit betekent dat men een instemming van iedereen verwacht, niet omdat men een algemene regel kan geven die het oordeel bepaalt, maar omdat men het eigen oordeel beschouwt als een voorbeeld van een regel die men niet kan formuleren. Het smaakoordeel pretendeert dus een zekere universaliteit, maar zonder dat deze op begrippen kan worden gebaseerd. Omdat een esthetisch oordeel geen objectief, cognitief oordeel is, kan de noodzakelijkheid ervan niet uit bepaalde begrippen worden afgeleid en is zij daarom niet apodictisch.

Kant benadrukt bovendien dat deze noodzakelijkheid niet uit ervaring kan worden afgeleid. Zelfs al zou men een grote mate van overeenstemming in oordelen over schoonheid aantreffen, dan nog zou dit geen grond kunnen vormen voor een noodzakelijk oordeel. Ervaring kan nooit een concept van noodzakelijkheid voortbrengen, en daarom kan de universaliteit die het smaakoordeel vooronderstelt niet empirisch worden gerechtvaardigd. De noodzakelijkheid van het esthetische oordeel blijft daarmee een bijzondere, niet-conceptuele en niet-empirische vorm van aanspraak op algemene geldigheid.

\bigskip
Kant maakt in deze paragraaf duidelijk dat het smaakoordeel een bijzondere vorm van noodzakelijkheid claimt. Schoonheid wordt voorgesteld alsof zij noodzakelijk met een gevoel van voldoening verbonden is, maar deze noodzakelijkheid is niet logisch, want zij kan niet uit begrippen worden afgeleid, en evenmin empirisch, want zij kan niet worden gebaseerd op feitelijke overeenstemming in ervaring. Toch spreekt het smaakoordeel alsof iedereen zou moeten instemmen met het oordeel dat iets mooi is. Deze spanning vormt de kern van Kants analyse: het smaakoordeel doet een normatieve aanspraak zonder dat er een wet, bewijs of dwang aan ten grondslag ligt.

Binnen het bredere betoog van de *Kritiek van het oordeelsvermogen* is dit een cruciale stap. Kant probeert te laten zien hoe normativiteit mogelijk is zonder dat zij op begrippen of morele wetten berust. Het smaakoordeel is een voorbeeld van een “moeten” dat niet uit een regel volgt, maar toch een aanspraak op universaliteit maakt. Daarmee vormt het een overgangsgebied tussen natuur en vrijheid: het is geen kennis en geen moraal, maar een vorm van oordelen waarin de menselijke vermogens zich op een vrije, maar toch gedeelde wijze ordenen. Deze structuur bereidt de weg voor Kants latere stelling dat schoonheid een symbolische relatie heeft tot moraliteit, omdat beide domeinen een vorm van noodzakelijkheid kennen die niet uit empirische feiten of logische regels voortkomt.

De centrale vraag van deze paragraaf is wat voor soort noodzakelijkheid het smaakoordeel precies claimt. Kant noemt deze noodzakelijkheid exemplarisch. Zij is niet dwingend, want niemand kan worden gedwongen iets mooi te vinden. Zij is niet bewijsbaar, want er bestaat geen concept of regel die het oordeel kan rechtvaardigen. Toch spreekt het smaakoordeel alsof anderen zouden moeten instemmen, alsof het oordeel een voorbeeld is van een regel die niet kan worden geformuleerd maar wel wordt verondersteld. Het is een “moet” zonder wet: een aanspraak op instemming die niet uit verplichting voortkomt, maar uit het vertrouwen dat anderen dezelfde vrije harmonie van hun vermogens zullen ervaren.

Een dergelijke noodzakelijkheid ervaren we vandaag nog steeds in domeinen waar we geen regels kunnen geven, maar toch een gedeelde reactie verwachten. Dit gebeurt bijvoorbeeld in kunstkritiek, wanneer we ervan uitgaan dat een bepaald kunstwerk “voor iedereen” indrukwekkend of geslaagd is, zonder dat we dit kunnen bewijzen. Het gebeurt ook in morele intuïties die nog niet tot expliciete regels zijn uitgewerkt, maar waarvan we toch verwachten dat anderen ze zullen delen. En het gebeurt in alledaagse esthetische ervaringen, zoals het bewonderen van een landschap, een muziekstuk of een architectonische ruimte, waarin we spontaan het gevoel hebben dat onze waardering niet louter persoonlijk is, maar een zekere universaliteit draagt. In al deze gevallen ervaren we een vorm van noodzakelijkheid die niet uit dwang of bewijs voortkomt, maar uit een gedeelde menselijke gevoeligheid die zich niet in regels laat vangen.

\end{mdframed}

\section{De subjectieve noodzakelijkheid die wij aan het smaakoordeel toeschrijven, is voorwaardelijk.} % §19

Het smaakoordeel schrijft instemming aan iedereen toe, en wie iets als mooi verklaart, wenst dat iedereen het object in kwestie \textbf{zou} goedkeuren en het evenzo mooi zou verklaren. Het \textbf{zou moeten} in esthetische smaakoordelen wordt dus slechts voorwaardelijk uitgesproken, zelfs wanneer alle gegevens die voor het oordelen vereist zijn, gegeven zijn. Men verzoekt om instemming van ieder ander, omdat men daarvoor een grond heeft die aan allen gemeen is; men zou zelfs op deze instemming kunnen rekenen, als men er maar altijd zeker van was dat het geval correct onder die grond als regel van goedkeuring werd gesubsumeerd.

\begin{mdframed}[leftmargin=0pt,rightmargin=0pt,backgroundcolor=blue!6,linecolor=white,innertopmargin=6pt,innerbottommargin=6pt]
Kant benadrukt in deze passage dat het smaakoordeel een aanspraak op algemene instemming bevat. Wie iets mooi noemt, verlangt dat iedereen dit object eveneens als mooi zal erkennen. Deze aanspraak wordt echter slechts voorwaardelijk uitgesproken, zelfs wanneer alle gegevens die voor het oordeel nodig zijn aanwezig zijn. Het smaakoordeel formuleert een “zou moeten” dat niet absoluut is, maar afhankelijk blijft van de juiste toepassing van een gemeenschappelijke grondslag. Men vraagt om instemming omdat men ervan uitgaat dat er een basis is die door alle mensen wordt gedeeld, een gemeenschappelijke structuur van het oordeelsvermogen die het mogelijk maakt dat anderen hetzelfde genoegen ervaren. Indien men er zeker van kon zijn dat het concrete geval correct onder deze gemeenschappelijke grond werd gebracht, zou men zelfs op die instemming kunnen rekenen. Toch blijft deze noodzakelijkheid voorwaardelijk, omdat het altijd mogelijk is dat de subsumptie onder de gedeelde grondslag niet volledig of niet juist is uitgevoerd.

\bigskip
Kant maakt in deze passage duidelijk dat het smaakoordeel alleen maar aanspraak kan maken op instemming van anderen omdat het uitgaat van een gedeelde structuur van onze vermogens. Wanneer iemand iets mooi noemt, doet hij dat niet als een privé-ervaring, maar vanuit de veronderstelling dat anderen over dezelfde wijze van waarnemen en oordelen beschikken. Deze gedeelde grond is geen empirisch feit en geen sociologische observatie; het is een noodzakelijke voorwaarde voor de mogelijkheid van een esthetisch oordeel zelf. Zonder deze impliciete aanname zou het woord “mooi” niets anders betekenen dan “ik vind het prettig”, en zou het smaakoordeel zijn normatieve karakter verliezen.

In het bredere betoog van de \emph{Kritiek van het oordeelsvermogen} vormt dit een belangrijke stap richting Kants visie op cultuur, opvoeding en ritueel. Esthetische oordelen zijn niet louter individuele reacties, maar momenten waarin onze vermogens zich op een manier ordenen die in principe door iedereen kan worden gedeeld. Daardoor krijgen zij een rol in de vorming van een gemeenschappelijke wereld: zij oefenen ons in het afstemmen van onze subjectieve ervaringen op een horizon van mogelijke instemming. In die zin is het esthetische oordeel een oefening in het samenleven, nog vóór er sprake is van expliciete regels of morele verplichtingen.

De centrale vraag van deze paragraaf is wat het esthetische oordeel communiceerbaar maakt. Kant antwoordt dat dit niet voortkomt uit feitelijke overeenstemming, noch uit culturele conventies, maar uit de structuur van het oordeel zelf. Het esthetische oordeel beroept zich op een gevoel waarvan men veronderstelt dat het door iedereen kan worden gedeeld, omdat het voortkomt uit een vrije harmonie van verbeelding en verstand die niet afhankelijk is van individuele voorkeuren. Deze gedeeldheid is niet empirisch aantoonbaar, maar wordt gedacht als een noodzakelijke mogelijkheid. Het smaakoordeel is communiceerbaar omdat het zich richt tot een “ieder” dat niet uit ervaring wordt afgeleid, maar uit de aard van het oordelen zelf.

Vanuit dit perspectief kan men zeggen dat esthetiek inderdaad een oefening in gemeenschap is. Niet omdat zij mensen feitelijk verenigt, maar omdat zij een ruimte opent waarin wij onze subjectieve ervaringen zo vormgeven dat zij in principe door anderen kunnen worden gedeeld. Esthetiek traint ons in het aanspreken van een gemeenschappelijke menselijkheid, zonder dat wij daarvoor een wet, een concept of een doel nodig hebben. Zij vormt een subtiele, maar fundamentele manier waarop wij leren onszelf te situeren in een wereld die wij met anderen delen.
\end{mdframed}

\section{De voorwaarde van de noodzakelijkheid die door een smaakoordeel wordt opgeëist, is het idee van een gemeenschappelijk zintuig.} % §20

Als oordelen van smaak (zoals cognitieve oordelen) een bepaald objectief principe zouden hebben, dan zou degene die ze overeenkomstig dat principe velt, aanspraak maken op de onvoorwaardelijke noodzakelijkheid van zijn oordeel. Als zij daarentegen helemaal geen principe hadden, zoals die van louter zintuiglijke smaak, dan zou men zelfs nooit op het idee van hun noodzakelijkheid komen. Zij moeten dus een subjectief principe hebben, dat bepaalt wat behaagt of mishaaagt uitsluitend door gevoel en niet door begrippen, maar toch met universele geldigheid. Een dergelijk principe kan echter alleen worden opgevat als een \textbf{gemeenschappelijke zin}, die wezenlijk verschilt van het gewone verstand, dat soms eveneens gezond verstand (\emph{sensus communis}) wordt genoemd, aangezien dit laatste niet door gevoel oordeelt maar altijd door begrippen, zij het doorgaans slechts in de vorm van vaag voorgestelde principes.

Alleen onder de vooronderstelling dat er een gemeenschappelijke zin is, waarmee wij echter geen uiterlijke zin bedoelen maar het effect van het vrije spel van onze kenvermogens, alleen onder de vooronderstelling van zo’n gemeenschappelijke zin, zeg ik, kan het oordeel van smaak worden geveld.

\begin{mdframed}[leftmargin=0pt,rightmargin=0pt,backgroundcolor=blue!6,linecolor=white,innertopmargin=6pt,innerbottommargin=6pt]
Kant stelt in deze passage dat, wanneer oordelen van smaak een vast en objectief principe zouden hebben zoals cognitieve oordelen, degene die volgens dat principe oordeelt aanspraak zou kunnen maken op een onvoorwaardelijke noodzakelijkheid van zijn oordeel. Als zulke oordelen daarentegen helemaal geen principe zouden hebben, zoals louter zintuiglijke voorkeuren, dan zou men nooit op het idee komen dat zij enige noodzakelijkheid bezitten. Daarom moeten oordelen van smaak een subjectief principe hebben dat bepaalt wat behaagt of mishaagt, niet door middel van begrippen maar uitsluitend via gevoel, en dat toch aanspraak maakt op algemene geldigheid. Een dergelijk principe kan volgens Kant alleen worden opgevat als een gemeenschappelijk zintuig, een vermogen dat wezenlijk verschilt van het gewone verstand dat soms eveneens ‘common sense’ wordt genoemd. Dat verstand oordeelt immers altijd door middel van begrippen, ook al zijn die vaak slechts vaag of impliciet aanwezig, terwijl het gemeenschappelijk zintuig waarop het smaakoordeel berust juist niet begripsmatig maar gevoelsmatig functioneert. Kant benadrukt dat het bestaan van dit gemeenschappelijk zintuig niet moet worden opgevat als een extern zintuig, maar als het effect van het vrije spel van onze kenvermogens. Alleen wanneer men dit gemeenschappelijk zintuig veronderstelt, is het mogelijk een smaakoordeel te vellen. Zonder deze vooronderstelling zou het smaakoordeel zijn aanspraak op universaliteit en noodzakelijkheid verliezen en zou het niet langer als een smaakoordeel kunnen worden begrepen.

\bigskip

\end{mdframed}

\section{Of men goede reden heeft om een gemeenschappelijk zintuig te veronderstellen.} % §21

Kenningen en oordelen moeten, samen met de overtuiging die hen vergezelt, universeel meedeelbaar kunnen zijn, want anders zouden zij geen overeenstemming met het object hebben: zij zouden dan allemaal slechts een louter subjectief spel van de voorstellingsvermogens zijn, precies zoals het scepticisme beweert. Maar als kenningen meedeelbaar moeten kunnen zijn, dan moet ook de gemoedstoestand, dat wil zeggen de gesteldheid van de kenvermogens voor een kenning in het algemeen, en wel die verhouding die geschikt is om uit een voorstelling (waardoor een object ons gegeven wordt) kenning te maken, universeel meedeelbaar kunnen zijn; want zonder deze, als subjectieve voorwaarde van het kennen, zou de kenning als effect niet kunnen ontstaan. En dit gebeurt inderdaad telkens wanneer een gegeven object, door middel van de zintuigen, de verbeeldingskracht in werking brengt voor de synthese van het veelvoudige, terwijl de verbeeldingskracht het verstand in werking brengt voor de eenmaking van het veelvoudige in begrippen. Maar deze gesteldheid van de kenvermogens heeft een verschillende verhouding naargelang het verschil van de gegeven objecten. Niettemin moet er één zijn waarin deze innerlijke verhouding optimaal is voor de bezieling van beide vermogens van de geest, het ene door het andere, met betrekking tot kenning van gegeven objecten in het algemeen; en deze gesteldheid kan niet anders dan door gevoel worden bepaald, niet door begrippen. Nu moet deze gesteldheid zelf universeel meedeelbaar kunnen zijn, en dus ook het gevoel ervan, in het geval van een gegeven voorstelling; maar aangezien de universele meedeelbaarheid van een gevoel een gemeenschappelijke zin veronderstelt, moet deze met goede reden kunnen worden aangenomen, en wel zonder beroep op psychologische waarnemingen, maar veeleer als de noodzakelijke voorwaarde van de universele meedeelbaarheid van onze kenning, die in elke logica en in elk principe van kenningen dat niet sceptisch is, wordt verondersteld.

\begin{mdframed}[leftmargin=0pt,rightmargin=0pt,backgroundcolor=blue!6,linecolor=white,innertopmargin=6pt,innerbottommargin=6pt]
\section*{Samenvatting}

Kant betoogt in deze passage dat zowel cognities als oordelen, samen met de overtuiging die hen begeleidt, in principe universeel mededeelbaar moeten zijn. Zonder deze mogelijkheid tot algemene communicatie zouden zij geen enkele relatie tot een object kunnen hebben en zouden zij slechts een subjectief spel van onze voorstellingsvermogens zijn, precies zoals het scepticisme beweert. Voor communicatie van cognities is het echter noodzakelijk dat ook de innerlijke toestand van de geest, dat wil zeggen de specifieke dispositie van de kenvermogens die een cognitie mogelijk maakt, universeel mededeelbaar is. Deze dispositie betreft de verhouding waarin verbeelding en verstand samenwerken om uit een gegeven voorstelling een object te kunnen kennen. Telkens wanneer een object via de zintuigen wordt aangeboden, activeert het de verbeelding om de veelheid van indrukken te synthetiseren, terwijl de verbeelding op haar beurt het verstand activeert om deze veelheid onder een concept te verenigen.

Hoewel deze verhouding tussen verbeelding en verstand varieert afhankelijk van het soort object dat wordt waargenomen, moet er volgens Kant een specifieke verhouding bestaan die optimaal is voor het tot stand brengen van cognitie in het algemeen. Deze optimale verhouding kan niet door begrippen worden vastgesteld, maar slechts door een gevoel worden ervaren. Omdat deze innerlijke dispositie zelf universeel mededeelbaar moet zijn, moet ook het gevoel dat ermee gepaard gaat universeel mededeelbaar zijn. De mogelijkheid van een dergelijk gedeeld gevoel veronderstelt echter het bestaan van een gemeenschappelijk zintuig. Dit gemeenschappelijk zintuig hoeft niet psychologisch te worden aangetoond, maar moet worden aangenomen als een noodzakelijke voorwaarde voor de universele mededeelbaarheid van onze cognities, een voorwaarde die in elke niet-sceptische logica en in elk beginsel van kennis impliciet wordt verondersteld.

\bigskip
Kant maakt in deze passage expliciet dat het smaakoordeel alleen mogelijk is onder de vooronderstelling van een gedeelde menselijkheid. Wanneer wij oordelen dat iets mooi is, doen wij dat niet als geïsoleerde individuen, maar vanuit het idee dat onze wijze van waarnemen en oordelen in principe door iedereen kan worden gedeeld. Deze gedeeldheid is geen empirisch feit en geen psychologische observatie, maar een noodzakelijke voorwaarde voor de mogelijkheid van kennis én voor de mogelijkheid van esthetische communicatie. Het esthetische oordeel veronderstelt dus een menselijk vermogen dat voorafgaat aan alle verschillen in cultuur, opvoeding of ervaring: een gemeenschappelijk zintuig dat niet uit begrippen bestaat, maar uit een gedeelde structuur van onze kenvermogens.

Binnen het bredere betoog van de \emph{Kritiek van het oordeelsvermogen} vormt dit de antropologische basis van Kants esthetiek. Schoonheid is niet zomaar een bijkomstigheid van menselijke ervaring, maar een aanwijzing voor hoe onze vermogens in hun meest vrije en harmonieuze vorm samenwerken. Esthetiek laat zien dat de mens niet alleen een kennend en handelend wezen is, maar ook een wezen dat zijn subjectieve ervaringen zo kan vormen dat zij in principe door anderen gedeeld kunnen worden. In die zin is het esthetische oordeel een oefening in menselijkheid: het toont hoe wij onszelf begrijpen als wezens die niet opgesloten zitten in hun eigen subjectiviteit, maar gericht zijn op een gedeelde wereld.

De centrale vraag van deze paragraaf is wat esthetiek zegt over wat het betekent mens te zijn. Kants antwoord is dat de mens een wezen is dat zijn subjectieve ervaring kan universaleren zonder haar te objectiveren. Esthetiek onthult een vermogen tot gemeenschap dat voorafgaat aan taal, moraal of kennis. Het laat zien dat wij in staat zijn onze innerlijke toestand zo te ervaren dat zij een aanspraak op universaliteit draagt, niet omdat zij door begrippen wordt geleid, maar omdat zij voortkomt uit een vrije harmonie van onze vermogens. Esthetiek is daarmee geen empirische psychologie, maar een reflectie op de voorwaarden waaronder menselijke ervaring überhaupt gedeeld kan worden.

Vanuit dit perspectief kan men zeggen dat schoonheid inderdaad een stille vorm van ethiek is. Niet omdat schoonheid morele regels bevat, maar omdat zij een structuur van gedeeldheid en universaliteit laat zien die verwant is aan de moraal. In de ervaring van het schone oefenen wij ons in een houding die ook voor de ethiek essentieel is: het vermogen om onze subjectieve standpunten te verbinden met een idee van wat voor iedereen zou kunnen gelden. Schoonheid spreekt geen geboden uit, maar zij vormt een stille voorafbeelding van morele gemeenschap. In het esthetische oordeel ervaren wij hoe vrijheid en universaliteit elkaar kunnen raken zonder dwang, en precies daarin toont schoonheid haar verwantschap met de ethiek.
\end{mdframed}

\section{De noodzakelijkheid van de algemene instemming die in een smaakoordeel wordt gedacht, is een subjectieve noodzakelijkheid, die als objectief wordt voorgesteld onder de vooronderstelling van een gemeenschappelijk zintuig.}

In alle oordelen waarin wij iets mooi noemen, staan wij niemand toe van een andere mening te zijn, zonder ons oordeel echter op begrippen te gronden, maar uitsluitend op ons gevoel, dat wij daarom niet als een privégevoel, maar als een gemeenschappelijk gevoel tot grondslag maken. Deze gemeenschappelijke zin kan voor dit doel echter niet op ervaring worden gegrond, want hij moet oordelen rechtvaardigen die een “zou moeten” bevatten: hij zegt niet dat iedereen met ons oordeel \textbf{zal} instemmen, maar dat iedereen ermee \textbf{zou moeten} instemmen. Dus is de gemeenschappelijke zin, waarvan ik hier mijn smaakoordeel als voorbeeld aanvoer en waaraan ik daarom \textbf{voorbeeldige} geldigheid toeschrijf, een louter ideale norm, onder de veronderstelling waarvan men terecht een oordeel dat ermee overeenstemt, en de daarin uitgedrukte voldoening aan een object, tot regel voor iedereen zou kunnen maken: aangezien het principe, hoewel slechts subjectief, toch als subjectief universeel, een voor iedereen noodzakelijke idee, wordt aangenomen, die, wat de overeenstemming van verschillende oordelaars betreft, evenzeer universele instemming zou kunnen eisen als een objectief principe, mits men er maar zeker van was dat men het geval er correct onder had gesubsumeerd.

Deze onbepaalde norm van een gemeenschappelijke zin wordt door ons inderdaad verondersteld: onze aanspraak bij het vellen van smaakoordelen bewijst dit. Of er werkelijk zo’n gemeenschappelijke zin bestaat, als constitutief principe van de mogelijkheid van ervaring, of dat een nog hoger principe van de rede hem slechts eerst tot een regulatief principe voor ons maakt, om een gemeenschappelijke zin in onszelf voort te brengen voor hogere doeleinden, en dus of smaak een oorspronkelijke en natuurlijke vermogensfaculteit is, of slechts de idee van een nog te verwerven, kunstmatige faculteit, zodat een smaakoordeel met zijn verwachting van universele instemming in feite slechts een eis van de rede is om zo’n overeenstemming in de wijze van voelen tot stand te brengen, en of het “zou moeten”, dat wil zeggen de objectieve noodzakelijkheid van het samenvallen van ieders gevoel met dat van ieder afzonderlijk, slechts de mogelijkheid betekent om tot overeenstemming te komen, en het smaakoordeel slechts een voorbeeld levert van de toepassing van dit principe – dit willen en kunnen wij hier nog niet onderzoeken; vooralsnog hebben wij slechts de taak het smaakvermogen in zijn elementen te ontleden en deze uiteindelijk te verenigen in de idee van een gemeenschappelijke zin.

\begin{mdframed}[leftmargin=0pt,rightmargin=0pt,backgroundcolor=blue!6,linecolor=white,innertopmargin=6pt,innerbottommargin=6pt]

Kant stelt in deze passage dat wanneer wij iets mooi noemen, wij impliciet verwachten dat niemand het daar fundamenteel mee oneens zal zijn. Deze aanspraak op instemming berust echter niet op begrippen of objectieve gronden, maar op een gevoel dat wij niet als louter particulier beschouwen, maar als iets dat in principe door iedereen gedeeld kan worden. Dit veronderstelde gemeenschappelijke gevoel kan niet uit ervaring worden afgeleid, omdat het bedoeld is om oordelen te rechtvaardigen die een normatief karakter hebben. Het smaakoordeel zegt immers niet dat iedereen feitelijk zal instemmen, maar dat iedereen zou moeten instemmen. Daarom moet het gemeenschappelijk zintuig dat aan deze aanspraak ten grondslag ligt worden opgevat als een louter ideale norm, een vooronderstelling die het mogelijk maakt om de eigen voldoening in een object als maatgevend voor iedereen te beschouwen. Hoewel het principe waarop dit berust slechts subjectief is, wordt het toch voorgesteld als subjectief universeel, als een idee die noodzakelijk is voor ieder mens en die, voor zover het om de overeenstemming van verschillende beoordelaars gaat, een aanspraak op algemene instemming kan doen die vergelijkbaar is met die van een objectief principe, mits men zeker zou zijn van de juiste toepassing ervan.

Kant benadrukt dat wij deze onbepaalde norm van een gemeenschappelijk zintuig daadwerkelijk veronderstellen, zoals blijkt uit de manier waarop wij smaakoordelen uitspreken. Of er werkelijk zo’n gemeenschappelijk zintuig bestaat als constitutieve voorwaarde voor ervaring, of dat een hoger principe van de rede het slechts als regulatief voorschrift invoert om ons ertoe aan te zetten een dergelijke gedeeldheid in onszelf te ontwikkelen, laat Kant hier in het midden. Evenmin beslist hij of smaak een oorspronkelijke en natuurlijke aanleg is, of slechts een idee van een vermogen dat nog moet worden verworven en dus kunstmatig is. In dat laatste geval zou de aanspraak op algemene instemming in het smaakoordeel slechts een eis van de rede zijn om een mogelijke overeenstemming in onze wijze van voelen tot stand te brengen, en zou het smaakoordeel slechts een voorbeeld zijn van de toepassing van dit principe. Kant stelt dat deze vragen op dit punt nog niet onderzocht kunnen worden. Zijn taak is voorlopig beperkt tot het analyseren van het vermogen tot oordelen van smaak en het uiteindelijk samenbrengen van zijn elementen in de idee van een gemeenschappelijk zintuig.

\bigskip
Kant formuleert in deze passage impliciet een samenvattende definitie van schoonheid: het schone is datgene waarvan wij vinden dat iedereen ermee zou moeten instemmen, hoewel dit oordeel uitsluitend op een gevoel berust. Schoonheid is dus geen eigenschap van het object, maar een bijzondere verhouding tussen onze vermogens die wij toch als normatief voorstellen. In die zin vat Kant hier de hele analyse van het smaakoordeel samen: schoonheid is een subjectieve voldoening die een aanspraak op universaliteit maakt, gedragen door de idee van een gemeenschappelijk zintuig.

Binnen het bredere betoog van de *Kritiek van het oordeelsvermogen* vormt dit een scharnierpunt. Kant heeft nu de structuur van het esthetische oordeel blootgelegd en kan daardoor de stap zetten naar de symbolische functie van schoonheid, die in §59 expliciet wordt uitgewerkt. De esthetische ervaring blijkt een model te bieden voor hoe onze vermogens vrij kunnen samenwerken en toch een aanspraak op universaliteit kunnen dragen. Precies dat maakt schoonheid tot een brug tussen natuur en vrijheid, en daarmee tot een noodzakelijke voorbereiding op de moraal. De analyse van het schone is dus geen afgesloten hoofdstuk, maar een fundament waarop Kant verder bouwt.

De centrale vraag is wat we nu precies gewonnen hebben met deze analyse. Kant heeft laten zien dat het smaakoordeel geen willekeurige voorkeur is, maar een gestructureerde vorm van oordelen die een bijzondere plaats inneemt tussen kennis en moraal. We begrijpen nu dat schoonheid niet verklaard kan worden door begrippen of empirische feiten, maar dat zij een noodzakelijke vooronderstelling van gedeelde menselijkheid veronderstelt. Daarmee wordt duidelijk waarom esthetiek niet gaat over objecten, maar over de voorwaarden waaronder mensen een gemeenschappelijke wereld kunnen ervaren. Wat we gewonnen hebben, is dus een inzicht in de architectuur van onze vermogens: schoonheid toont hoe subjectiviteit zich kan openen naar universaliteit zonder dwang.

Vanuit dit perspectief wordt duidelijk waarom schoonheid voor Kant onmisbaar is in het filosofisch systeem. Zij vormt het enige domein waarin wij ervaren hoe vrijheid en natuur elkaar kunnen raken zonder dat de ene de andere opheft. In de esthetische ervaring zien we hoe onze vermogens spontaan in harmonie kunnen samenwerken, en hoe een subjectief gevoel toch een aanspraak op universaliteit kan dragen. Deze structuur is een voorafbeelding van de moraal, waarin vrijheid zich eveneens moet verbinden met universaliteit. Schoonheid is daarom geen ornament van de filosofie, maar een noodzakelijke schakel: zij laat zien hoe een vrije, gedeelde wereld überhaupt denkbaar is.
\end{mdframed}

\paragraph{Definitie van schoonheid afgeleid uit het vierde moment.} |textbf{Schoonheid} is datgene wat zonder begrip wordt gekend als het object van een \textbf{noodzakelijke} voldoening.

\section*{Algemene opmerking over het Eerste Deel van de Analytiek}

Als men de conclusie trekt uit bovenstaande analyses, blijkt dat alles voortvloeit uit het begrip van smaak als een vermogen om een object te oordelen in relatie tot de \textbf{vrije wettigheid} van de verbeelding. Maar als in het smaakoordeel de verbeelding in haar vrijheid moet worden beschouwd, dan wordt zij in eerste instantie niet opgevat als reproductief, als onderworpen aan de wetten van associatie, maar als productief en zelfactief (als de auteur van vrijwillige vormen van mogelijke intuïties); en hoewel zij bij het waarnemen van een gegeven zintuiglijk object uiteraard gebonden is aan een bepaalde vorm van dit object en daardoor geen vrije speelruimte heeft (zoals bij uitvinding), is het toch heel goed voorstelbaar dat het object haar een vorm kan verschaffen die juist een zodanige compositie van het veelvoud bevat zoals de verbeelding die in harmonie met de \textbf{wettigheid van het verstand} in het algemeen zou ontwerpen als zij geheel vrij gelaten werd. Toch is het een contradictie dat de \textbf{verbeelding vrij} en tegelijkertijd \textbf{wettig door zichzelf} is, dat wil zeggen dat zij autonomie met zich meebrengt. Alleen het verstand geeft de wet. Maar wanneer de verbeelding gedwongen wordt te handelen volgens een bepaalde wet, wordt bepaald hoe haar product, wat de vorm betreft, moet zijn door begrippen; maar dan, zoals hierboven werd gezegd, is de voldoening niet die in het schone, maar in het goede (van volmaaktheid, in elk geval slechts de formele soort), en is het oordeel geen oordeel door middel van smaak. Zo is alleen een wettigheid zonder wet en een subjectieve overeenkomst van de verbeelding met het verstand zonder een objectieve – waarbij de voorstelling betrekking heeft op een bepaald begrip van een object – verenigbaar met de vrije wettigheid van het verstand (die ook doelgerichtheid zonder doel wordt genoemd) en met de eigenaardigheid van een smaakoordeel.

Nu worden geometrisch regelmatige vormen – een cirkel, een vierkant, een kubus, enzovoort – vaak aangevoerd door critici van smaak als de eenvoudigste en meest onbetwijfelbare voorbeelden van schoonheid; en toch worden ze juist regelmatig genoemd omdat ze niet kunnen worden voorgesteld anders dan door ze op te vatten als louter voorstellingen van een bepaald begrip, dat de regel voor die vorm voorschrijft (waarmee ze alleen mogelijk is). Dus moet één van de twee fout zijn: of het oordeel van de critici dat zulke vormen schoonheid toeschrijft, of het onze, dat doelmatigheid zonder begrip noodzakelijk acht voor schoonheid.

Waarschijnlijk zal niemand denken dat het noodzakelijk is dat iemand smaak heeft om meer voldoening te vinden in de vorm van een cirkel dan in een krabbellijn, of meer in een gelijkzijdig en gelijkhoekig vierkant dan in een scheef en onregelmatig, als het ware misvormd; dit vereist slechts gezond verstand en helemaal geen smaak. Waar een doel voor ogen staat, bijvoorbeeld het beoordelen van de grootte van een oppervlakte of het begrijpen van de verhouding van de delen in een verdeling tot elkaar en tot het geheel, zijn regelmatige vormen, en inderdaad de eenvoudigste soort, noodzakelijk, en de voldoening rust dan niet direct op de aanblik van de vorm, maar op het nut ervan voor allerlei mogelijke doelen. Een kamer waarvan de wanden schuine hoeken vormen, een tuin van een soortgelijke aard, zelfs elke schending van symmetrie in de vorm van dieren (bijvoorbeeld één oog) evenals in gebouwen of bloemstukken, leidt tot ongenoegen, omdat het contraproposioneel is, niet alleen praktisch, met betrekking tot een bepaald gebruik van deze dingen, maar ook voor het oordeel ten aanzien van allerlei mogelijke doelen; dit is niet het geval in het smaakoordeel, dat, als het zuiver is, voldoening of ongenoegen onmiddellijk verbindt met de loutere \textbf{beschouwing} van het object, zonder acht te slaan op nut of doel.

De regelmaat die tot het begrip van een object leidt, is uiteraard de onmisbare voorwaarde (\emph{conditio sine qua non}) om het object in één voorstelling te vatten en het veelvoud in zijn vorm te bepalen. Deze bepaling is een doel met betrekking tot kennis; en in relatie daartoe is zij ook altijd verbonden met voldoening (die het gevolg is van het bereiken van elk doel, zelfs een louter problematisch doel). Maar dan is het slechts de goedkeuring van de oplossing die een probleem beantwoordt, en niet een vrije en onbepaald doelgerichte vermaak van de geestelijke vermogens met datgene wat wij het schone noemen, waarbij het verstand in dienst staat van de verbeelding en niet omgekeerd.

In een ding dat alleen mogelijk is door een intentie, in een gebouw, zelfs in een dier, moet de regelmaat die bestaat uit symmetrie de eenheid van de intuïtie uitdrukken, die het begrip van het doel begeleidt en tot de kennis behoort. Maar waar alleen een vrije speelruimte van de representatievermogens (alhoewel onder de voorwaarde dat het verstand daardoor geen schade lijdt) gehandhaafd moet worden, in lusthoven, in de inrichting van kamers, in allerlei smaakvolle gebruiksvoorwerpen en dergelijke, moet regelmaat die als dwang overkomt, zo veel mogelijk vermeden worden; vandaar dat de Engelse smaak in tuinen of de barokke smaak in meubelen de vrijheid van de verbeelding bijna tot het groteske drijft, en deze abstractie van alle dwang door regels het eigenlijke geval maakt waarin de smaak haar grootste volmaaktheid in projecten van de verbeelding kan tonen.

Alle stijve regelmaat (wat ook maar de nabijheid van wiskundige regelmaat benadert) is op zichzelf tegengesteld aan smaak: de beschouwing ervan verschaft geen duurzaam vermaak, maar integendeel, voor zover zij niet uitdrukkelijk kennis of een bepaald praktisch doel tot oogmerk heeft, wekt zij verveling op. Daarentegen is datgene waarmee de verbeelding op een ongedwongen en doelgerichte manier kan spelen, altijd nieuw voor ons, en raken wij er nooit op uitgekeken. In zijn beschrijving van Sumatra merkt \textbf{Marsden} op dat de vrije schoonheden van de natuur de waarnemer daar overal omringen en hem daarom weinig meer aantrekken; daarentegen had een pepertuin waarin de palen waarop de planten waren geleid parallelle rijen vormden, veel charme voor hem toen hij deze midden in een bos tegenkwam; en hieruit concludeert hij dat wilde, schijnbaar onregelmatige schoonheid alleen als afwisseling aangenaam is voor iemand die genoeg heeft van de regelmatige soort. Maar hij had slechts de proef hoeven nemen van één dag in zijn pepertuin om te beseffen dat zodra het verstand door middel van regelmaat op de orde is ingesteld die het altijd vereist, het object hem niet langer zou vermaken, maar eerder een belastende dwang aan de verbeelding zou opleggen, terwijl de natuur, die daar in haar variëteiten tot op het punt van overvloed uitbundig is, en aan geen kunstmatige regels onderworpen, zijn smaak duurzaam kon voeden. – Zelfs het gezang van de vogel, dat wij onder geen enkele muzikale regel kunnen brengen, lijkt meer vrijheid te bevatten en dus meer vermaak voor de smaak dan zelfs een menselijke zang die volgens alle regels van de muziek wordt uitgevoerd: men raakt namelijk veel sneller uitgekeken op het laatste als het vaak en langdurig wordt herhaald. Maar hier kunnen wij onze sympathie voor het plezier van een geliefd klein schepsel gemakkelijk verwarren met de schoonheid van zijn zang, die, wanneer zij exact door een mens wordt nagedaan (zoals soms met de noten van de nachtegaal gebeurt), ons oor volstrekt smaakloos treft.

Verder moeten mooie objecten onderscheiden worden van mooie aanzichten van objecten (die vanwege de afstand vaak niet meer duidelijk kunnen worden gekend). In het laatste lijkt de smaak zich niet zozeer te richten op wat de verbeelding in dit veld \textbf{waarneemt}, als wel op wat haar aanleiding geeft tot \textbf{verbeelding}, dat wil zeggen op wat strikt genomen de fantasieën zijn waarmee de geest zich vermaakt terwijl hij continu wordt aangespoord door het veelvoud dat het oog treft, zoals bijvoorbeeld bij het kijken naar de veranderende vormen van een vuur in een haard of van een kabbelend beekje, die beiden geen schoonheid bezitten, maar beiden charme voor de verbeelding bevatten, omdat zij haar vrije spel in stand houden.

\begin{mdframed}[leftmargin=0pt,rightmargin=0pt,backgroundcolor=blue!6,linecolor=white,innertopmargin=6pt,innerbottommargin=6pt]
Kant stelt in deze algemene opmerking dat oordelen van smaak een aanspraak maken op instemming van iedereen, ook al zijn ze niet gebaseerd op begrippen maar op gevoel. Wanneer we iets mooi noemen, doen we dat niet als een privémening, maar alsof iedereen het zou moeten kunnen beamen. Die aanspraak op algemeen geldige instemming kan niet worden verklaard als een puur objectieve noodzaak (zoals bij kennis) en ook niet als louter subjectief welbehagen (zoals bij zintuiglijke voorkeuren). Daarom veronderstelt het smaak-oordeel een subjectief principe met algemene geldigheid.

Dat principe noemt Kant een sensus communis, een gemeenschappelijke zin. Daarmee bedoelt hij niet het alledaagse verstand of gedeelde meningen, maar een gemeenschappelijke gesteldheid van onze kenvermogens: een bepaalde harmonie tussen verbeelding en verstand, die wordt ervaren als een gevoel van welbehagen. Deze innerlijke afstemming is niet via begrippen vast te stellen, maar alleen via gevoel, en toch moet zij in principe voor iedereen mededeelbaar zijn. Alleen zo kan een oordeel van smaak aanspraak maken op algemene instemming.

Kant benadrukt dat deze gemeenschappelijke zin geen empirisch gegeven is dat we uit ervaring afleiden. Zij fungeert als een ideaal of norm: een vooronderstelling die het mogelijk maakt om van anderen instemming te verlangen. Of deze sensus communis een aangeboren vermogen is, of eerder een ideaal dat door cultuur en oefening tot stand moet komen, laat hij hier in het midden. Voor dit moment is het voldoende dat het oordeel van smaak alleen begrijpelijk wordt wanneer we het idee van een gemeenschappelijke zin veronderstellen.

\bigskip
In deze algemene opmerking brengt Kant expliciet samen wat in de voorafgaande paragrafen is ontwikkeld. Hij herneemt dat schoonheid niet wordt beoordeeld via begrippen of kennisclaims, maar via een vrij en belangeloos gevoel van welbehagen. Dat gevoel is onmiskenbaar subjectief, en toch gaat het oordeel van smaak steeds gepaard met een aanspraak op algemene instemming. Die paradox wordt begrijpelijk doordat het esthetisch genoegen voortkomt uit het vrije spel van verbeelding en verstand, een innerlijke harmonie die niet dwingt maar uitnodigt. Juist omdat deze harmonie niet door belangen, doelen of concepten wordt bepaald, kan zij gelden als voorbeeldig voor iedereen. In deze samenvatting wordt duidelijk dat schoonheid bij Kant geen geïsoleerd fenomeen is, maar een specifieke manier waarop onze kenvermogens zich tot elkaar verhouden.

Tegelijk heeft deze opmerking een duidelijk voorbereidende functie. Kant positioneert het esthetisch oordeel als overgangsgebied tussen het louter subjectieve en het domein van natuur en moraal. Door te benadrukken dat het oordeel van smaak een reflectief oordeel is, bereidt hij de lezer voor op de gedachte dat natuurlijke schoonheid meer kan betekenen dan louter aangenaamheid. In de volgende paragrafen zal zij verschijnen als een symbool van moraliteit, juist omdat zij zonder concepten toch een orde en doelmatigheid laat ervaren. De algemene opmerking maakt zo duidelijk dat §§1–22 niet als afgesloten analyse moeten worden gelezen, maar als fundament voor een bredere systematische beweging waarin esthetiek, teleologie en ethiek met elkaar verbonden raken.

De kernvraag die hier naar voren komt, is wat we door de analyse van schoonheid hebben geleerd over onszelf als oordelende wezens. We hebben gezien dat een vrij, subjectief oordeel niet hoeft te verzanden in relativisme, zolang het berust op een gemeenschappelijke structuur van onze vermogens. Het idee van een sensus communis maakt het mogelijk dat individuele gevoelens toch normatieve pretentie hebben. Vanuit dit inzicht wordt begrijpelijk hoe natuurlijke schoonheid symbolisch kan zijn: zij toont, zonder morele voorschriften te geven, een wereld die ervaren kan worden als passend bij onze rede en onze morele bestemming. Het vrije spel van verbeelding en verstand fungeert daarbij als scharnier tussen individuele ervaring en universele idealen.

Dat Kant juist hier een algemene opmerking invoegt, is structureel betekenisvol. Hij onderbreekt de voortgang niet om iets nieuws te introduceren, maar om de lezer te heroriënteren en het bereikte overzichtelijk te maken voordat de tekst een nieuwe richting inslaat. Dit helpt om de spanning tussen subjectiviteit en universaliteit niet als een probleem te blijven zien, maar als een constitutief kenmerk van esthetische ervaring. Vanuit hedendaags perspectief laat dit zich herkennen in ervaringen van schoonheid in kunst, natuur of ritueel, waar men zich persoonlijk geraakt voelt en tegelijk het gevoel heeft dat die ervaring gedeeld zou moeten kunnen worden. Ook in symbolisch geladen praktijken, zoals rituelen waarin vorm, orde en betekenis samenkomen, kan het Kantiaanse idee van vrij esthetisch spel worden herkend als een ruimte waarin rede zich oefent zonder dwang, en waarin subjectieve ervaring een brug slaat naar gemeenschappelijke zin en morele oriëntatie.

\end{mdframed}

